% Artículo para la Revista Internacional de Investigación Multidisciplinaria (RIIM).
% Compilar con XeLaTeX (dos pasadas):
%   versión ciega para OJS:  xelatex articulo_RIIM.tex
%   versión con autores:     xelatex -jobname=articulo_RIIM_autores "\def\conautores{}% Artículo para la Revista Internacional de Investigación Multidisciplinaria (RIIM).
% Compilar con XeLaTeX (dos pasadas):
%   versión ciega para OJS:  xelatex articulo_RIIM.tex
%   versión con autores:     xelatex -jobname=articulo_RIIM_autores "\def\conautores{}% Artículo para la Revista Internacional de Investigación Multidisciplinaria (RIIM).
% Compilar con XeLaTeX (dos pasadas):
%   versión ciega para OJS:  xelatex articulo_RIIM.tex
%   versión con autores:     xelatex -jobname=articulo_RIIM_autores "\def\conautores{}% Artículo para la Revista Internacional de Investigación Multidisciplinaria (RIIM).
% Compilar con XeLaTeX (dos pasadas):
%   versión ciega para OJS:  xelatex articulo_RIIM.tex
%   versión con autores:     xelatex -jobname=articulo_RIIM_autores "\def\conautores{}\input{articulo_RIIM}"
\documentclass[10pt]{article}

% ── Revisión doble ciega ────────────────────────────────────────────
\newif\ifblind
\ifdefined\conautores \blindfalse \else \blindtrue \fi

% ── Página (plantilla RIIM: carta; márgenes 2.0 / 2.0 / 3.5 / 2.25 cm;
%    encabezado y pie a 1.25 cm del borde) ──────────────────────────
\usepackage[letterpaper, left=2cm, right=2cm, top=3.5cm, bottom=2.25cm,
            headheight=1.45cm, headsep=0.8cm, footskip=0.62cm]{geometry}

% ── Tipografía ──────────────────────────────────────────────────────
\usepackage{fontspec}
\setmainfont{Montserrat}[
  Path=fonts/, Extension=.ttf,
  UprightFont=*-Regular, ItalicFont=*-Italic,
  BoldFont=*-SemiBold, BoldItalicFont=*-SemiBoldItalic]
\newfontfamily\semibold{Montserrat}[
  Path=fonts/, Extension=.ttf,
  UprightFont=*-SemiBold, ItalicFont=*-SemiBoldItalic, Ligatures=TeX]
\newfontfamily\arialfont{Arial}[Ligatures=TeX]
\usepackage[spanish, mexico, es-tabla, es-nodecimaldot, es-noshorthands]{babel}
\spanishplainpercent
\usepackage{amsmath, amssymb}
\usepackage{xcolor}
\definecolor{riimazul}{HTML}{003D7C}
\definecolor{cfcfs}{HTML}{1E88E5}
\definecolor{crh}{HTML}{F59E0B}

\usepackage{graphicx}
\usepackage{booktabs}
\usepackage{tabularx}
\usepackage{float}
\usepackage{flafter}
% Permite que tablas y figuras compartan página con el texto en lugar de ir a páginas solo de flotantes.
\renewcommand{\topfraction}{0.85}
\renewcommand{\bottomfraction}{0.7}
\renewcommand{\textfraction}{0.1}
\renewcommand{\floatpagefraction}{0.8}
\setcounter{topnumber}{3}
\setcounter{bottomnumber}{2}
\setcounter{totalnumber}{4}
\setlength{\textfloatsep}{12pt plus 2pt minus 2pt}
\setlength{\floatsep}{10pt plus 2pt minus 2pt}
\setlength{\intextsep}{10pt plus 2pt minus 2pt}
\usepackage{tikz}
\usetikzlibrary{arrows.meta, positioning, fit, backgrounds}
\usepackage{fancyhdr}
\usepackage{caption}
\usepackage[hidelinks]{hyperref}

% Cuerpo: Montserrat 10, interlineado sencillo, justificado, sin sangría, 8 pt después.
\renewcommand{\normalsize}{\fontsize{10}{12.2}\selectfont}
\normalsize
\setlength{\parindent}{0pt}
\setlength{\parskip}{8pt plus 1pt}
\setlength{\abovedisplayskip}{6pt}
\setlength{\belowdisplayskip}{6pt}
\setlength{\abovedisplayshortskip}{4pt}
\setlength{\belowdisplayshortskip}{4pt}

% Títulos de sección sin numerar: Montserrat SemiBold 12 pt azul; subtítulos SemiBold itálica 11 pt.
\makeatletter
\setcounter{secnumdepth}{0}
\renewcommand{\section}{\@startsection{section}{1}{0pt}{8pt}{0.1pt}
  {\semibold\fontsize{12}{14.6}\selectfont\color{riimazul}\raggedright}}
\renewcommand{\subsection}{\@startsection{subsection}{2}{0pt}{8pt}{0.1pt}
  {\semibold\itshape\fontsize{11}{13.4}\selectfont\addfontfeatures{LetterSpace=6.8}\raggedright}}
\makeatother

% Pies de figura y de tabla: 9 pt, centrados, "Figura N." / "Tabla N.".
\captionsetup{font={footnotesize}, labelfont={footnotesize}, labelsep=period,
  justification=centering, skip=4pt}
\captionsetup[table]{position=top}
\renewcommand{\footnotesize}{\fontsize{9}{11}\selectfont}

% Encabezado: logotipo RIIM (imagen de la plantilla recortada al 59.3 % izquierdo, como en Word) a la izquierda y "RIIM, Vol 3, Num 2, 2026" (Arial 8) a la derecha.
% Pie: número de página centrado; ISSN a la izquierda y URL a la derecha (9 pt).
\pagestyle{fancy}
\fancyhf{}
\renewcommand{\headrulewidth}{0pt}
\renewcommand{\footrulewidth}{0pt}
\fancyhead[C]{\makebox[\textwidth]{%
  \raisebox{-0.5\height}{\includegraphics[trim=0 0 335.53bp 0, clip, width=7.37cm, height=1.40cm]{img/riim_header.png}}%
  \hfill\raisebox{-0.5\height}{\arialfont\fontsize{8}{9.6}\selectfont RIIM, Vol 3, Num 2, 2026}}}
\fancyfoot[C]{\parbox[t]{\textwidth}{\centering\fontsize{9}{11}\selectfont\thepage\\[1pt]
  ISSN 3122-4806\hfill http://orizaba.tecnm.mx/revistaojs}}

% Listas compactas.
\makeatletter
\renewcommand{\@listI}{\leftmargin=1.2em \topsep=0pt \partopsep=0pt \itemsep=2pt \parsep=0pt}
\let\@listi\@listI
\makeatother

% Referencias IEEE (Arial 10, sangría francesa de 1.27 cm).
\makeatletter
\renewenvironment{thebibliography}[1]{%
  \section{Referencias}
  \list{\@biblabel{\@arabic\c@enumiv}}{%
    \arialfont\fontsize{10}{12}\selectfont
    \settowidth\labelwidth{\@biblabel{#1}}%
    \leftmargin=1.27cm \labelwidth=1.27cm \labelsep=0pt \itemindent=0pt
    \itemsep=0pt \parsep=0pt \topsep=0pt
    \def\makelabel##1{##1\hfil}%
    \usecounter{enumiv}\let\p@enumiv\@empty
    \renewcommand\theenumiv{\@arabic\c@enumiv}}%
  \sloppy\clubpenalty4000\widowpenalty4000\sfcode`\.\@m}
  {\def\@noitemerr{\@latex@warning{Empty `thebibliography' environment}}\endlist}
\makeatother

\hypersetup{pdftitle={Logística de última milla en redes viales reales: MDRP reproducible en La Paz, B.C.S.},
  pdfauthor={}, pdfsubject={RIIM}, pdfkeywords={enrutamiento dinámico, horizonte rodante, OpenStreetMap, simulación}}
\ifblind\else\hypersetup{pdfauthor={Alan José Marcos Sánchez Martínez}}\fi

% Notación.
\newcommand{\CtD}{\text{CtD}}
\newcommand{\Cdisp}{\mathcal{C}_{disp}(t)}

% Diagrama de caja: \caja{x}{mín}{Q1}{mediana}{Q3}{máx}{color}
\newcommand{\caja}[7]{%
  \draw (#1,#2) -- (#1,#3) (#1,#5) -- (#1,#6);
  \draw (#1-0.14,#2) -- (#1+0.14,#2) (#1-0.14,#6) -- (#1+0.14,#6);
  \filldraw[fill=#7!75, draw=black!70] (#1-0.3,#3) rectangle (#1+0.3,#5);
  \draw[line width=1pt] (#1-0.3,#4) -- (#1+0.3,#4);}

\begin{document}

% ── Título, autores y área ──────────────────────────────────────────
\begin{center}
  \setlength{\parskip}{0pt}
  \vspace*{6pt}
  {\semibold\fontsize{16}{19.5}\selectfont\addfontfeatures{LetterSpace=-3.1}%
   LOGÍSTICA DE ÚLTIMA MILLA EN REDES VIALES REALES: MDRP REPRODUCIBLE EN LA PAZ, B.C.S.\par}
  \vspace{18pt}
  \ifblind\else
  {\itshape\fontsize{9}{11}\selectfont
   Sánchez Martínez, Alan José Marcos\textsuperscript{1*}\par
   \textsuperscript{1}Tecnológico Nacional de México / Instituto Tecnológico de La Paz, m16310322@lapaz.tecnm.mx\par}
  \fi
  {\fontsize{9}{11}\selectfont Área de participación: Sistemas Computacionales\par}
\end{center}

\section{Resumen}
Las plataformas de entrega de comida deciden sin conocer los pedidos futuros, si asignan cada orden
al repartidor libre más cercano o si esperan unos minutos para agrupar órdenes y coordinar la
asignación. Este trabajo evalúa en qué condiciones conviene coordinar en una ciudad media mexicana,
La Paz, B.C.S., y cuanto cambia la evaluación al usar la red vial real. Se construyó un simulador de
código abierto que obtiene los tiempos de viaje de OSRM sobre OpenStreetMap y compara una política
reactiva (FCFS) con una heurística de horizonte rodante (RH) en cuatro niveles de saturación de la
flota, con diez instancias sintéticas por nivel. Con una flota holgada, FCFS entrega 1.4~min antes;
con 26 o más pedidos por repartidor en la jornada, RH reduce el tiempo de entrega en todas las
instancias, hasta en 67\,\%. Las distancias euclidianas subestiman los tiempos de entrega y el
beneficio de coordinar.

{\semibold\fontsize{12}{14.6}\selectfont\color{riimazul}Palabras clave}: \textit{enrutamiento dinámico, horizonte rodante, OpenStreetMap, simulación}

\section{Abstract}
\begin{otherlanguage}{english}\frenchspacing
Food delivery platforms must decide, without knowing future orders, whether to assign each order to
the nearest idle courier or to wait a few minutes, batch orders and coordinate the assignment. This
work evaluates under which conditions coordination pays off in a mid-sized Mexican city, La Paz,
B.C.S., and how much the evaluation changes when the real road network is used. An open-source
simulator was built that obtains travel times from OSRM over OpenStreetMap and compares a reactive
policy (FCFS) with a rolling-horizon heuristic (RH) at four fleet saturation levels, with ten
synthetic instances per level. With a slack fleet, FCFS delivers 1.4~min earlier; with 26 or more
orders per courier over the day, RH reduces delivery time in every instance, by up to 67\,\%.
Euclidean distances underestimate both delivery times and the benefit of coordination.

{\semibold\fontsize{12}{14.6}\selectfont\color{riimazul}Key words}: \textit{dynamic routing, rolling horizon, OpenStreetMap, simulation}
\end{otherlanguage}

\section{Introducción}
La última milla, el tramo que va del punto de recolección al domicilio del cliente, es uno de los
eslabones mas costosos y difíciles de la logística urbana. En zonas residenciales concentra grandes
volúmenes de entregas con horarios estrechos y condiciones que cambian durante el día
\cite{Jazemi2023}. La entrega de comida a domicilio agrava estas exigencias: los pedidos llegan sin
aviso, la comida pierde calidad con cada minuto de espera y el cliente espera recibirla en menos de
una hora.

En este contexto, una plataforma debe decidir de forma continua qué repartidor atiende cada orden y
en qué momento. Este problema, formalizado como \textit{Meal Delivery Routing Problem} (MDRP) por
Reyes \textit{et al.} \cite{Reyes2018}, admite dos estrategias básicas. La primera es reactiva: cada
orden se asigna en cuanto aparece al repartidor libre más cercano (\textit{First-Come, First-Served},
FCFS). La segunda es coordinada: el sistema espera unos minutos, agrupa órdenes del mismo restaurante
y decide de manera conjunta qué repartidor atiende cada agrupamiento, como en la heurística de
horizonte rodante (\textit{Rolling Horizon}, RH) propuesta en \cite{Reyes2018}. Coordinar tiene un
costo, porque algunas órdenes esperan a que se forme el agrupamiento o a que se entreguen otras antes
que ellas, y un beneficio, porque cada viaje atiende más pedidos. El balance entre ambos depende de
cuántos repartidores hay frente a la demanda.

La evidencia disponible sobre este balance proviene de entornos muy distintos a una ciudad media. RH
se evaluó con instancias que reproducen áreas metropolitanas de Estados Unidos y con tiempos de viaje
calculados como la distancia euclidiana multiplicada por un factor constante \cite{Reyes2018}. En
instancias pequeñas resueltas a optimalidad, su \textit{Click-to-Door} promedio queda apenas un
4\,\% por encima del óptimo \cite{Yildiz2019}, pero no se ha reportado cómo cambia su ventaja frente
a una política reactiva conforme varía el tamaño de la flota. Otras formulaciones del problema
incorporan la incertidumbre en los tiempos de preparación \cite{Ulmer2021} o la estimación del tiempo
de llegada, cuya desviación respecto a lo prometido pesa en la satisfacción del cliente tanto como la
rapidez misma \cite{Hildebrandt2022}.

La forma de medir los tiempos de viaje añade una segunda fuente de error. Las distancias por la red
vial suelen ser entre 1.2 y 1.4 veces mayores que en línea recta, y en problemas con ventanas de
tiempo estrechas, como la entrega de comida, una ruta factible con distancias euclidianas puede
resultar infactible sobre la red real \cite{Boyaci2021}. Una evaluación que ignore la red puede
equivocarse no solo en los tiempos de entrega, sino también en la carga efectiva de la flota y, por
lo tanto, en el beneficio de coordinar. Los servicios comerciales de rutas resuelven la precisión,
pero su costo de licencia limita su uso en investigación \cite{Pereira2021} y su funcionamiento
cerrado dificulta reproducir los experimentos.

La Paz, B.C.S., ilustra esta brecha: es una ciudad media sin datos públicos de operaciones de entrega
de comida, cuya red vial, condicionada por la costa y por pocas avenidas principales, obliga a
recorridos más largos que la línea recta. De ahí surge la pregunta que guía este trabajo:
\textit{en una ciudad media, ¿en qué condiciones de saturación de la flota conviene coordinar los
despachos en lugar de asignar cada pedido al llegar, y cuánto cambia esa evaluación al usar la red
vial real en lugar de distancias euclidianas?} La pregunta se divide en tres: cómo varía la ventaja
de RH con la saturación y a partir de qué nivel cambia la política más conveniente (PI1); cuánto más largos son los
trayectos por la red y cuánto cambia la comparación al ignorarla (PI2); y si la evaluación puede
hacerse de forma reproducible, con recursos modestos y solo con herramientas de código abierto (PI3).

Las contribuciones corresponden a esas tres preguntas: (i)~una comparación pareada de FCFS y RH en
cuatro niveles de saturación, con prueba estadística por instancia, que ubica el punto a partir del
cual coordinar compensa; (ii)~la medición de la circuidad de la red urbana de La Paz y dos escenarios
de control con distancias euclidianas que separan el efecto de la red del efecto de la política; y
(iii)~un entorno de simulación reproducible, construido sobre OpenStreetMap y OSRM, que ejecuta un
día de operación en alrededor de un minuto en una computadora portátil. El estudio no propone un
algoritmo nuevo ni compara RH con políticas del estado del arte: su objetivo es caracterizar cuándo
conviene coordinar en un entorno realista.

\section{Metodología}

\subsection{Formulación del problema y métricas}
Se adopta la formulación del MDRP de \cite{Reyes2018}. Un conjunto $\mathcal{R}$ de restaurantes con
ubicación fija recibe órdenes $o \in \mathcal{O}$, cada una con restaurante de origen $r_o$, instante
de colocación $a_o$, instante en que la comida está lista $e_o \geq a_o$ y punto de entrega $\ell_o$.
La orden solo se conoce a partir de $a_o$. Una flota $\mathcal{C}$ de repartidores, con turno
$[e_c, l_c]$, recoge y entrega las órdenes. Órdenes del mismo restaurante pueden formar un
agrupamiento $B$ que se recoge en una sola visita y se entrega en una ruta de varias paradas. Los
tiempos de servicio son $s_r$ en el restaurante, una vez por agrupamiento, y $s_o$ con cada cliente.

La métrica principal es el \textit{Click-to-Door} de cada orden,
$\CtD_o = t_{delivery,o} - a_o$. A partir de ella se reportan su media y su percentil 90, y el
cumplimiento del nivel de servicio (SLA, \textit{Service Level Agreement}), entendido como el
porcentaje de órdenes entregadas dentro del objetivo $\tau$:
\begin{equation}
  \text{SLA} = \frac{|\{o \in \mathcal{O}_{del} : \CtD_o \leq \tau\}|}{|\mathcal{O}_{del}|} \times 100,
  \label{eq:sla}
\end{equation}
donde $\mathcal{O}_{del}$ es el conjunto de órdenes entregadas. Como ambas métricas se calculan
sobre órdenes entregadas, se reporta también el porcentaje de órdenes sin entregar al cierre. Se
complementan con el tiempo que la comida lista espera a ser recogida (RtP, \textit{Ready-to-Pickup}),
$\text{RtP}_o = t_{pickup,o} - e_o$, la
distancia recorrida por pedido, la utilización de la flota (fracción del turno en conducción o en
servicio), el tamaño de agrupamiento y el costo por pedido, con un pago de $p_1$ por orden entregada
y un mínimo garantizado de $p_2$ por hora de turno \cite{Reyes2018}.

\subsection{Políticas de despacho}
El simulador avanza en epocas de $f = 5$~min para ambas políticas: en cada época, FCFS asigna una por
una las órdenes nuevas y RH resuelve su optimización. Las dos calculan los instantes de recolección y
entrega con las mismas reglas. Un agrupamiento $B$ solo puede recogerse cuando ya está lista la
última de sus órdenes, es decir, a partir del instante $e_B = \max_{o\in B} e_o$. El repartidor $c$
llega al restaurante en el instante $t_{arr,r}$ y, como en \cite{Reyes2018}, cada tiempo de servicio
se reparte mitad antes y mitad después del momento de recoger o de entregar. Si la comida ya está
lista, la recolección ocurre al cumplirse la primera mitad del servicio en el restaurante; si no, el
repartidor espera hasta que esté lista. Por eso el instante de recolección es
$t_{pickup}(c,B) = \max\bigl(e_B,\; t_{arr,r} + s_r/2\bigr)$. Después, el repartidor sale del
restaurante al terminar la segunda mitad del servicio, recorre el tramo hasta el primer cliente y le
entrega al cumplirse la primera mitad del servicio con él; cada entrega siguiente se calcula igual a
partir de la anterior.

\textbf{FCFS (línea base reactiva).} En cada época, las órdenes reveladas sin repartidor se procesan
por orden de colocación. Cada una se asigna al repartidor libre más cercano en línea recta al
restaurante (esta distancia solo sirve para escoger al repartidor); la ruta repartidor $\to$ restaurante
$\to$ cliente y sus tiempos se obtienen de OSRM, y el repartidor queda comprometido de inmediato. Cada
viaje lleva una sola orden.

\textbf{RH (política coordinada).} Sigue el esquema de \cite{Reyes2018} y en cada época $t$ ejecuta
cinco pasos. Primero, delimita las órdenes que entran a la optimización: de las órdenes ya reveladas y
aún no entregadas, $\mathcal{O}_t$, solo considera las que estarán listas dentro del horizonte de
asignación $\Delta_u$, es decir, $U_t = \{o \in \mathcal{O}_t : e_o \leq t + \Delta_u\}$; planear
órdenes cuya comida tardará mucho en estar lista sería poco útil en un entorno que cambia cada pocos
minutos.

Segundo, estima cuántas órdenes conviene juntar en cada viaje. El tamaño objetivo de agrupamiento
$Z_t$ compara la demanda próxima con la oferta próxima: el numerador cuenta las órdenes que estarán
listas en los siguientes $\Delta_1$ minutos, y el denominador, los repartidores en turno, libres u
ocupados, que quedarán disponibles en los siguientes $\Delta_2$ minutos, donde $e^{avail}_c$ es el
instante en que el repartidor $c$ termina su ruta actual:
\begin{equation}
  Z_t = \left\lceil \frac{|\{o \in \mathcal{O}_t : e_o \leq t + \Delta_1\}|}
                        {|\{c \in \mathcal{C}_{turno}(t) : e^{avail}_c \leq t + \Delta_2\}|} \right\rceil.
  \label{eq:zt}
\end{equation}
Si hay tantos repartidores como ordenes, $Z_t = 1$ y cada orden tiende a viajar sola; si la flota se
satura, $Z_t$ crece y se forman agrupamientos más grandes. $Z_t$ se acota en 5, y ningún agrupamiento
puede superar $B_{max} = 4$ órdenes.

Tercero, forma los agrupamientos de cada restaurante $r$. Sus órdenes $U_{t,r}$ se reparten en
\[
  m_r = \min\bigl(\lceil |U_{t,r}|/Z_t \rceil,\; |\Cdisp|\bigr)
\]
agrupamientos, de modo que cada uno reciba alrededor de $Z_t$ órdenes sin crear más agrupamientos que
repartidores disponibles. Las órdenes se toman por \textit{ready time}, y cada una se inserta en el
agrupamiento y la posición de la ruta que menos incrementan el costo
\begin{equation}
  \text{costo}(B, o, p) = T_{ruta}(B \oplus_p o) + \beta\, T_{servicio}(B \oplus_p o),
  \label{eq:costo}
\end{equation}
donde $B \oplus_p o$ es el agrupamiento con la orden $o$ insertada en la posición $p$, $T_{ruta}$ es
la duración de la ruta según OSRM y $T_{servicio}$ el tiempo de servicio acumulado. El peso $\beta$, que en
\cite{Reyes2018} representa la pérdida de frescura y se calibra en $[0, 1]$, se fija en 0.1 para que la
duración de la ruta domine el costo. Después, dos iteraciones de extracción y reinserción mueven
cada orden al agrupamiento donde su costo es menor.

Cuarto, asigna agrupamientos a repartidores. Los agrupamientos se clasifican en tres grupos de
prioridad: urgentes, cuando ningún repartidor puede entregar su orden más antigua dentro del objetivo
$\tau$; retrasados, cuando el objetivo aún es alcanzable pero ningún repartidor llega al restaurante
antes de que la comida esté lista; y normales. Los grupos se resuelven en ese orden, para que los
urgentes obtengan primero a los repartidores más cercanos. En cada grupo se resuelve un problema de
asignación lineal, en el que cada repartidor recibe a lo más un agrupamiento, con el peso
\begin{equation}
  w_{c,B} = \frac{|B|}{T_{delivery}(c,B)} - \theta\,\bigl(t_{pickup}(c,B) - e_B\bigr) - P_g.
  \label{eq:peso}
\end{equation}
El primer término mide la eficiencia del viaje en órdenes entregadas por hora, donde $T_{delivery}(c,B)$
es el tiempo desde que $c$ queda libre hasta la última entrega de $B$. El segundo penaliza, con
$\theta = 0.1$, los minutos que la comida espera lista en el restaurante, y $P_g$ resta 100 si el par
excede el máximo absoluto $\tau_{max}$ y 50 si la recolección se retrasa. La asignación de máximo peso
se obtiene de forma exacta con un algoritmo de tipo húngaro \cite{Kuhn1955}.

Quinto, el compromiso en dos etapas de \cite{Reyes2018} decide qué parte de la asignación se
ejecuta. Si el repartidor llega al restaurante y todas las órdenes estarán listas dentro de la época,
o si alguna orden lleva más de $\Delta_x$ minutos lista, se fija la ruta completa. En otro caso, el
repartidor solo se dirige al restaurante, como un preposicionamiento, y en la época siguiente el
agrupamiento puede crecer o cambiar de repartidor con las órdenes que se hayan revelado.

\subsection{Entorno de simulación reproducible}
La Figura~\ref{fig:arquitectura} resume el entorno. El extracto de OpenStreetMap de Baja California
Sur, modelo de datos viales abierto y de cobertura mundial \cite{Boeing2022}, se preprocesa con el
algoritmo \textit{Contraction Hierarchies} \cite{Geisberger2008}, que permite resolver cada consulta
de camino mínimo en milisegundos. El servidor OSRM \cite{LuxenVetter2011}, versión 5.27.1, atiende
las consultas de ruta del simulador por HTTP; su imagen oficial de Docker permite reproducir el
entorno en cualquier sistema operativo. El simulador, escrito en Python, avanza por épocas, guarda en
memoria las rutas ya consultadas y produce, para cada corrida, los registros por orden y por
repartidor de los que se derivan los indicadores.

\begin{figure}[htbp]
\centering
\begin{tikzpicture}[
  font=\footnotesize,
  caja/.style={draw=black!70, rounded corners=2pt, align=center, minimum height=1.05cm,
               text width=#1, fill=white},
  caja/.default=2.8cm,
  flecha/.style={-{Stealth[length=5pt]}, thick, black!70},
  capa/.style={draw=#1!70!black, rounded corners=4pt, fill=#1!8, inner xsep=7pt, inner ysep=7pt}]
  \node[caja=2.8cm] (osm)   at (0,0)      {Extracto\\OpenStreetMap};
  \node[caja=3.0cm] (ch)    at (3.55,0)    {Preprocesamiento\\\textit{Contraction}\\\textit{Hierarchies}};
  \node[caja=2.8cm] (osrm)  at (7.1,0)    {Servidor OSRM\\(API de rutas)};
  \node[caja=2.8cm] (gen)   at (0,-2.3)   {Generador de\\instancias (PPNH)};
  \node[caja=2.8cm] (sim)   at (3.55,-2.3) {Simulador por\\épocas ($f = 5$ min)};
  \node[caja=2.8cm] (pol)   at (7.1,-2.3) {Política\\FCFS o RH};
  \node[caja=2.8cm] (kpi)   at (10.9,-2.3) {Registros por orden\\y 21 indicadores};
  \node[caja=2.8cm] (est)   at (10.9,0)   {Análisis estadístico\\(Wilcoxon, Friedman)};
  \draw[flecha] (osm) -- (ch);
  \draw[flecha] (ch) -- (osrm);
  \draw[flecha] (gen) -- (sim);
  \draw[flecha, {Stealth[length=5pt]}-{Stealth[length=5pt]}] (sim) -- (pol);
  \draw[flecha, dashed, {Stealth[length=5pt]}-{Stealth[length=5pt]}] (osrm) --
     node[right, font=\scriptsize, align=left] {tiempos de\\viaje (HTTP)} (pol);
  \draw[flecha] (pol) -- (kpi);
  \draw[flecha] (kpi) -- (est);
  \begin{scope}[on background layer]
    \node[capa=blue, fit=(osm)(ch)(osrm)] {};
    \node[capa=green, fit=(gen)(sim)(pol)] {};
    \node[capa=orange, fit=(kpi)(est)] {};
  \end{scope}
\end{tikzpicture}
\caption{Arquitectura del entorno: enrutamiento sobre la red vial (azul), simulación (verde) y
análisis (naranja).}
\label{fig:arquitectura}
\end{figure}

\subsection{Generación de instancias}
Al no existir datos públicos de entregas de comida en La Paz, las instancias se generan de forma
sintética. Las llegadas siguen un proceso de Poisson no homogéneo con intensidad bimodal
\begin{equation}
  \lambda(t) = \lambda_0
  + A_m \exp\!\left(-\frac{(t-\mu_m)^2}{2\sigma^2}\right)
  + A_e \exp\!\left(-\frac{(t-\mu_e)^2}{2\sigma^2}\right),
  \label{eq:lambda}
\end{equation}
siguiendo el patrón bimodal observado en entregas de última milla \cite{Wu2024}. A falta de datos
locales, los parámetros son supuestos: ventana operativa de 08:00 a 19:00, picos en $\mu_m = $ 09:00
y $\mu_e = $ 17:00, $\lambda_0 = 0.6$ órdenes/min, $A_m = 2.4$, $A_e = 2.0$ y $\sigma = 45$~min. El perfil se escala para
obtener alrededor de mil órdenes por día (entre 988 y 1\,047 según la semilla). El tiempo de
preparación es $\max(5, \mathcal{N}(12, 3^2))$~min. Los 25 restaurantes y los clientes se ubican
por muestreo uniforme dentro de un polígono de 8.4~km$^2$ del centro urbano, y cada cliente se asocia
a un restaurante situado a entre 2 y 5~km en línea recta. Todos los repartidores comparten un turno
que inicia 15~min antes del primer pedido y termina una hora después del último \textit{ready time},
y parten del centroide de los restaurantes. La flota se dimensiona a partir de la razón de
saturación $\rho$, el número objetivo de pedidos por repartidor en la jornada:
\begin{equation}
  |\mathcal{C}| = \max\!\left(n_{\min},\; \left\lceil |\mathcal{O}| / \rho \right\rceil\right),
  \qquad n_{\min} = 10.
  \label{eq:flota}
\end{equation}
Fijar $\rho$, y no el número de repartidores, mantiene la misma carga relativa en todas las semillas
de un escenario aunque el número de órdenes varíe.

\subsection{Diseño experimental}
La Tabla~\ref{tab:parametros} reúne los parámetros, comunes a ambas políticas, y los cuatro
escenarios de saturación. Cada semilla (2025 a 2034) genera una instancia, y sobre ella se ejecutan
FCFS y RH con la misma flota; como el simulador no tiene componentes aleatorios propios, lo único que
cambia dentro de cada par es la política. Esto da 80 corridas sobre la red vial. Para evitar la
pseudorreplicación, cada métrica se promedia por semilla y la comparación se hace con la prueba de
rangos con signo de Wilcoxon pareada \cite{Wilcoxon1945} sobre las diez diferencias de cada
escenario, con intervalo de confianza $t$ al 95\,\% y la correlación biserial de rangos $r_{rb}$ como
tamaño de efecto. La prueba de Friedman \cite{Friedman1937} sobre la matriz de diferencias
$10 \times 4$ evalúa si el efecto cambia entre escenarios.

\begin{table}[htbp]
\centering
\footnotesize
\caption{Parámetros del simulador y escenarios de saturación.}
\label{tab:parametros}
\setlength{\tabcolsep}{4pt}
\begin{tabular}{@{}l c c @{\hspace{0.4cm}} l c c@{}}
\toprule
Parámetro & Símbolo & Valor & Parámetro & Símbolo & Valor \\
\midrule
Frecuencia de optimización & $f$ & 5 min & Penalización de frescura & $\beta$ & 0.1 \\
Horizonte de asignación & $\Delta_u$ & 20 min & Penalización por espera & $\theta$ & 0.1 \\
Anticipación (demanda, oferta) & $\Delta_1$, $\Delta_2$ & 20 min & Espera para forzar compromiso & $\Delta_x$ & 10 min \\
CtD objetivo / máximo & $\tau$, $\tau_{max}$ & 40/90 min & Tamaño máximo de agrupamiento & $B_{max}$ & 4 \\
Tiempos de servicio & $s_r$, $s_o$ & 4 min & Pago por orden / mínimo por hora & $p_1$, $p_2$ & \$10 / \$15 \\
\midrule
\multicolumn{6}{@{}l}{Escenarios: S1, $\rho = 18$ (55--59 repartidores); S2, $\rho = 26$ (38--41); S3, $\rho = 35$ (29--30); S4, $\rho = 52$ (19--21).} \\
\bottomrule
\end{tabular}
\end{table}

Para medir el efecto de la red vial se calculó la circuidad $\kappa = d_{OSRM}/d_{recta}$ sobre todos
los trayectos restaurante--cliente, y el experimento completo se repitió con dos controles en los que
el tiempo de viaje es la distancia en línea recta (Haversine) a velocidad constante: la velocidad
media que OSRM asigna a esos trayectos (651~m/min), que deja como única diferencia la geometría de
las rutas, y esa velocidad dividida entre $\kappa$ (472~m/min), que iguala la duración media de los
viajes. Cada control suma otras 80 corridas.

\section{Resultados y discusión}

\subsection{Calidad de servicio y cambio de la política más conveniente}
La Tabla~\ref{tab:calidad} muestra los indicadores de calidad de servicio. Con la flota holgada de S1,
FCFS entrega en promedio 1.4~min antes que RH: con repartidores de sobra, llevar cada pedido directo
del restaurante al cliente es más rápido que esperar a formar un agrupamiento, y en un agrupamiento
cada entrega espera a las anteriores. Aun así, el cumplimiento del SLA es prácticamente igual en ambas
políticas. A partir de S2 la relación se invierte: RH reduce el \textit{Click-to-Door} medio en
36.8\,\%, el percentil 90 en 46.0\,\% y eleva el cumplimiento del SLA en 31.5 puntos porcentuales. La
ventaja crece en S3 y S4, donde FCFS deja sin entregar al cierre el 11.3\,\% y el 39.9\,\% de los
pedidos, mientras que RH entrega prácticamente todos. Como el CtD y el SLA se calculan sobre órdenes
entregadas, la tabla favorece a FCFS en esos escenarios: medido sobre todas las órdenes generadas, su
cumplimiento cae a 6.4\,\% en S3 y 3.8\,\% en S4, frente a 57.6\,\% y 42.4\,\% bajo RH.

\begin{table}[htbp]
\centering
\footnotesize
\caption{Calidad de servicio, promedio de 10 semillas. $\Delta$ es el cambio de RH respecto a FCFS.}
\label{tab:calidad}
\setlength{\tabcolsep}{3pt}
\begin{tabular*}{\textwidth}{@{\extracolsep{\fill}}l rrr rrr rrr rr@{}}
\toprule
& \multicolumn{3}{c}{CtD medio (minutos)} & \multicolumn{3}{c}{CtD percentil 90 (minutos)} & \multicolumn{3}{c}{SLA (\%)} & \multicolumn{2}{c}{Sin entregar (\%)} \\
\cmidrule(lr){2-4} \cmidrule(lr){5-7} \cmidrule(lr){8-10} \cmidrule(l){11-12}
Esc. & FCFS & RH & $\Delta$ (\%) & FCFS & RH & $\Delta$ (\%) & FCFS & RH & $\Delta$ (pp) & FCFS & RH \\
\midrule
S1 & 28.93  & 30.31 & $+4.8$  & 43.33  & 43.80  & $+1.1$  & 84.6 & 83.4 & $-1.1$  & 0.0  & 0.0 \\
S2 & 56.07  & 35.42 & $-36.8$ & 97.91  & 52.87  & $-46.0$ & 38.1 & 69.6 & $+31.5$ & 0.1  & 0.0 \\
S3 & 108.49 & 45.36 & $-58.2$ & 169.54 & 79.10  & $-53.3$ & 7.3  & 57.6 & $+50.3$ & 11.3 & 0.0 \\
S4 & 224.88 & 73.63 & $-67.3$ & 321.10 & 186.18 & $-42.0$ & 6.2  & 42.5 & $+36.3$ & 39.9 & 0.2 \\
\bottomrule
\end{tabular*}
\end{table}

La Figura~\ref{fig:boxplot} muestra la distribución completa del \textit{Click-to-Door}. En S1 ambas
cajas son casi idénticas y quedan por debajo del objetivo de 40~min; desde S2, la distribución de
FCFS se alarga hacia valores de varias horas, mientras que la de RH se mantiene compacta y con la
mediana cerca del objetivo.

\begin{figure}[htbp]
\centering
\begin{tikzpicture}[y=0.0100cm, x=1cm, font=\footnotesize]
  \foreach \p/\esc in {0/S1 ($\rho=18$), 1/S2 ($\rho=26$), 2/S3 ($\rho=35$), 3/S4 ($\rho=52$)} {
    \begin{scope}[xshift=\p*4.15cm]
      \draw[black!60] (0,0) rectangle (3.6,390);
      \foreach \y in {0,100,200,300} \draw[black!12] (0,\y) -- (3.6,\y);
      \draw[red!70, dashed] (0,40) -- (3.6,40);
      \node[above] at (1.8,390) {\esc};
      \node[below] at (1.1,0) {FCFS};
      \node[below] at (2.5,0) {RH};
    \end{scope}}
  \foreach \y in {0,100,200,300} \node[left] at (0,\y) {\y};
  \node[rotate=90] at (-1.25,195) {\textit{Click-to-Door} (minutos)};
  % S1
  \caja{1.1}{10.9}{21.7}{26.4}{35.1}{55.1}{cfcfs}
  \caja{2.5}{11.4}{22.8}{28.5}{36.1}{56.1}{crh}
  % S2
  \begin{scope}[xshift=4.15cm]
  \caja{1.1}{11.9}{27.3}{54.3}{80.6}{121.7}{cfcfs}
  \caja{2.5}{11.4}{24.5}{32.5}{42.2}{68.7}{crh}
  \end{scope}
  % S3
  \begin{scope}[xshift=8.3cm]
  \caja{1.1}{11.9}{72.1}{109.1}{148.2}{199.9}{cfcfs}
  \caja{2.5}{11.4}{26.6}{36.5}{50.3}{85.8}{crh}
  \end{scope}
  % S4
  \begin{scope}[xshift=12.45cm]
  \caja{1.1}{14.9}{141.0}{266.0}{306.5}{367.5}{cfcfs}
  \caja{2.5}{12.3}{31.1}{43.8}{79.3}{151.5}{crh}
  \end{scope}
\end{tikzpicture}
\caption{Distribución del \textit{Click-to-Door} de las órdenes entregadas en las diez semillas
(cajas: cuartiles; bigotes: 1.5 veces el rango intercuartílico). La línea discontinua marca el
objetivo $\tau = 40$~min.}
\label{fig:boxplot}
\end{figure}

La Tabla~\ref{tab:wilcoxon} confirma, a nivel de instancia, el cambio de la política más conveniente. En S2, S3 y S4, RH
supera a FCFS en las diez semillas, tanto en CtD como en SLA, con $p = 0.0020$, el valor mínimo que
admite la prueba con $n = 10$, y $|r_{rb}| = 1$. En S1 la ventaja de FCFS en CtD es significativa
pero pequeña, y la diferencia en SLA no lo es. La prueba de Friedman confirma que la magnitud del
efecto cambia con la saturación ($\chi^2_F = 30.0$, $p < 0.001$). La Figura~\ref{fig:saturacion}
resume la respuesta a PI1: las curvas de CtD se cruzan entre $\rho = 18$ y $\rho = 26$, y por encima
de ese punto la ventaja de coordinar crece con la carga.

\begin{table}[htbp]
\centering
\footnotesize
\caption{Wilcoxon pareada por semilla ($n = 10$), $\bar{\Delta} = \bar{x}_{FCFS} - \bar{x}_{RH}$: favorecen a RH
$\bar{\Delta} > 0$ en CtD y $\bar{\Delta} < 0$ en SLA.}
\label{tab:wilcoxon}
\begin{tabular}{@{}ll rr l c r@{}}
\toprule
Esc. & Métrica & $\bar{x}_{FCFS}$ & $\bar{x}_{RH}$ & $\bar{\Delta}$ [IC 95\,\%] & $p$ & $r_{rb}$ \\
\midrule
S1 & CtD (minutos) & 28.93  & 30.31 & $-1.38$ [$-2.35$, $-0.41$]        & 0.0137 & $-0.85$ \\
S1 & SLA (\%)  & 84.6   & 83.4  & $+1.1$ [$-1.8$, $+4.1$]           & 0.2324 & $+0.45$ \\
S2 & CtD (minutos) & 56.07  & 35.42 & $+20.64$ [$+18.98$, $+22.31$]     & 0.0020 & $+1.00$ \\
S2 & SLA (\%)  & 38.1   & 69.6  & $-31.5$ [$-33.7$, $-29.2$]        & 0.0020 & $-1.00$ \\
S3 & CtD (minutos) & 108.49 & 45.36 & $+63.13$ [$+56.14$, $+70.13$]     & 0.0020 & $+1.00$ \\
S3 & SLA (\%)  & 7.3    & 57.6  & $-50.3$ [$-52.6$, $-48.0$]        & 0.0020 & $-1.00$ \\
S4 & CtD (minutos) & 224.88 & 73.63 & $+151.24$ [$+143.71$, $+158.78$]  & 0.0020 & $+1.00$ \\
S4 & SLA (\%)  & 6.2    & 42.5  & $-36.3$ [$-37.2$, $-35.3$]        & 0.0020 & $-1.00$ \\
\bottomrule
\end{tabular}
\end{table}

\begin{figure}[htbp]
\centering
\begin{tikzpicture}[font=\footnotesize]
  % Panel izquierdo: CtD medio
  \begin{scope}[x=0.16cm, y=0.018cm]
    \draw[black!60] (14,0) rectangle (56,240);
    \foreach \y in {0,50,100,150,200} {\draw[black!12] (14,\y) -- (56,\y); \node[left] at (14,\y) {\y};}
    \foreach \x in {18,26,35,52} {\draw[black!12] (\x,0) -- (\x,240); \node[below] at (\x,0) {\x};}
    \draw[red!70, dashed] (14,40) -- (56,40);
    \draw[cfcfs, very thick] (18,28.93) -- (26,56.07) -- (35,108.49) -- (52,224.88);
    \foreach \x/\y in {18/28.93, 26/56.07, 35/108.49, 52/224.88} \fill[cfcfs] (\x,\y) circle (2pt);
    \draw[crh, very thick] (18,30.31) -- (26,35.42) -- (35,45.36) -- (52,73.63);
    \foreach \x/\y in {18/30.31, 26/35.42, 35/45.36, 52/73.63} \fill[crh] (\x-0.45,\y-4.5) rectangle (\x+0.45,\y+4.5);
    \node[rotate=90] at (8.5,120) {CtD medio (minutos)};
    \node at (35,-30) {Pedidos por repartidor en la jornada ($\rho$)};
  \end{scope}
  % Panel derecho: SLA
  \begin{scope}[xshift=8.6cm, x=0.16cm, y=0.043cm]
    \draw[black!60] (14,0) rectangle (56,100);
    \foreach \y in {0,20,40,60,80,100} {\draw[black!12] (14,\y) -- (56,\y); \node[left] at (14,\y) {\y};}
    \foreach \x in {18,26,35,52} {\draw[black!12] (\x,0) -- (\x,100); \node[below] at (\x,0) {\x};}
    \draw[cfcfs, very thick] (18,84.6) -- (26,38.1) -- (35,7.3) -- (52,6.2);
    \foreach \x/\y in {18/84.6, 26/38.1, 35/7.3, 52/6.2} \fill[cfcfs] (\x,\y) circle (2pt);
    \draw[crh, very thick] (18,83.4) -- (26,69.6) -- (35,57.6) -- (52,42.5);
    \foreach \x/\y in {18/83.4, 26/69.6, 35/57.6, 52/42.5} \fill[crh] (\x-0.45,\y-1.9) rectangle (\x+0.45,\y+1.9);
    \node[rotate=90] at (8.5,50) {Cumplimiento del SLA (\%)};
    \node at (35,-12.5) {Pedidos por repartidor en la jornada ($\rho$)};
    % Leyenda
    \draw[cfcfs, very thick] (34,93) -- (38,93); \fill[cfcfs] (36,93) circle (2pt); \node[right] at (38,93) {FCFS};
    \draw[crh, very thick] (46,93) -- (50,93); \fill[crh] (47.55,91.1) rectangle (48.45,94.9); \node[right] at (50,93) {RH};
  \end{scope}
\end{tikzpicture}
\caption{\textit{Click-to-Door} medio y cumplimiento del SLA según la saturación de la flota. La línea
discontinua marca el objetivo $\tau = 40$~min.}
\label{fig:saturacion}
\end{figure}

\subsection{Eficiencia operativa y agrupamiento}
La Tabla~\ref{tab:eficiencia} explica el origen de la ventaja de RH. Aun en S1, RH reduce la espera
de la comida lista en el restaurante (RtP), porque en los picos de demanda la flota se ocupa por completo y
FCFS, con un pedido por viaje, acumula órdenes sin recoger, mientras que RH recoge varias en una sola
visita. Al agrupar, RH recorre entre 13.5\,\% y 49.4\,\% menos distancia por pedido y mantiene una
utilización de la flota menor que FCFS en S1--S3: con la misma flota, entrega lo mismo con menos
tiempo de conducción y, en S4, convierte ese margen en pedidos entregados. El tamaño de agrupamiento
crece con la saturación; en S1 una cuarta parte de los pedidos (25.4\,\%) viaja sola y en S4 el
80.9\,\% viaja en agrupamientos de tres o cuatro. El costo por pedido solo difiere en S1, donde muchos
repartidores reciben el pago mínimo y RH concentra las entregas en menos de ellos; desde S2 ambos
convergen al pago por orden.

\begin{table}[htbp]
\centering
\footnotesize
\caption{Eficiencia operativa y tamaño del agrupamiento en que viaja cada pedido (promedio de 10 semillas).}
\label{tab:eficiencia}
\setlength{\tabcolsep}{3pt}
\begin{tabular*}{\textwidth}{@{\extracolsep{\fill}}l rr rr rr r rr@{}}
\toprule
& \multicolumn{2}{c}{RtP medio (minutos)} & \multicolumn{2}{c}{Distancia/pedido (km)} & \multicolumn{2}{c}{Utilización (\%)} & Agrupamiento & \multicolumn{2}{c}{Costo/pedido (\$)} \\
\cmidrule(lr){2-3} \cmidrule(lr){4-5} \cmidrule(lr){6-7} \cmidrule(lr){8-8} \cmidrule(l){9-10}
Esc. & FCFS & RH & FCFS & RH & FCFS & RH & RH & FCFS & RH \\
\midrule
S1 & 7.78   & 4.94  & 4.71 & 4.08 & 37.7 & 31.2 & 2.54 & 10.80 & 11.59 \\
S2 & 34.92  & 9.33  & 4.88 & 3.60 & 55.2 & 41.4 & 2.80 & 10.00 & 10.03 \\
S3 & 87.35  & 18.62 & 5.07 & 3.17 & 66.8 & 51.4 & 3.03 & 10.00 & 10.00 \\
S4 & 203.75 & 46.20 & 5.20 & 2.63 & 67.3 & 67.8 & 3.30 & 10.00 & 10.00 \\
\bottomrule
\end{tabular*}
\end{table}

\subsection{Efecto de la red vial}
Sobre los 10\,184 trayectos restaurante--cliente de las diez instancias, la distancia por la red es
en promedio 3.30~km y en línea recta 2.39~km: los trayectos reales son 38\,\% más largos (circuidad
$\kappa = 1.38$), dentro del rango de 1.2 a 1.4 que reporta la literatura \cite{Boyaci2021}.

Para medir cuánto cambia la comparación, se repitió el experimento con tiempos de viaje en línea
recta. A la velocidad media de la red (651~m/min, primer control), ambas políticas parecen más
rápidas: en S2, el CtD de FCFS baja de 56.1 a 41.8~min y el de RH, de 35.4 a 30.3~min. Como FCFS
mejora más, la ventaja de RH se reduce en los cuatro escenarios ($p = 0.002$); en S3, por ejemplo,
pasa de 63.1 a 37.1~min. El segundo control baja la velocidad hasta que los viajes en línea recta
duran en promedio lo mismo que por la red ($651/1.38 \approx 472$~m/min); con él desaparece entre el
62 y el 84\,\% de esa reducción. La mayor parte del efecto se debe, por tanto, a que los trayectos
reales son más largos y cargan más a la flota, y el resto, a la forma de la red ($p \leq 0.004$).
Ignorar la red hace parecer más rápidas a ambas políticas y subestima el beneficio de coordinar, lo
que responde a la pregunta PI2.

\subsection{Reproducibilidad y costo de cómputo}
Toda la cadena (OpenStreetMap, OSRM y Python) es de código abierto, las instancias se regeneran a
partir de semillas fijas y cada corrida pasa verificaciones automáticas de integridad: ninguna orden
se recoge antes de estar lista ni se entrega antes de recogerse, no hay entregas duplicadas y ninguna
ruta recurrió a la estimación de respaldo en línea recta (de haber fallado OSRM). Con 1\,038 órdenes, simular un día toma
entre 1.8 y 2.7~segundos con FCFS y entre 57 y 80~segundos con RH, incluidas las consultas a OSRM, con menos de
225~MB de memoria por proceso, en una computadora portátil con procesador Apple M1 Pro y 16~GB de
RAM. Esto responde a PI3: la evaluación es reproducible y no requiere infraestructura especializada.

\subsection{Limitaciones}
El reloj del simulador avanza en épocas de 5~min. En cada época primero se asignan los pedidos y
después se registran las rutas terminadas, de modo que un repartidor que termina a mitad de una época
queda libre hasta la siguiente y espera entre 5 y 10~min su nuevo viaje. Como esta espera ocurre una
vez por viaje, perjudica más a FCFS, que hace un viaje por pedido. Para comprobar que no determina el
resultado, se repitieron las 80 corridas sobre la red vial con el orden inverso (primero se registran
las rutas terminadas y luego se asigna): FCFS sigue ganando en S1 y RH en S2--S4 en todas las
semillas, y el cruce permanece entre S1 y S2. La reducción del CtD baja a 14.6\,\% en S2, 41.4\,\% en
S3 y 63.0\,\% en S4, por lo que las magnitudes dependen de este detalle de implementación, pero las
conclusiones se mantienen. Además, la demanda es
sintética: sigue el patrón bimodal de la recolección de paquetes en LaDe con parámetros supuestos, no
calibrados con la demanda local de comida, y el tipo de ciudad influye en el
comportamiento de la última milla \cite{MunozVillamizar2021}. Los tiempos de preparación y de
servicio son homogéneos, la red se considera sin congestión y las ubicaciones se muestrean de forma
uniforme dentro del dominio urbano.

\section{Trabajo a futuro}
Quedan pendientes seis líneas de trabajo: calibrar la demanda, los tiempos de preparación y la
ubicación de restaurantes y clientes con datos locales de una plataforma de entregas o de encuestas;
explorar valores de $\rho$ entre 18 y 26 para ubicar con precisión el punto de cruce y evaluar
políticas híbridas que activen la coordinación solo cuando la carga lo justifique; incorporar
velocidades dependientes de la hora para representar la congestión; ampliar el dominio a toda la
mancha urbana, incluidos los trayectos que rodean la ensenada de La Paz; comparar RH con políticas
anticipatorias basadas en aprendizaje automático sobre el mismo entorno; y evaluar la escalabilidad
del simulador con instancias de mayor tamaño.

\section{Conclusiones}
En las instancias evaluadas del centro urbano de La Paz, coordinar los despachos no compensa cuando
la flota es holgada, pero sí a partir de una saturación moderada. Con una flota dimensionada para 18
pedidos por repartidor en la jornada, FCFS entrega 1.4~min antes que RH y el cumplimiento del SLA es
equivalente. Con 26 o más pedidos por repartidor, RH reduce el tiempo de entrega y mejora el
cumplimiento del servicio en todas las semillas, y cuando la flota es escasa entrega prácticamente
todos los pedidos que FCFS deja pendientes. Para un operador, esto sugiere que la decisión de
coordinar debe tomarse en función de la carga esperada por repartidor y no de forma fija.

La red vial importa para esta decisión. Los trayectos reales son en promedio 38\,\% más largos que
en línea recta, y un modelo euclidiano subestima tanto los tiempos de entrega como el beneficio de
coordinar. Finalmente, toda la evaluación puede reproducirse con herramientas abiertas en una
computadora portátil, lo que la hace accesible para operadores y desarrolladores de ciudades medias
que no cuentan con servicios comerciales de rutas.

\begin{thebibliography}{99}

\bibitem{Jazemi2023}
R. Jazemi, E. Alidadiani, K. Ahn y J. Jang, ``A review of literature on vehicle routing problems of
last-mile delivery in urban areas,'' \textit{Applied Sciences}, vol.~13, no.~24, art.~13015, 2023,
doi: 10.3390/app132413015.

\bibitem{Reyes2018}
D. Reyes, A. Erera, M. Savelsbergh, S. Sahasrabudhe y R. O'Neil, ``The meal delivery routing
problem,'' preprint, \textit{Optimization Online}, 2018. [En línea]. Disponible:
https://optimization-online.org/2018/04/6571/

\bibitem{Yildiz2019}
B. Yıldız y M. Savelsbergh, ``Provably high-quality solutions for the meal delivery routing
problem,'' \textit{Transportation Science}, vol.~53, no.~5, pp.~1372--1388, 2019,
doi: 10.1287/trsc.2018.0887.

\bibitem{Ulmer2021}
M. W. Ulmer, B. W. Thomas, A. M. Campbell y N. Woyak, ``The restaurant meal delivery problem: Dynamic
pickup and delivery with deadlines and random ready times,'' \textit{Transportation Science},
vol.~55, no.~1, pp.~75--100, 2021, doi: 10.1287/trsc.2020.1000.

\bibitem{Hildebrandt2022}
F. D. Hildebrandt y M. W. Ulmer, ``Supervised learning for arrival time estimations in restaurant
meal delivery,'' \textit{Transportation Science}, vol.~56, no.~4, pp.~1058--1084, 2022,
doi: 10.1287/trsc.2021.1095.

\bibitem{Boyaci2021}
B. Boyacı, T. H. Dang y A. N. Letchford, ``Vehicle routing on road networks: How good is Euclidean
approximation?,'' \textit{Computers \& Operations Research}, vol.~129, art.~105197, 2021,
doi: 10.1016/j.cor.2020.105197.

\bibitem{Pereira2021}
R. H. M. Pereira, M. Saraiva, D. Herszenhut, C. K. V. Braga y M. W. Conway, ``r5r: Rapid realistic
routing on multimodal transport networks with R5 in R,'' \textit{Findings}, 2021,
doi: 10.32866/001c.21262.

\bibitem{Kuhn1955}
H. W. Kuhn, ``The Hungarian method for the assignment problem,'' \textit{Naval Research Logistics
Quarterly}, vol.~2, no.~1--2, pp.~83--97, 1955, doi: 10.1002/nav.3800020109.

\bibitem{Boeing2022}
G. Boeing, ``Street network models and indicators for every urban area in the world,''
\textit{Geographical Analysis}, vol.~54, no.~3, pp.~519--535, 2022, doi: 10.1111/gean.12281.

\bibitem{Geisberger2008}
R. Geisberger, P. Sanders, D. Schultes y D. Delling, ``Contraction hierarchies: Faster and simpler
hierarchical routing in road networks,'' en \textit{Experimental Algorithms (WEA 2008)}, Lecture
Notes in Computer Science, vol.~5038, Berlín: Springer, 2008, pp.~319--333,
doi: 10.1007/978-3-540-68552-4\_24.

\bibitem{LuxenVetter2011}
D. Luxen y C. Vetter, ``Real-time routing with OpenStreetMap data,'' en \textit{Proc. 19th ACM
SIGSPATIAL Int. Conf. Advances in Geographic Information Systems}, Chicago, IL, EE.~UU., 2011,
pp.~513--516, doi: 10.1145/2093973.2094062.

\bibitem{Wu2024}
L. Wu \textit{et al.}, ``LaDe: The first comprehensive last-mile express dataset from industry,''
en \textit{Proc. 30th ACM SIGKDD Conf. Knowledge Discovery and Data Mining}, 2024,
pp.~5991--6002, doi: 10.1145/3637528.3671548.

\bibitem{Wilcoxon1945}
F. Wilcoxon, ``Individual comparisons by ranking methods,'' \textit{Biometrics Bulletin}, vol.~1,
no.~6, pp.~80--83, 1945, doi: 10.2307/3001968.

\bibitem{Friedman1937}
M. Friedman, ``The use of ranks to avoid the assumption of normality implicit in the analysis of
variance,'' \textit{Journal of the American Statistical Association}, vol.~32, no.~200,
pp.~675--701, 1937, doi: 10.1080/01621459.1937.10503522.

\bibitem{MunozVillamizar2021}
A. Muñoz-Villamizar, E. L. Solano-Charris, L. Reyes-Rubiano y J. Faulin, ``Measuring disruptions in
last-mile delivery operations,'' \textit{Logistics}, vol.~5, no.~1, art.~17, 2021,
doi: 10.3390/logistics5010017.

\end{thebibliography}

\end{document}
"
\documentclass[10pt]{article}

% ── Revisión doble ciega ────────────────────────────────────────────
\newif\ifblind
\ifdefined\conautores \blindfalse \else \blindtrue \fi

% ── Página (plantilla RIIM: carta; márgenes 2.0 / 2.0 / 3.5 / 2.25 cm;
%    encabezado y pie a 1.25 cm del borde) ──────────────────────────
\usepackage[letterpaper, left=2cm, right=2cm, top=3.5cm, bottom=2.25cm,
            headheight=1.45cm, headsep=0.8cm, footskip=0.62cm]{geometry}

% ── Tipografía ──────────────────────────────────────────────────────
\usepackage{fontspec}
\setmainfont{Montserrat}[
  Path=fonts/, Extension=.ttf,
  UprightFont=*-Regular, ItalicFont=*-Italic,
  BoldFont=*-SemiBold, BoldItalicFont=*-SemiBoldItalic]
\newfontfamily\semibold{Montserrat}[
  Path=fonts/, Extension=.ttf,
  UprightFont=*-SemiBold, ItalicFont=*-SemiBoldItalic, Ligatures=TeX]
\newfontfamily\arialfont{Arial}[Ligatures=TeX]
\usepackage[spanish, mexico, es-tabla, es-nodecimaldot, es-noshorthands]{babel}
\spanishplainpercent
\usepackage{amsmath, amssymb}
\usepackage{xcolor}
\definecolor{riimazul}{HTML}{003D7C}
\definecolor{cfcfs}{HTML}{1E88E5}
\definecolor{crh}{HTML}{F59E0B}

\usepackage{graphicx}
\usepackage{booktabs}
\usepackage{tabularx}
\usepackage{float}
\usepackage{flafter}
% Permite que tablas y figuras compartan página con el texto en lugar de ir a páginas solo de flotantes.
\renewcommand{\topfraction}{0.85}
\renewcommand{\bottomfraction}{0.7}
\renewcommand{\textfraction}{0.1}
\renewcommand{\floatpagefraction}{0.8}
\setcounter{topnumber}{3}
\setcounter{bottomnumber}{2}
\setcounter{totalnumber}{4}
\setlength{\textfloatsep}{12pt plus 2pt minus 2pt}
\setlength{\floatsep}{10pt plus 2pt minus 2pt}
\setlength{\intextsep}{10pt plus 2pt minus 2pt}
\usepackage{tikz}
\usetikzlibrary{arrows.meta, positioning, fit, backgrounds}
\usepackage{fancyhdr}
\usepackage{caption}
\usepackage[hidelinks]{hyperref}

% Cuerpo: Montserrat 10, interlineado sencillo, justificado, sin sangría, 8 pt después.
\renewcommand{\normalsize}{\fontsize{10}{12.2}\selectfont}
\normalsize
\setlength{\parindent}{0pt}
\setlength{\parskip}{8pt plus 1pt}
\setlength{\abovedisplayskip}{6pt}
\setlength{\belowdisplayskip}{6pt}
\setlength{\abovedisplayshortskip}{4pt}
\setlength{\belowdisplayshortskip}{4pt}

% Títulos de sección sin numerar: Montserrat SemiBold 12 pt azul; subtítulos SemiBold itálica 11 pt.
\makeatletter
\setcounter{secnumdepth}{0}
\renewcommand{\section}{\@startsection{section}{1}{0pt}{8pt}{0.1pt}
  {\semibold\fontsize{12}{14.6}\selectfont\color{riimazul}\raggedright}}
\renewcommand{\subsection}{\@startsection{subsection}{2}{0pt}{8pt}{0.1pt}
  {\semibold\itshape\fontsize{11}{13.4}\selectfont\addfontfeatures{LetterSpace=6.8}\raggedright}}
\makeatother

% Pies de figura y de tabla: 9 pt, centrados, "Figura N." / "Tabla N.".
\captionsetup{font={footnotesize}, labelfont={footnotesize}, labelsep=period,
  justification=centering, skip=4pt}
\captionsetup[table]{position=top}
\renewcommand{\footnotesize}{\fontsize{9}{11}\selectfont}

% Encabezado: logotipo RIIM (imagen de la plantilla recortada al 59.3 % izquierdo, como en Word) a la izquierda y "RIIM, Vol 3, Num 2, 2026" (Arial 8) a la derecha.
% Pie: número de página centrado; ISSN a la izquierda y URL a la derecha (9 pt).
\pagestyle{fancy}
\fancyhf{}
\renewcommand{\headrulewidth}{0pt}
\renewcommand{\footrulewidth}{0pt}
\fancyhead[C]{\makebox[\textwidth]{%
  \raisebox{-0.5\height}{\includegraphics[trim=0 0 335.53bp 0, clip, width=7.37cm, height=1.40cm]{img/riim_header.png}}%
  \hfill\raisebox{-0.5\height}{\arialfont\fontsize{8}{9.6}\selectfont RIIM, Vol 3, Num 2, 2026}}}
\fancyfoot[C]{\parbox[t]{\textwidth}{\centering\fontsize{9}{11}\selectfont\thepage\\[1pt]
  ISSN 3122-4806\hfill http://orizaba.tecnm.mx/revistaojs}}

% Listas compactas.
\makeatletter
\renewcommand{\@listI}{\leftmargin=1.2em \topsep=0pt \partopsep=0pt \itemsep=2pt \parsep=0pt}
\let\@listi\@listI
\makeatother

% Referencias IEEE (Arial 10, sangría francesa de 1.27 cm).
\makeatletter
\renewenvironment{thebibliography}[1]{%
  \section{Referencias}
  \list{\@biblabel{\@arabic\c@enumiv}}{%
    \arialfont\fontsize{10}{12}\selectfont
    \settowidth\labelwidth{\@biblabel{#1}}%
    \leftmargin=1.27cm \labelwidth=1.27cm \labelsep=0pt \itemindent=0pt
    \itemsep=0pt \parsep=0pt \topsep=0pt
    \def\makelabel##1{##1\hfil}%
    \usecounter{enumiv}\let\p@enumiv\@empty
    \renewcommand\theenumiv{\@arabic\c@enumiv}}%
  \sloppy\clubpenalty4000\widowpenalty4000\sfcode`\.\@m}
  {\def\@noitemerr{\@latex@warning{Empty `thebibliography' environment}}\endlist}
\makeatother

\hypersetup{pdftitle={Logística de última milla en redes viales reales: MDRP reproducible en La Paz, B.C.S.},
  pdfauthor={}, pdfsubject={RIIM}, pdfkeywords={enrutamiento dinámico, horizonte rodante, OpenStreetMap, simulación}}
\ifblind\else\hypersetup{pdfauthor={Alan José Marcos Sánchez Martínez}}\fi

% Notación.
\newcommand{\CtD}{\text{CtD}}
\newcommand{\Cdisp}{\mathcal{C}_{disp}(t)}

% Diagrama de caja: \caja{x}{mín}{Q1}{mediana}{Q3}{máx}{color}
\newcommand{\caja}[7]{%
  \draw (#1,#2) -- (#1,#3) (#1,#5) -- (#1,#6);
  \draw (#1-0.14,#2) -- (#1+0.14,#2) (#1-0.14,#6) -- (#1+0.14,#6);
  \filldraw[fill=#7!75, draw=black!70] (#1-0.3,#3) rectangle (#1+0.3,#5);
  \draw[line width=1pt] (#1-0.3,#4) -- (#1+0.3,#4);}

\begin{document}

% ── Título, autores y área ──────────────────────────────────────────
\begin{center}
  \setlength{\parskip}{0pt}
  \vspace*{6pt}
  {\semibold\fontsize{16}{19.5}\selectfont\addfontfeatures{LetterSpace=-3.1}%
   LOGÍSTICA DE ÚLTIMA MILLA EN REDES VIALES REALES: MDRP REPRODUCIBLE EN LA PAZ, B.C.S.\par}
  \vspace{18pt}
  \ifblind\else
  {\itshape\fontsize{9}{11}\selectfont
   Sánchez Martínez, Alan José Marcos\textsuperscript{1*}\par
   \textsuperscript{1}Tecnológico Nacional de México / Instituto Tecnológico de La Paz, m16310322@lapaz.tecnm.mx\par}
  \fi
  {\fontsize{9}{11}\selectfont Área de participación: Sistemas Computacionales\par}
\end{center}

\section{Resumen}
Las plataformas de entrega de comida deciden sin conocer los pedidos futuros, si asignan cada orden
al repartidor libre más cercano o si esperan unos minutos para agrupar órdenes y coordinar la
asignación. Este trabajo evalúa en qué condiciones conviene coordinar en una ciudad media mexicana,
La Paz, B.C.S., y cuanto cambia la evaluación al usar la red vial real. Se construyó un simulador de
código abierto que obtiene los tiempos de viaje de OSRM sobre OpenStreetMap y compara una política
reactiva (FCFS) con una heurística de horizonte rodante (RH) en cuatro niveles de saturación de la
flota, con diez instancias sintéticas por nivel. Con una flota holgada, FCFS entrega 1.4~min antes;
con 26 o más pedidos por repartidor en la jornada, RH reduce el tiempo de entrega en todas las
instancias, hasta en 67\,\%. Las distancias euclidianas subestiman los tiempos de entrega y el
beneficio de coordinar.

{\semibold\fontsize{12}{14.6}\selectfont\color{riimazul}Palabras clave}: \textit{enrutamiento dinámico, horizonte rodante, OpenStreetMap, simulación}

\section{Abstract}
\begin{otherlanguage}{english}\frenchspacing
Food delivery platforms must decide, without knowing future orders, whether to assign each order to
the nearest idle courier or to wait a few minutes, batch orders and coordinate the assignment. This
work evaluates under which conditions coordination pays off in a mid-sized Mexican city, La Paz,
B.C.S., and how much the evaluation changes when the real road network is used. An open-source
simulator was built that obtains travel times from OSRM over OpenStreetMap and compares a reactive
policy (FCFS) with a rolling-horizon heuristic (RH) at four fleet saturation levels, with ten
synthetic instances per level. With a slack fleet, FCFS delivers 1.4~min earlier; with 26 or more
orders per courier over the day, RH reduces delivery time in every instance, by up to 67\,\%.
Euclidean distances underestimate both delivery times and the benefit of coordination.

{\semibold\fontsize{12}{14.6}\selectfont\color{riimazul}Key words}: \textit{dynamic routing, rolling horizon, OpenStreetMap, simulation}
\end{otherlanguage}

\section{Introducción}
La última milla, el tramo que va del punto de recolección al domicilio del cliente, es uno de los
eslabones mas costosos y difíciles de la logística urbana. En zonas residenciales concentra grandes
volúmenes de entregas con horarios estrechos y condiciones que cambian durante el día
\cite{Jazemi2023}. La entrega de comida a domicilio agrava estas exigencias: los pedidos llegan sin
aviso, la comida pierde calidad con cada minuto de espera y el cliente espera recibirla en menos de
una hora.

En este contexto, una plataforma debe decidir de forma continua qué repartidor atiende cada orden y
en qué momento. Este problema, formalizado como \textit{Meal Delivery Routing Problem} (MDRP) por
Reyes \textit{et al.} \cite{Reyes2018}, admite dos estrategias básicas. La primera es reactiva: cada
orden se asigna en cuanto aparece al repartidor libre más cercano (\textit{First-Come, First-Served},
FCFS). La segunda es coordinada: el sistema espera unos minutos, agrupa órdenes del mismo restaurante
y decide de manera conjunta qué repartidor atiende cada agrupamiento, como en la heurística de
horizonte rodante (\textit{Rolling Horizon}, RH) propuesta en \cite{Reyes2018}. Coordinar tiene un
costo, porque algunas órdenes esperan a que se forme el agrupamiento o a que se entreguen otras antes
que ellas, y un beneficio, porque cada viaje atiende más pedidos. El balance entre ambos depende de
cuántos repartidores hay frente a la demanda.

La evidencia disponible sobre este balance proviene de entornos muy distintos a una ciudad media. RH
se evaluó con instancias que reproducen áreas metropolitanas de Estados Unidos y con tiempos de viaje
calculados como la distancia euclidiana multiplicada por un factor constante \cite{Reyes2018}. En
instancias pequeñas resueltas a optimalidad, su \textit{Click-to-Door} promedio queda apenas un
4\,\% por encima del óptimo \cite{Yildiz2019}, pero no se ha reportado cómo cambia su ventaja frente
a una política reactiva conforme varía el tamaño de la flota. Otras formulaciones del problema
incorporan la incertidumbre en los tiempos de preparación \cite{Ulmer2021} o la estimación del tiempo
de llegada, cuya desviación respecto a lo prometido pesa en la satisfacción del cliente tanto como la
rapidez misma \cite{Hildebrandt2022}.

La forma de medir los tiempos de viaje añade una segunda fuente de error. Las distancias por la red
vial suelen ser entre 1.2 y 1.4 veces mayores que en línea recta, y en problemas con ventanas de
tiempo estrechas, como la entrega de comida, una ruta factible con distancias euclidianas puede
resultar infactible sobre la red real \cite{Boyaci2021}. Una evaluación que ignore la red puede
equivocarse no solo en los tiempos de entrega, sino también en la carga efectiva de la flota y, por
lo tanto, en el beneficio de coordinar. Los servicios comerciales de rutas resuelven la precisión,
pero su costo de licencia limita su uso en investigación \cite{Pereira2021} y su funcionamiento
cerrado dificulta reproducir los experimentos.

La Paz, B.C.S., ilustra esta brecha: es una ciudad media sin datos públicos de operaciones de entrega
de comida, cuya red vial, condicionada por la costa y por pocas avenidas principales, obliga a
recorridos más largos que la línea recta. De ahí surge la pregunta que guía este trabajo:
\textit{en una ciudad media, ¿en qué condiciones de saturación de la flota conviene coordinar los
despachos en lugar de asignar cada pedido al llegar, y cuánto cambia esa evaluación al usar la red
vial real en lugar de distancias euclidianas?} La pregunta se divide en tres: cómo varía la ventaja
de RH con la saturación y a partir de qué nivel cambia la política más conveniente (PI1); cuánto más largos son los
trayectos por la red y cuánto cambia la comparación al ignorarla (PI2); y si la evaluación puede
hacerse de forma reproducible, con recursos modestos y solo con herramientas de código abierto (PI3).

Las contribuciones corresponden a esas tres preguntas: (i)~una comparación pareada de FCFS y RH en
cuatro niveles de saturación, con prueba estadística por instancia, que ubica el punto a partir del
cual coordinar compensa; (ii)~la medición de la circuidad de la red urbana de La Paz y dos escenarios
de control con distancias euclidianas que separan el efecto de la red del efecto de la política; y
(iii)~un entorno de simulación reproducible, construido sobre OpenStreetMap y OSRM, que ejecuta un
día de operación en alrededor de un minuto en una computadora portátil. El estudio no propone un
algoritmo nuevo ni compara RH con políticas del estado del arte: su objetivo es caracterizar cuándo
conviene coordinar en un entorno realista.

\section{Metodología}

\subsection{Formulación del problema y métricas}
Se adopta la formulación del MDRP de \cite{Reyes2018}. Un conjunto $\mathcal{R}$ de restaurantes con
ubicación fija recibe órdenes $o \in \mathcal{O}$, cada una con restaurante de origen $r_o$, instante
de colocación $a_o$, instante en que la comida está lista $e_o \geq a_o$ y punto de entrega $\ell_o$.
La orden solo se conoce a partir de $a_o$. Una flota $\mathcal{C}$ de repartidores, con turno
$[e_c, l_c]$, recoge y entrega las órdenes. Órdenes del mismo restaurante pueden formar un
agrupamiento $B$ que se recoge en una sola visita y se entrega en una ruta de varias paradas. Los
tiempos de servicio son $s_r$ en el restaurante, una vez por agrupamiento, y $s_o$ con cada cliente.

La métrica principal es el \textit{Click-to-Door} de cada orden,
$\CtD_o = t_{delivery,o} - a_o$. A partir de ella se reportan su media y su percentil 90, y el
cumplimiento del nivel de servicio (SLA, \textit{Service Level Agreement}), entendido como el
porcentaje de órdenes entregadas dentro del objetivo $\tau$:
\begin{equation}
  \text{SLA} = \frac{|\{o \in \mathcal{O}_{del} : \CtD_o \leq \tau\}|}{|\mathcal{O}_{del}|} \times 100,
  \label{eq:sla}
\end{equation}
donde $\mathcal{O}_{del}$ es el conjunto de órdenes entregadas. Como ambas métricas se calculan
sobre órdenes entregadas, se reporta también el porcentaje de órdenes sin entregar al cierre. Se
complementan con el tiempo que la comida lista espera a ser recogida (RtP, \textit{Ready-to-Pickup}),
$\text{RtP}_o = t_{pickup,o} - e_o$, la
distancia recorrida por pedido, la utilización de la flota (fracción del turno en conducción o en
servicio), el tamaño de agrupamiento y el costo por pedido, con un pago de $p_1$ por orden entregada
y un mínimo garantizado de $p_2$ por hora de turno \cite{Reyes2018}.

\subsection{Políticas de despacho}
El simulador avanza en epocas de $f = 5$~min para ambas políticas: en cada época, FCFS asigna una por
una las órdenes nuevas y RH resuelve su optimización. Las dos calculan los instantes de recolección y
entrega con las mismas reglas. Un agrupamiento $B$ solo puede recogerse cuando ya está lista la
última de sus órdenes, es decir, a partir del instante $e_B = \max_{o\in B} e_o$. El repartidor $c$
llega al restaurante en el instante $t_{arr,r}$ y, como en \cite{Reyes2018}, cada tiempo de servicio
se reparte mitad antes y mitad después del momento de recoger o de entregar. Si la comida ya está
lista, la recolección ocurre al cumplirse la primera mitad del servicio en el restaurante; si no, el
repartidor espera hasta que esté lista. Por eso el instante de recolección es
$t_{pickup}(c,B) = \max\bigl(e_B,\; t_{arr,r} + s_r/2\bigr)$. Después, el repartidor sale del
restaurante al terminar la segunda mitad del servicio, recorre el tramo hasta el primer cliente y le
entrega al cumplirse la primera mitad del servicio con él; cada entrega siguiente se calcula igual a
partir de la anterior.

\textbf{FCFS (línea base reactiva).} En cada época, las órdenes reveladas sin repartidor se procesan
por orden de colocación. Cada una se asigna al repartidor libre más cercano en línea recta al
restaurante (esta distancia solo sirve para escoger al repartidor); la ruta repartidor $\to$ restaurante
$\to$ cliente y sus tiempos se obtienen de OSRM, y el repartidor queda comprometido de inmediato. Cada
viaje lleva una sola orden.

\textbf{RH (política coordinada).} Sigue el esquema de \cite{Reyes2018} y en cada época $t$ ejecuta
cinco pasos. Primero, delimita las órdenes que entran a la optimización: de las órdenes ya reveladas y
aún no entregadas, $\mathcal{O}_t$, solo considera las que estarán listas dentro del horizonte de
asignación $\Delta_u$, es decir, $U_t = \{o \in \mathcal{O}_t : e_o \leq t + \Delta_u\}$; planear
órdenes cuya comida tardará mucho en estar lista sería poco útil en un entorno que cambia cada pocos
minutos.

Segundo, estima cuántas órdenes conviene juntar en cada viaje. El tamaño objetivo de agrupamiento
$Z_t$ compara la demanda próxima con la oferta próxima: el numerador cuenta las órdenes que estarán
listas en los siguientes $\Delta_1$ minutos, y el denominador, los repartidores en turno, libres u
ocupados, que quedarán disponibles en los siguientes $\Delta_2$ minutos, donde $e^{avail}_c$ es el
instante en que el repartidor $c$ termina su ruta actual:
\begin{equation}
  Z_t = \left\lceil \frac{|\{o \in \mathcal{O}_t : e_o \leq t + \Delta_1\}|}
                        {|\{c \in \mathcal{C}_{turno}(t) : e^{avail}_c \leq t + \Delta_2\}|} \right\rceil.
  \label{eq:zt}
\end{equation}
Si hay tantos repartidores como ordenes, $Z_t = 1$ y cada orden tiende a viajar sola; si la flota se
satura, $Z_t$ crece y se forman agrupamientos más grandes. $Z_t$ se acota en 5, y ningún agrupamiento
puede superar $B_{max} = 4$ órdenes.

Tercero, forma los agrupamientos de cada restaurante $r$. Sus órdenes $U_{t,r}$ se reparten en
\[
  m_r = \min\bigl(\lceil |U_{t,r}|/Z_t \rceil,\; |\Cdisp|\bigr)
\]
agrupamientos, de modo que cada uno reciba alrededor de $Z_t$ órdenes sin crear más agrupamientos que
repartidores disponibles. Las órdenes se toman por \textit{ready time}, y cada una se inserta en el
agrupamiento y la posición de la ruta que menos incrementan el costo
\begin{equation}
  \text{costo}(B, o, p) = T_{ruta}(B \oplus_p o) + \beta\, T_{servicio}(B \oplus_p o),
  \label{eq:costo}
\end{equation}
donde $B \oplus_p o$ es el agrupamiento con la orden $o$ insertada en la posición $p$, $T_{ruta}$ es
la duración de la ruta según OSRM y $T_{servicio}$ el tiempo de servicio acumulado. El peso $\beta$, que en
\cite{Reyes2018} representa la pérdida de frescura y se calibra en $[0, 1]$, se fija en 0.1 para que la
duración de la ruta domine el costo. Después, dos iteraciones de extracción y reinserción mueven
cada orden al agrupamiento donde su costo es menor.

Cuarto, asigna agrupamientos a repartidores. Los agrupamientos se clasifican en tres grupos de
prioridad: urgentes, cuando ningún repartidor puede entregar su orden más antigua dentro del objetivo
$\tau$; retrasados, cuando el objetivo aún es alcanzable pero ningún repartidor llega al restaurante
antes de que la comida esté lista; y normales. Los grupos se resuelven en ese orden, para que los
urgentes obtengan primero a los repartidores más cercanos. En cada grupo se resuelve un problema de
asignación lineal, en el que cada repartidor recibe a lo más un agrupamiento, con el peso
\begin{equation}
  w_{c,B} = \frac{|B|}{T_{delivery}(c,B)} - \theta\,\bigl(t_{pickup}(c,B) - e_B\bigr) - P_g.
  \label{eq:peso}
\end{equation}
El primer término mide la eficiencia del viaje en órdenes entregadas por hora, donde $T_{delivery}(c,B)$
es el tiempo desde que $c$ queda libre hasta la última entrega de $B$. El segundo penaliza, con
$\theta = 0.1$, los minutos que la comida espera lista en el restaurante, y $P_g$ resta 100 si el par
excede el máximo absoluto $\tau_{max}$ y 50 si la recolección se retrasa. La asignación de máximo peso
se obtiene de forma exacta con un algoritmo de tipo húngaro \cite{Kuhn1955}.

Quinto, el compromiso en dos etapas de \cite{Reyes2018} decide qué parte de la asignación se
ejecuta. Si el repartidor llega al restaurante y todas las órdenes estarán listas dentro de la época,
o si alguna orden lleva más de $\Delta_x$ minutos lista, se fija la ruta completa. En otro caso, el
repartidor solo se dirige al restaurante, como un preposicionamiento, y en la época siguiente el
agrupamiento puede crecer o cambiar de repartidor con las órdenes que se hayan revelado.

\subsection{Entorno de simulación reproducible}
La Figura~\ref{fig:arquitectura} resume el entorno. El extracto de OpenStreetMap de Baja California
Sur, modelo de datos viales abierto y de cobertura mundial \cite{Boeing2022}, se preprocesa con el
algoritmo \textit{Contraction Hierarchies} \cite{Geisberger2008}, que permite resolver cada consulta
de camino mínimo en milisegundos. El servidor OSRM \cite{LuxenVetter2011}, versión 5.27.1, atiende
las consultas de ruta del simulador por HTTP; su imagen oficial de Docker permite reproducir el
entorno en cualquier sistema operativo. El simulador, escrito en Python, avanza por épocas, guarda en
memoria las rutas ya consultadas y produce, para cada corrida, los registros por orden y por
repartidor de los que se derivan los indicadores.

\begin{figure}[htbp]
\centering
\begin{tikzpicture}[
  font=\footnotesize,
  caja/.style={draw=black!70, rounded corners=2pt, align=center, minimum height=1.05cm,
               text width=#1, fill=white},
  caja/.default=2.8cm,
  flecha/.style={-{Stealth[length=5pt]}, thick, black!70},
  capa/.style={draw=#1!70!black, rounded corners=4pt, fill=#1!8, inner xsep=7pt, inner ysep=7pt}]
  \node[caja=2.8cm] (osm)   at (0,0)      {Extracto\\OpenStreetMap};
  \node[caja=3.0cm] (ch)    at (3.55,0)    {Preprocesamiento\\\textit{Contraction}\\\textit{Hierarchies}};
  \node[caja=2.8cm] (osrm)  at (7.1,0)    {Servidor OSRM\\(API de rutas)};
  \node[caja=2.8cm] (gen)   at (0,-2.3)   {Generador de\\instancias (PPNH)};
  \node[caja=2.8cm] (sim)   at (3.55,-2.3) {Simulador por\\épocas ($f = 5$ min)};
  \node[caja=2.8cm] (pol)   at (7.1,-2.3) {Política\\FCFS o RH};
  \node[caja=2.8cm] (kpi)   at (10.9,-2.3) {Registros por orden\\y 21 indicadores};
  \node[caja=2.8cm] (est)   at (10.9,0)   {Análisis estadístico\\(Wilcoxon, Friedman)};
  \draw[flecha] (osm) -- (ch);
  \draw[flecha] (ch) -- (osrm);
  \draw[flecha] (gen) -- (sim);
  \draw[flecha, {Stealth[length=5pt]}-{Stealth[length=5pt]}] (sim) -- (pol);
  \draw[flecha, dashed, {Stealth[length=5pt]}-{Stealth[length=5pt]}] (osrm) --
     node[right, font=\scriptsize, align=left] {tiempos de\\viaje (HTTP)} (pol);
  \draw[flecha] (pol) -- (kpi);
  \draw[flecha] (kpi) -- (est);
  \begin{scope}[on background layer]
    \node[capa=blue, fit=(osm)(ch)(osrm)] {};
    \node[capa=green, fit=(gen)(sim)(pol)] {};
    \node[capa=orange, fit=(kpi)(est)] {};
  \end{scope}
\end{tikzpicture}
\caption{Arquitectura del entorno: enrutamiento sobre la red vial (azul), simulación (verde) y
análisis (naranja).}
\label{fig:arquitectura}
\end{figure}

\subsection{Generación de instancias}
Al no existir datos públicos de entregas de comida en La Paz, las instancias se generan de forma
sintética. Las llegadas siguen un proceso de Poisson no homogéneo con intensidad bimodal
\begin{equation}
  \lambda(t) = \lambda_0
  + A_m \exp\!\left(-\frac{(t-\mu_m)^2}{2\sigma^2}\right)
  + A_e \exp\!\left(-\frac{(t-\mu_e)^2}{2\sigma^2}\right),
  \label{eq:lambda}
\end{equation}
siguiendo el patrón bimodal observado en entregas de última milla \cite{Wu2024}. A falta de datos
locales, los parámetros son supuestos: ventana operativa de 08:00 a 19:00, picos en $\mu_m = $ 09:00
y $\mu_e = $ 17:00, $\lambda_0 = 0.6$ órdenes/min, $A_m = 2.4$, $A_e = 2.0$ y $\sigma = 45$~min. El perfil se escala para
obtener alrededor de mil órdenes por día (entre 988 y 1\,047 según la semilla). El tiempo de
preparación es $\max(5, \mathcal{N}(12, 3^2))$~min. Los 25 restaurantes y los clientes se ubican
por muestreo uniforme dentro de un polígono de 8.4~km$^2$ del centro urbano, y cada cliente se asocia
a un restaurante situado a entre 2 y 5~km en línea recta. Todos los repartidores comparten un turno
que inicia 15~min antes del primer pedido y termina una hora después del último \textit{ready time},
y parten del centroide de los restaurantes. La flota se dimensiona a partir de la razón de
saturación $\rho$, el número objetivo de pedidos por repartidor en la jornada:
\begin{equation}
  |\mathcal{C}| = \max\!\left(n_{\min},\; \left\lceil |\mathcal{O}| / \rho \right\rceil\right),
  \qquad n_{\min} = 10.
  \label{eq:flota}
\end{equation}
Fijar $\rho$, y no el número de repartidores, mantiene la misma carga relativa en todas las semillas
de un escenario aunque el número de órdenes varíe.

\subsection{Diseño experimental}
La Tabla~\ref{tab:parametros} reúne los parámetros, comunes a ambas políticas, y los cuatro
escenarios de saturación. Cada semilla (2025 a 2034) genera una instancia, y sobre ella se ejecutan
FCFS y RH con la misma flota; como el simulador no tiene componentes aleatorios propios, lo único que
cambia dentro de cada par es la política. Esto da 80 corridas sobre la red vial. Para evitar la
pseudorreplicación, cada métrica se promedia por semilla y la comparación se hace con la prueba de
rangos con signo de Wilcoxon pareada \cite{Wilcoxon1945} sobre las diez diferencias de cada
escenario, con intervalo de confianza $t$ al 95\,\% y la correlación biserial de rangos $r_{rb}$ como
tamaño de efecto. La prueba de Friedman \cite{Friedman1937} sobre la matriz de diferencias
$10 \times 4$ evalúa si el efecto cambia entre escenarios.

\begin{table}[htbp]
\centering
\footnotesize
\caption{Parámetros del simulador y escenarios de saturación.}
\label{tab:parametros}
\setlength{\tabcolsep}{4pt}
\begin{tabular}{@{}l c c @{\hspace{0.4cm}} l c c@{}}
\toprule
Parámetro & Símbolo & Valor & Parámetro & Símbolo & Valor \\
\midrule
Frecuencia de optimización & $f$ & 5 min & Penalización de frescura & $\beta$ & 0.1 \\
Horizonte de asignación & $\Delta_u$ & 20 min & Penalización por espera & $\theta$ & 0.1 \\
Anticipación (demanda, oferta) & $\Delta_1$, $\Delta_2$ & 20 min & Espera para forzar compromiso & $\Delta_x$ & 10 min \\
CtD objetivo / máximo & $\tau$, $\tau_{max}$ & 40/90 min & Tamaño máximo de agrupamiento & $B_{max}$ & 4 \\
Tiempos de servicio & $s_r$, $s_o$ & 4 min & Pago por orden / mínimo por hora & $p_1$, $p_2$ & \$10 / \$15 \\
\midrule
\multicolumn{6}{@{}l}{Escenarios: S1, $\rho = 18$ (55--59 repartidores); S2, $\rho = 26$ (38--41); S3, $\rho = 35$ (29--30); S4, $\rho = 52$ (19--21).} \\
\bottomrule
\end{tabular}
\end{table}

Para medir el efecto de la red vial se calculó la circuidad $\kappa = d_{OSRM}/d_{recta}$ sobre todos
los trayectos restaurante--cliente, y el experimento completo se repitió con dos controles en los que
el tiempo de viaje es la distancia en línea recta (Haversine) a velocidad constante: la velocidad
media que OSRM asigna a esos trayectos (651~m/min), que deja como única diferencia la geometría de
las rutas, y esa velocidad dividida entre $\kappa$ (472~m/min), que iguala la duración media de los
viajes. Cada control suma otras 80 corridas.

\section{Resultados y discusión}

\subsection{Calidad de servicio y cambio de la política más conveniente}
La Tabla~\ref{tab:calidad} muestra los indicadores de calidad de servicio. Con la flota holgada de S1,
FCFS entrega en promedio 1.4~min antes que RH: con repartidores de sobra, llevar cada pedido directo
del restaurante al cliente es más rápido que esperar a formar un agrupamiento, y en un agrupamiento
cada entrega espera a las anteriores. Aun así, el cumplimiento del SLA es prácticamente igual en ambas
políticas. A partir de S2 la relación se invierte: RH reduce el \textit{Click-to-Door} medio en
36.8\,\%, el percentil 90 en 46.0\,\% y eleva el cumplimiento del SLA en 31.5 puntos porcentuales. La
ventaja crece en S3 y S4, donde FCFS deja sin entregar al cierre el 11.3\,\% y el 39.9\,\% de los
pedidos, mientras que RH entrega prácticamente todos. Como el CtD y el SLA se calculan sobre órdenes
entregadas, la tabla favorece a FCFS en esos escenarios: medido sobre todas las órdenes generadas, su
cumplimiento cae a 6.4\,\% en S3 y 3.8\,\% en S4, frente a 57.6\,\% y 42.4\,\% bajo RH.

\begin{table}[htbp]
\centering
\footnotesize
\caption{Calidad de servicio, promedio de 10 semillas. $\Delta$ es el cambio de RH respecto a FCFS.}
\label{tab:calidad}
\setlength{\tabcolsep}{3pt}
\begin{tabular*}{\textwidth}{@{\extracolsep{\fill}}l rrr rrr rrr rr@{}}
\toprule
& \multicolumn{3}{c}{CtD medio (minutos)} & \multicolumn{3}{c}{CtD percentil 90 (minutos)} & \multicolumn{3}{c}{SLA (\%)} & \multicolumn{2}{c}{Sin entregar (\%)} \\
\cmidrule(lr){2-4} \cmidrule(lr){5-7} \cmidrule(lr){8-10} \cmidrule(l){11-12}
Esc. & FCFS & RH & $\Delta$ (\%) & FCFS & RH & $\Delta$ (\%) & FCFS & RH & $\Delta$ (pp) & FCFS & RH \\
\midrule
S1 & 28.93  & 30.31 & $+4.8$  & 43.33  & 43.80  & $+1.1$  & 84.6 & 83.4 & $-1.1$  & 0.0  & 0.0 \\
S2 & 56.07  & 35.42 & $-36.8$ & 97.91  & 52.87  & $-46.0$ & 38.1 & 69.6 & $+31.5$ & 0.1  & 0.0 \\
S3 & 108.49 & 45.36 & $-58.2$ & 169.54 & 79.10  & $-53.3$ & 7.3  & 57.6 & $+50.3$ & 11.3 & 0.0 \\
S4 & 224.88 & 73.63 & $-67.3$ & 321.10 & 186.18 & $-42.0$ & 6.2  & 42.5 & $+36.3$ & 39.9 & 0.2 \\
\bottomrule
\end{tabular*}
\end{table}

La Figura~\ref{fig:boxplot} muestra la distribución completa del \textit{Click-to-Door}. En S1 ambas
cajas son casi idénticas y quedan por debajo del objetivo de 40~min; desde S2, la distribución de
FCFS se alarga hacia valores de varias horas, mientras que la de RH se mantiene compacta y con la
mediana cerca del objetivo.

\begin{figure}[htbp]
\centering
\begin{tikzpicture}[y=0.0100cm, x=1cm, font=\footnotesize]
  \foreach \p/\esc in {0/S1 ($\rho=18$), 1/S2 ($\rho=26$), 2/S3 ($\rho=35$), 3/S4 ($\rho=52$)} {
    \begin{scope}[xshift=\p*4.15cm]
      \draw[black!60] (0,0) rectangle (3.6,390);
      \foreach \y in {0,100,200,300} \draw[black!12] (0,\y) -- (3.6,\y);
      \draw[red!70, dashed] (0,40) -- (3.6,40);
      \node[above] at (1.8,390) {\esc};
      \node[below] at (1.1,0) {FCFS};
      \node[below] at (2.5,0) {RH};
    \end{scope}}
  \foreach \y in {0,100,200,300} \node[left] at (0,\y) {\y};
  \node[rotate=90] at (-1.25,195) {\textit{Click-to-Door} (minutos)};
  % S1
  \caja{1.1}{10.9}{21.7}{26.4}{35.1}{55.1}{cfcfs}
  \caja{2.5}{11.4}{22.8}{28.5}{36.1}{56.1}{crh}
  % S2
  \begin{scope}[xshift=4.15cm]
  \caja{1.1}{11.9}{27.3}{54.3}{80.6}{121.7}{cfcfs}
  \caja{2.5}{11.4}{24.5}{32.5}{42.2}{68.7}{crh}
  \end{scope}
  % S3
  \begin{scope}[xshift=8.3cm]
  \caja{1.1}{11.9}{72.1}{109.1}{148.2}{199.9}{cfcfs}
  \caja{2.5}{11.4}{26.6}{36.5}{50.3}{85.8}{crh}
  \end{scope}
  % S4
  \begin{scope}[xshift=12.45cm]
  \caja{1.1}{14.9}{141.0}{266.0}{306.5}{367.5}{cfcfs}
  \caja{2.5}{12.3}{31.1}{43.8}{79.3}{151.5}{crh}
  \end{scope}
\end{tikzpicture}
\caption{Distribución del \textit{Click-to-Door} de las órdenes entregadas en las diez semillas
(cajas: cuartiles; bigotes: 1.5 veces el rango intercuartílico). La línea discontinua marca el
objetivo $\tau = 40$~min.}
\label{fig:boxplot}
\end{figure}

La Tabla~\ref{tab:wilcoxon} confirma, a nivel de instancia, el cambio de la política más conveniente. En S2, S3 y S4, RH
supera a FCFS en las diez semillas, tanto en CtD como en SLA, con $p = 0.0020$, el valor mínimo que
admite la prueba con $n = 10$, y $|r_{rb}| = 1$. En S1 la ventaja de FCFS en CtD es significativa
pero pequeña, y la diferencia en SLA no lo es. La prueba de Friedman confirma que la magnitud del
efecto cambia con la saturación ($\chi^2_F = 30.0$, $p < 0.001$). La Figura~\ref{fig:saturacion}
resume la respuesta a PI1: las curvas de CtD se cruzan entre $\rho = 18$ y $\rho = 26$, y por encima
de ese punto la ventaja de coordinar crece con la carga.

\begin{table}[htbp]
\centering
\footnotesize
\caption{Wilcoxon pareada por semilla ($n = 10$), $\bar{\Delta} = \bar{x}_{FCFS} - \bar{x}_{RH}$: favorecen a RH
$\bar{\Delta} > 0$ en CtD y $\bar{\Delta} < 0$ en SLA.}
\label{tab:wilcoxon}
\begin{tabular}{@{}ll rr l c r@{}}
\toprule
Esc. & Métrica & $\bar{x}_{FCFS}$ & $\bar{x}_{RH}$ & $\bar{\Delta}$ [IC 95\,\%] & $p$ & $r_{rb}$ \\
\midrule
S1 & CtD (minutos) & 28.93  & 30.31 & $-1.38$ [$-2.35$, $-0.41$]        & 0.0137 & $-0.85$ \\
S1 & SLA (\%)  & 84.6   & 83.4  & $+1.1$ [$-1.8$, $+4.1$]           & 0.2324 & $+0.45$ \\
S2 & CtD (minutos) & 56.07  & 35.42 & $+20.64$ [$+18.98$, $+22.31$]     & 0.0020 & $+1.00$ \\
S2 & SLA (\%)  & 38.1   & 69.6  & $-31.5$ [$-33.7$, $-29.2$]        & 0.0020 & $-1.00$ \\
S3 & CtD (minutos) & 108.49 & 45.36 & $+63.13$ [$+56.14$, $+70.13$]     & 0.0020 & $+1.00$ \\
S3 & SLA (\%)  & 7.3    & 57.6  & $-50.3$ [$-52.6$, $-48.0$]        & 0.0020 & $-1.00$ \\
S4 & CtD (minutos) & 224.88 & 73.63 & $+151.24$ [$+143.71$, $+158.78$]  & 0.0020 & $+1.00$ \\
S4 & SLA (\%)  & 6.2    & 42.5  & $-36.3$ [$-37.2$, $-35.3$]        & 0.0020 & $-1.00$ \\
\bottomrule
\end{tabular}
\end{table}

\begin{figure}[htbp]
\centering
\begin{tikzpicture}[font=\footnotesize]
  % Panel izquierdo: CtD medio
  \begin{scope}[x=0.16cm, y=0.018cm]
    \draw[black!60] (14,0) rectangle (56,240);
    \foreach \y in {0,50,100,150,200} {\draw[black!12] (14,\y) -- (56,\y); \node[left] at (14,\y) {\y};}
    \foreach \x in {18,26,35,52} {\draw[black!12] (\x,0) -- (\x,240); \node[below] at (\x,0) {\x};}
    \draw[red!70, dashed] (14,40) -- (56,40);
    \draw[cfcfs, very thick] (18,28.93) -- (26,56.07) -- (35,108.49) -- (52,224.88);
    \foreach \x/\y in {18/28.93, 26/56.07, 35/108.49, 52/224.88} \fill[cfcfs] (\x,\y) circle (2pt);
    \draw[crh, very thick] (18,30.31) -- (26,35.42) -- (35,45.36) -- (52,73.63);
    \foreach \x/\y in {18/30.31, 26/35.42, 35/45.36, 52/73.63} \fill[crh] (\x-0.45,\y-4.5) rectangle (\x+0.45,\y+4.5);
    \node[rotate=90] at (8.5,120) {CtD medio (minutos)};
    \node at (35,-30) {Pedidos por repartidor en la jornada ($\rho$)};
  \end{scope}
  % Panel derecho: SLA
  \begin{scope}[xshift=8.6cm, x=0.16cm, y=0.043cm]
    \draw[black!60] (14,0) rectangle (56,100);
    \foreach \y in {0,20,40,60,80,100} {\draw[black!12] (14,\y) -- (56,\y); \node[left] at (14,\y) {\y};}
    \foreach \x in {18,26,35,52} {\draw[black!12] (\x,0) -- (\x,100); \node[below] at (\x,0) {\x};}
    \draw[cfcfs, very thick] (18,84.6) -- (26,38.1) -- (35,7.3) -- (52,6.2);
    \foreach \x/\y in {18/84.6, 26/38.1, 35/7.3, 52/6.2} \fill[cfcfs] (\x,\y) circle (2pt);
    \draw[crh, very thick] (18,83.4) -- (26,69.6) -- (35,57.6) -- (52,42.5);
    \foreach \x/\y in {18/83.4, 26/69.6, 35/57.6, 52/42.5} \fill[crh] (\x-0.45,\y-1.9) rectangle (\x+0.45,\y+1.9);
    \node[rotate=90] at (8.5,50) {Cumplimiento del SLA (\%)};
    \node at (35,-12.5) {Pedidos por repartidor en la jornada ($\rho$)};
    % Leyenda
    \draw[cfcfs, very thick] (34,93) -- (38,93); \fill[cfcfs] (36,93) circle (2pt); \node[right] at (38,93) {FCFS};
    \draw[crh, very thick] (46,93) -- (50,93); \fill[crh] (47.55,91.1) rectangle (48.45,94.9); \node[right] at (50,93) {RH};
  \end{scope}
\end{tikzpicture}
\caption{\textit{Click-to-Door} medio y cumplimiento del SLA según la saturación de la flota. La línea
discontinua marca el objetivo $\tau = 40$~min.}
\label{fig:saturacion}
\end{figure}

\subsection{Eficiencia operativa y agrupamiento}
La Tabla~\ref{tab:eficiencia} explica el origen de la ventaja de RH. Aun en S1, RH reduce la espera
de la comida lista en el restaurante (RtP), porque en los picos de demanda la flota se ocupa por completo y
FCFS, con un pedido por viaje, acumula órdenes sin recoger, mientras que RH recoge varias en una sola
visita. Al agrupar, RH recorre entre 13.5\,\% y 49.4\,\% menos distancia por pedido y mantiene una
utilización de la flota menor que FCFS en S1--S3: con la misma flota, entrega lo mismo con menos
tiempo de conducción y, en S4, convierte ese margen en pedidos entregados. El tamaño de agrupamiento
crece con la saturación; en S1 una cuarta parte de los pedidos (25.4\,\%) viaja sola y en S4 el
80.9\,\% viaja en agrupamientos de tres o cuatro. El costo por pedido solo difiere en S1, donde muchos
repartidores reciben el pago mínimo y RH concentra las entregas en menos de ellos; desde S2 ambos
convergen al pago por orden.

\begin{table}[htbp]
\centering
\footnotesize
\caption{Eficiencia operativa y tamaño del agrupamiento en que viaja cada pedido (promedio de 10 semillas).}
\label{tab:eficiencia}
\setlength{\tabcolsep}{3pt}
\begin{tabular*}{\textwidth}{@{\extracolsep{\fill}}l rr rr rr r rr@{}}
\toprule
& \multicolumn{2}{c}{RtP medio (minutos)} & \multicolumn{2}{c}{Distancia/pedido (km)} & \multicolumn{2}{c}{Utilización (\%)} & Agrupamiento & \multicolumn{2}{c}{Costo/pedido (\$)} \\
\cmidrule(lr){2-3} \cmidrule(lr){4-5} \cmidrule(lr){6-7} \cmidrule(lr){8-8} \cmidrule(l){9-10}
Esc. & FCFS & RH & FCFS & RH & FCFS & RH & RH & FCFS & RH \\
\midrule
S1 & 7.78   & 4.94  & 4.71 & 4.08 & 37.7 & 31.2 & 2.54 & 10.80 & 11.59 \\
S2 & 34.92  & 9.33  & 4.88 & 3.60 & 55.2 & 41.4 & 2.80 & 10.00 & 10.03 \\
S3 & 87.35  & 18.62 & 5.07 & 3.17 & 66.8 & 51.4 & 3.03 & 10.00 & 10.00 \\
S4 & 203.75 & 46.20 & 5.20 & 2.63 & 67.3 & 67.8 & 3.30 & 10.00 & 10.00 \\
\bottomrule
\end{tabular*}
\end{table}

\subsection{Efecto de la red vial}
Sobre los 10\,184 trayectos restaurante--cliente de las diez instancias, la distancia por la red es
en promedio 3.30~km y en línea recta 2.39~km: los trayectos reales son 38\,\% más largos (circuidad
$\kappa = 1.38$), dentro del rango de 1.2 a 1.4 que reporta la literatura \cite{Boyaci2021}.

Para medir cuánto cambia la comparación, se repitió el experimento con tiempos de viaje en línea
recta. A la velocidad media de la red (651~m/min, primer control), ambas políticas parecen más
rápidas: en S2, el CtD de FCFS baja de 56.1 a 41.8~min y el de RH, de 35.4 a 30.3~min. Como FCFS
mejora más, la ventaja de RH se reduce en los cuatro escenarios ($p = 0.002$); en S3, por ejemplo,
pasa de 63.1 a 37.1~min. El segundo control baja la velocidad hasta que los viajes en línea recta
duran en promedio lo mismo que por la red ($651/1.38 \approx 472$~m/min); con él desaparece entre el
62 y el 84\,\% de esa reducción. La mayor parte del efecto se debe, por tanto, a que los trayectos
reales son más largos y cargan más a la flota, y el resto, a la forma de la red ($p \leq 0.004$).
Ignorar la red hace parecer más rápidas a ambas políticas y subestima el beneficio de coordinar, lo
que responde a la pregunta PI2.

\subsection{Reproducibilidad y costo de cómputo}
Toda la cadena (OpenStreetMap, OSRM y Python) es de código abierto, las instancias se regeneran a
partir de semillas fijas y cada corrida pasa verificaciones automáticas de integridad: ninguna orden
se recoge antes de estar lista ni se entrega antes de recogerse, no hay entregas duplicadas y ninguna
ruta recurrió a la estimación de respaldo en línea recta (de haber fallado OSRM). Con 1\,038 órdenes, simular un día toma
entre 1.8 y 2.7~segundos con FCFS y entre 57 y 80~segundos con RH, incluidas las consultas a OSRM, con menos de
225~MB de memoria por proceso, en una computadora portátil con procesador Apple M1 Pro y 16~GB de
RAM. Esto responde a PI3: la evaluación es reproducible y no requiere infraestructura especializada.

\subsection{Limitaciones}
El reloj del simulador avanza en épocas de 5~min. En cada época primero se asignan los pedidos y
después se registran las rutas terminadas, de modo que un repartidor que termina a mitad de una época
queda libre hasta la siguiente y espera entre 5 y 10~min su nuevo viaje. Como esta espera ocurre una
vez por viaje, perjudica más a FCFS, que hace un viaje por pedido. Para comprobar que no determina el
resultado, se repitieron las 80 corridas sobre la red vial con el orden inverso (primero se registran
las rutas terminadas y luego se asigna): FCFS sigue ganando en S1 y RH en S2--S4 en todas las
semillas, y el cruce permanece entre S1 y S2. La reducción del CtD baja a 14.6\,\% en S2, 41.4\,\% en
S3 y 63.0\,\% en S4, por lo que las magnitudes dependen de este detalle de implementación, pero las
conclusiones se mantienen. Además, la demanda es
sintética: sigue el patrón bimodal de la recolección de paquetes en LaDe con parámetros supuestos, no
calibrados con la demanda local de comida, y el tipo de ciudad influye en el
comportamiento de la última milla \cite{MunozVillamizar2021}. Los tiempos de preparación y de
servicio son homogéneos, la red se considera sin congestión y las ubicaciones se muestrean de forma
uniforme dentro del dominio urbano.

\section{Trabajo a futuro}
Quedan pendientes seis líneas de trabajo: calibrar la demanda, los tiempos de preparación y la
ubicación de restaurantes y clientes con datos locales de una plataforma de entregas o de encuestas;
explorar valores de $\rho$ entre 18 y 26 para ubicar con precisión el punto de cruce y evaluar
políticas híbridas que activen la coordinación solo cuando la carga lo justifique; incorporar
velocidades dependientes de la hora para representar la congestión; ampliar el dominio a toda la
mancha urbana, incluidos los trayectos que rodean la ensenada de La Paz; comparar RH con políticas
anticipatorias basadas en aprendizaje automático sobre el mismo entorno; y evaluar la escalabilidad
del simulador con instancias de mayor tamaño.

\section{Conclusiones}
En las instancias evaluadas del centro urbano de La Paz, coordinar los despachos no compensa cuando
la flota es holgada, pero sí a partir de una saturación moderada. Con una flota dimensionada para 18
pedidos por repartidor en la jornada, FCFS entrega 1.4~min antes que RH y el cumplimiento del SLA es
equivalente. Con 26 o más pedidos por repartidor, RH reduce el tiempo de entrega y mejora el
cumplimiento del servicio en todas las semillas, y cuando la flota es escasa entrega prácticamente
todos los pedidos que FCFS deja pendientes. Para un operador, esto sugiere que la decisión de
coordinar debe tomarse en función de la carga esperada por repartidor y no de forma fija.

La red vial importa para esta decisión. Los trayectos reales son en promedio 38\,\% más largos que
en línea recta, y un modelo euclidiano subestima tanto los tiempos de entrega como el beneficio de
coordinar. Finalmente, toda la evaluación puede reproducirse con herramientas abiertas en una
computadora portátil, lo que la hace accesible para operadores y desarrolladores de ciudades medias
que no cuentan con servicios comerciales de rutas.

\begin{thebibliography}{99}

\bibitem{Jazemi2023}
R. Jazemi, E. Alidadiani, K. Ahn y J. Jang, ``A review of literature on vehicle routing problems of
last-mile delivery in urban areas,'' \textit{Applied Sciences}, vol.~13, no.~24, art.~13015, 2023,
doi: 10.3390/app132413015.

\bibitem{Reyes2018}
D. Reyes, A. Erera, M. Savelsbergh, S. Sahasrabudhe y R. O'Neil, ``The meal delivery routing
problem,'' preprint, \textit{Optimization Online}, 2018. [En línea]. Disponible:
https://optimization-online.org/2018/04/6571/

\bibitem{Yildiz2019}
B. Yıldız y M. Savelsbergh, ``Provably high-quality solutions for the meal delivery routing
problem,'' \textit{Transportation Science}, vol.~53, no.~5, pp.~1372--1388, 2019,
doi: 10.1287/trsc.2018.0887.

\bibitem{Ulmer2021}
M. W. Ulmer, B. W. Thomas, A. M. Campbell y N. Woyak, ``The restaurant meal delivery problem: Dynamic
pickup and delivery with deadlines and random ready times,'' \textit{Transportation Science},
vol.~55, no.~1, pp.~75--100, 2021, doi: 10.1287/trsc.2020.1000.

\bibitem{Hildebrandt2022}
F. D. Hildebrandt y M. W. Ulmer, ``Supervised learning for arrival time estimations in restaurant
meal delivery,'' \textit{Transportation Science}, vol.~56, no.~4, pp.~1058--1084, 2022,
doi: 10.1287/trsc.2021.1095.

\bibitem{Boyaci2021}
B. Boyacı, T. H. Dang y A. N. Letchford, ``Vehicle routing on road networks: How good is Euclidean
approximation?,'' \textit{Computers \& Operations Research}, vol.~129, art.~105197, 2021,
doi: 10.1016/j.cor.2020.105197.

\bibitem{Pereira2021}
R. H. M. Pereira, M. Saraiva, D. Herszenhut, C. K. V. Braga y M. W. Conway, ``r5r: Rapid realistic
routing on multimodal transport networks with R5 in R,'' \textit{Findings}, 2021,
doi: 10.32866/001c.21262.

\bibitem{Kuhn1955}
H. W. Kuhn, ``The Hungarian method for the assignment problem,'' \textit{Naval Research Logistics
Quarterly}, vol.~2, no.~1--2, pp.~83--97, 1955, doi: 10.1002/nav.3800020109.

\bibitem{Boeing2022}
G. Boeing, ``Street network models and indicators for every urban area in the world,''
\textit{Geographical Analysis}, vol.~54, no.~3, pp.~519--535, 2022, doi: 10.1111/gean.12281.

\bibitem{Geisberger2008}
R. Geisberger, P. Sanders, D. Schultes y D. Delling, ``Contraction hierarchies: Faster and simpler
hierarchical routing in road networks,'' en \textit{Experimental Algorithms (WEA 2008)}, Lecture
Notes in Computer Science, vol.~5038, Berlín: Springer, 2008, pp.~319--333,
doi: 10.1007/978-3-540-68552-4\_24.

\bibitem{LuxenVetter2011}
D. Luxen y C. Vetter, ``Real-time routing with OpenStreetMap data,'' en \textit{Proc. 19th ACM
SIGSPATIAL Int. Conf. Advances in Geographic Information Systems}, Chicago, IL, EE.~UU., 2011,
pp.~513--516, doi: 10.1145/2093973.2094062.

\bibitem{Wu2024}
L. Wu \textit{et al.}, ``LaDe: The first comprehensive last-mile express dataset from industry,''
en \textit{Proc. 30th ACM SIGKDD Conf. Knowledge Discovery and Data Mining}, 2024,
pp.~5991--6002, doi: 10.1145/3637528.3671548.

\bibitem{Wilcoxon1945}
F. Wilcoxon, ``Individual comparisons by ranking methods,'' \textit{Biometrics Bulletin}, vol.~1,
no.~6, pp.~80--83, 1945, doi: 10.2307/3001968.

\bibitem{Friedman1937}
M. Friedman, ``The use of ranks to avoid the assumption of normality implicit in the analysis of
variance,'' \textit{Journal of the American Statistical Association}, vol.~32, no.~200,
pp.~675--701, 1937, doi: 10.1080/01621459.1937.10503522.

\bibitem{MunozVillamizar2021}
A. Muñoz-Villamizar, E. L. Solano-Charris, L. Reyes-Rubiano y J. Faulin, ``Measuring disruptions in
last-mile delivery operations,'' \textit{Logistics}, vol.~5, no.~1, art.~17, 2021,
doi: 10.3390/logistics5010017.

\end{thebibliography}

\end{document}
"
\documentclass[10pt]{article}

% ── Revisión doble ciega ────────────────────────────────────────────
\newif\ifblind
\ifdefined\conautores \blindfalse \else \blindtrue \fi

% ── Página (plantilla RIIM: carta; márgenes 2.0 / 2.0 / 3.5 / 2.25 cm;
%    encabezado y pie a 1.25 cm del borde) ──────────────────────────
\usepackage[letterpaper, left=2cm, right=2cm, top=3.5cm, bottom=2.25cm,
            headheight=1.45cm, headsep=0.8cm, footskip=0.62cm]{geometry}

% ── Tipografía ──────────────────────────────────────────────────────
\usepackage{fontspec}
\setmainfont{Montserrat}[
  Path=fonts/, Extension=.ttf,
  UprightFont=*-Regular, ItalicFont=*-Italic,
  BoldFont=*-SemiBold, BoldItalicFont=*-SemiBoldItalic]
\newfontfamily\semibold{Montserrat}[
  Path=fonts/, Extension=.ttf,
  UprightFont=*-SemiBold, ItalicFont=*-SemiBoldItalic, Ligatures=TeX]
\newfontfamily\arialfont{Arial}[Ligatures=TeX]
\usepackage[spanish, mexico, es-tabla, es-nodecimaldot, es-noshorthands]{babel}
\spanishplainpercent
\usepackage{amsmath, amssymb}
\usepackage{xcolor}
\definecolor{riimazul}{HTML}{003D7C}
\definecolor{cfcfs}{HTML}{1E88E5}
\definecolor{crh}{HTML}{F59E0B}

\usepackage{graphicx}
\usepackage{booktabs}
\usepackage{tabularx}
\usepackage{float}
\usepackage{flafter}
% Permite que tablas y figuras compartan página con el texto en lugar de ir a páginas solo de flotantes.
\renewcommand{\topfraction}{0.85}
\renewcommand{\bottomfraction}{0.7}
\renewcommand{\textfraction}{0.1}
\renewcommand{\floatpagefraction}{0.8}
\setcounter{topnumber}{3}
\setcounter{bottomnumber}{2}
\setcounter{totalnumber}{4}
\setlength{\textfloatsep}{12pt plus 2pt minus 2pt}
\setlength{\floatsep}{10pt plus 2pt minus 2pt}
\setlength{\intextsep}{10pt plus 2pt minus 2pt}
\usepackage{tikz}
\usetikzlibrary{arrows.meta, positioning, fit, backgrounds}
\usepackage{fancyhdr}
\usepackage{caption}
\usepackage[hidelinks]{hyperref}

% Cuerpo: Montserrat 10, interlineado sencillo, justificado, sin sangría, 8 pt después.
\renewcommand{\normalsize}{\fontsize{10}{12.2}\selectfont}
\normalsize
\setlength{\parindent}{0pt}
\setlength{\parskip}{8pt plus 1pt}
\setlength{\abovedisplayskip}{6pt}
\setlength{\belowdisplayskip}{6pt}
\setlength{\abovedisplayshortskip}{4pt}
\setlength{\belowdisplayshortskip}{4pt}

% Títulos de sección sin numerar: Montserrat SemiBold 12 pt azul; subtítulos SemiBold itálica 11 pt.
\makeatletter
\setcounter{secnumdepth}{0}
\renewcommand{\section}{\@startsection{section}{1}{0pt}{8pt}{0.1pt}
  {\semibold\fontsize{12}{14.6}\selectfont\color{riimazul}\raggedright}}
\renewcommand{\subsection}{\@startsection{subsection}{2}{0pt}{8pt}{0.1pt}
  {\semibold\itshape\fontsize{11}{13.4}\selectfont\addfontfeatures{LetterSpace=6.8}\raggedright}}
\makeatother

% Pies de figura y de tabla: 9 pt, centrados, "Figura N." / "Tabla N.".
\captionsetup{font={footnotesize}, labelfont={footnotesize}, labelsep=period,
  justification=centering, skip=4pt}
\captionsetup[table]{position=top}
\renewcommand{\footnotesize}{\fontsize{9}{11}\selectfont}

% Encabezado: logotipo RIIM (imagen de la plantilla recortada al 59.3 % izquierdo, como en Word) a la izquierda y "RIIM, Vol 3, Num 2, 2026" (Arial 8) a la derecha.
% Pie: número de página centrado; ISSN a la izquierda y URL a la derecha (9 pt).
\pagestyle{fancy}
\fancyhf{}
\renewcommand{\headrulewidth}{0pt}
\renewcommand{\footrulewidth}{0pt}
\fancyhead[C]{\makebox[\textwidth]{%
  \raisebox{-0.5\height}{\includegraphics[trim=0 0 335.53bp 0, clip, width=7.37cm, height=1.40cm]{img/riim_header.png}}%
  \hfill\raisebox{-0.5\height}{\arialfont\fontsize{8}{9.6}\selectfont RIIM, Vol 3, Num 2, 2026}}}
\fancyfoot[C]{\parbox[t]{\textwidth}{\centering\fontsize{9}{11}\selectfont\thepage\\[1pt]
  ISSN 3122-4806\hfill http://orizaba.tecnm.mx/revistaojs}}

% Listas compactas.
\makeatletter
\renewcommand{\@listI}{\leftmargin=1.2em \topsep=0pt \partopsep=0pt \itemsep=2pt \parsep=0pt}
\let\@listi\@listI
\makeatother

% Referencias IEEE (Arial 10, sangría francesa de 1.27 cm).
\makeatletter
\renewenvironment{thebibliography}[1]{%
  \section{Referencias}
  \list{\@biblabel{\@arabic\c@enumiv}}{%
    \arialfont\fontsize{10}{12}\selectfont
    \settowidth\labelwidth{\@biblabel{#1}}%
    \leftmargin=1.27cm \labelwidth=1.27cm \labelsep=0pt \itemindent=0pt
    \itemsep=0pt \parsep=0pt \topsep=0pt
    \def\makelabel##1{##1\hfil}%
    \usecounter{enumiv}\let\p@enumiv\@empty
    \renewcommand\theenumiv{\@arabic\c@enumiv}}%
  \sloppy\clubpenalty4000\widowpenalty4000\sfcode`\.\@m}
  {\def\@noitemerr{\@latex@warning{Empty `thebibliography' environment}}\endlist}
\makeatother

\hypersetup{pdftitle={Logística de última milla en redes viales reales: MDRP reproducible en La Paz, B.C.S.},
  pdfauthor={}, pdfsubject={RIIM}, pdfkeywords={enrutamiento dinámico, horizonte rodante, OpenStreetMap, simulación}}
\ifblind\else\hypersetup{pdfauthor={Alan José Marcos Sánchez Martínez}}\fi

% Notación.
\newcommand{\CtD}{\text{CtD}}
\newcommand{\Cdisp}{\mathcal{C}_{disp}(t)}

% Diagrama de caja: \caja{x}{mín}{Q1}{mediana}{Q3}{máx}{color}
\newcommand{\caja}[7]{%
  \draw (#1,#2) -- (#1,#3) (#1,#5) -- (#1,#6);
  \draw (#1-0.14,#2) -- (#1+0.14,#2) (#1-0.14,#6) -- (#1+0.14,#6);
  \filldraw[fill=#7!75, draw=black!70] (#1-0.3,#3) rectangle (#1+0.3,#5);
  \draw[line width=1pt] (#1-0.3,#4) -- (#1+0.3,#4);}

\begin{document}

% ── Título, autores y área ──────────────────────────────────────────
\begin{center}
  \setlength{\parskip}{0pt}
  \vspace*{6pt}
  {\semibold\fontsize{16}{19.5}\selectfont\addfontfeatures{LetterSpace=-3.1}%
   LOGÍSTICA DE ÚLTIMA MILLA EN REDES VIALES REALES: MDRP REPRODUCIBLE EN LA PAZ, B.C.S.\par}
  \vspace{18pt}
  \ifblind\else
  {\itshape\fontsize{9}{11}\selectfont
   Sánchez Martínez, Alan José Marcos\textsuperscript{1*}\par
   \textsuperscript{1}Tecnológico Nacional de México / Instituto Tecnológico de La Paz, m16310322@lapaz.tecnm.mx\par}
  \fi
  {\fontsize{9}{11}\selectfont Área de participación: Sistemas Computacionales\par}
\end{center}

\section{Resumen}
Las plataformas de entrega de comida deciden sin conocer los pedidos futuros, si asignan cada orden
al repartidor libre más cercano o si esperan unos minutos para agrupar órdenes y coordinar la
asignación. Este trabajo evalúa en qué condiciones conviene coordinar en una ciudad media mexicana,
La Paz, B.C.S., y cuanto cambia la evaluación al usar la red vial real. Se construyó un simulador de
código abierto que obtiene los tiempos de viaje de OSRM sobre OpenStreetMap y compara una política
reactiva (FCFS) con una heurística de horizonte rodante (RH) en cuatro niveles de saturación de la
flota, con diez instancias sintéticas por nivel. Con una flota holgada, FCFS entrega 1.4~min antes;
con 26 o más pedidos por repartidor en la jornada, RH reduce el tiempo de entrega en todas las
instancias, hasta en 67\,\%. Las distancias euclidianas subestiman los tiempos de entrega y el
beneficio de coordinar.

{\semibold\fontsize{12}{14.6}\selectfont\color{riimazul}Palabras clave}: \textit{enrutamiento dinámico, horizonte rodante, OpenStreetMap, simulación}

\section{Abstract}
\begin{otherlanguage}{english}\frenchspacing
Food delivery platforms must decide, without knowing future orders, whether to assign each order to
the nearest idle courier or to wait a few minutes, batch orders and coordinate the assignment. This
work evaluates under which conditions coordination pays off in a mid-sized Mexican city, La Paz,
B.C.S., and how much the evaluation changes when the real road network is used. An open-source
simulator was built that obtains travel times from OSRM over OpenStreetMap and compares a reactive
policy (FCFS) with a rolling-horizon heuristic (RH) at four fleet saturation levels, with ten
synthetic instances per level. With a slack fleet, FCFS delivers 1.4~min earlier; with 26 or more
orders per courier over the day, RH reduces delivery time in every instance, by up to 67\,\%.
Euclidean distances underestimate both delivery times and the benefit of coordination.

{\semibold\fontsize{12}{14.6}\selectfont\color{riimazul}Key words}: \textit{dynamic routing, rolling horizon, OpenStreetMap, simulation}
\end{otherlanguage}

\section{Introducción}
La última milla, el tramo que va del punto de recolección al domicilio del cliente, es uno de los
eslabones mas costosos y difíciles de la logística urbana. En zonas residenciales concentra grandes
volúmenes de entregas con horarios estrechos y condiciones que cambian durante el día
\cite{Jazemi2023}. La entrega de comida a domicilio agrava estas exigencias: los pedidos llegan sin
aviso, la comida pierde calidad con cada minuto de espera y el cliente espera recibirla en menos de
una hora.

En este contexto, una plataforma debe decidir de forma continua qué repartidor atiende cada orden y
en qué momento. Este problema, formalizado como \textit{Meal Delivery Routing Problem} (MDRP) por
Reyes \textit{et al.} \cite{Reyes2018}, admite dos estrategias básicas. La primera es reactiva: cada
orden se asigna en cuanto aparece al repartidor libre más cercano (\textit{First-Come, First-Served},
FCFS). La segunda es coordinada: el sistema espera unos minutos, agrupa órdenes del mismo restaurante
y decide de manera conjunta qué repartidor atiende cada agrupamiento, como en la heurística de
horizonte rodante (\textit{Rolling Horizon}, RH) propuesta en \cite{Reyes2018}. Coordinar tiene un
costo, porque algunas órdenes esperan a que se forme el agrupamiento o a que se entreguen otras antes
que ellas, y un beneficio, porque cada viaje atiende más pedidos. El balance entre ambos depende de
cuántos repartidores hay frente a la demanda.

La evidencia disponible sobre este balance proviene de entornos muy distintos a una ciudad media. RH
se evaluó con instancias que reproducen áreas metropolitanas de Estados Unidos y con tiempos de viaje
calculados como la distancia euclidiana multiplicada por un factor constante \cite{Reyes2018}. En
instancias pequeñas resueltas a optimalidad, su \textit{Click-to-Door} promedio queda apenas un
4\,\% por encima del óptimo \cite{Yildiz2019}, pero no se ha reportado cómo cambia su ventaja frente
a una política reactiva conforme varía el tamaño de la flota. Otras formulaciones del problema
incorporan la incertidumbre en los tiempos de preparación \cite{Ulmer2021} o la estimación del tiempo
de llegada, cuya desviación respecto a lo prometido pesa en la satisfacción del cliente tanto como la
rapidez misma \cite{Hildebrandt2022}.

La forma de medir los tiempos de viaje añade una segunda fuente de error. Las distancias por la red
vial suelen ser entre 1.2 y 1.4 veces mayores que en línea recta, y en problemas con ventanas de
tiempo estrechas, como la entrega de comida, una ruta factible con distancias euclidianas puede
resultar infactible sobre la red real \cite{Boyaci2021}. Una evaluación que ignore la red puede
equivocarse no solo en los tiempos de entrega, sino también en la carga efectiva de la flota y, por
lo tanto, en el beneficio de coordinar. Los servicios comerciales de rutas resuelven la precisión,
pero su costo de licencia limita su uso en investigación \cite{Pereira2021} y su funcionamiento
cerrado dificulta reproducir los experimentos.

La Paz, B.C.S., ilustra esta brecha: es una ciudad media sin datos públicos de operaciones de entrega
de comida, cuya red vial, condicionada por la costa y por pocas avenidas principales, obliga a
recorridos más largos que la línea recta. De ahí surge la pregunta que guía este trabajo:
\textit{en una ciudad media, ¿en qué condiciones de saturación de la flota conviene coordinar los
despachos en lugar de asignar cada pedido al llegar, y cuánto cambia esa evaluación al usar la red
vial real en lugar de distancias euclidianas?} La pregunta se divide en tres: cómo varía la ventaja
de RH con la saturación y a partir de qué nivel cambia la política más conveniente (PI1); cuánto más largos son los
trayectos por la red y cuánto cambia la comparación al ignorarla (PI2); y si la evaluación puede
hacerse de forma reproducible, con recursos modestos y solo con herramientas de código abierto (PI3).

Las contribuciones corresponden a esas tres preguntas: (i)~una comparación pareada de FCFS y RH en
cuatro niveles de saturación, con prueba estadística por instancia, que ubica el punto a partir del
cual coordinar compensa; (ii)~la medición de la circuidad de la red urbana de La Paz y dos escenarios
de control con distancias euclidianas que separan el efecto de la red del efecto de la política; y
(iii)~un entorno de simulación reproducible, construido sobre OpenStreetMap y OSRM, que ejecuta un
día de operación en alrededor de un minuto en una computadora portátil. El estudio no propone un
algoritmo nuevo ni compara RH con políticas del estado del arte: su objetivo es caracterizar cuándo
conviene coordinar en un entorno realista.

\section{Metodología}

\subsection{Formulación del problema y métricas}
Se adopta la formulación del MDRP de \cite{Reyes2018}. Un conjunto $\mathcal{R}$ de restaurantes con
ubicación fija recibe órdenes $o \in \mathcal{O}$, cada una con restaurante de origen $r_o$, instante
de colocación $a_o$, instante en que la comida está lista $e_o \geq a_o$ y punto de entrega $\ell_o$.
La orden solo se conoce a partir de $a_o$. Una flota $\mathcal{C}$ de repartidores, con turno
$[e_c, l_c]$, recoge y entrega las órdenes. Órdenes del mismo restaurante pueden formar un
agrupamiento $B$ que se recoge en una sola visita y se entrega en una ruta de varias paradas. Los
tiempos de servicio son $s_r$ en el restaurante, una vez por agrupamiento, y $s_o$ con cada cliente.

La métrica principal es el \textit{Click-to-Door} de cada orden,
$\CtD_o = t_{delivery,o} - a_o$. A partir de ella se reportan su media y su percentil 90, y el
cumplimiento del nivel de servicio (SLA, \textit{Service Level Agreement}), entendido como el
porcentaje de órdenes entregadas dentro del objetivo $\tau$:
\begin{equation}
  \text{SLA} = \frac{|\{o \in \mathcal{O}_{del} : \CtD_o \leq \tau\}|}{|\mathcal{O}_{del}|} \times 100,
  \label{eq:sla}
\end{equation}
donde $\mathcal{O}_{del}$ es el conjunto de órdenes entregadas. Como ambas métricas se calculan
sobre órdenes entregadas, se reporta también el porcentaje de órdenes sin entregar al cierre. Se
complementan con el tiempo que la comida lista espera a ser recogida (RtP, \textit{Ready-to-Pickup}),
$\text{RtP}_o = t_{pickup,o} - e_o$, la
distancia recorrida por pedido, la utilización de la flota (fracción del turno en conducción o en
servicio), el tamaño de agrupamiento y el costo por pedido, con un pago de $p_1$ por orden entregada
y un mínimo garantizado de $p_2$ por hora de turno \cite{Reyes2018}.

\subsection{Políticas de despacho}
El simulador avanza en epocas de $f = 5$~min para ambas políticas: en cada época, FCFS asigna una por
una las órdenes nuevas y RH resuelve su optimización. Las dos calculan los instantes de recolección y
entrega con las mismas reglas. Un agrupamiento $B$ solo puede recogerse cuando ya está lista la
última de sus órdenes, es decir, a partir del instante $e_B = \max_{o\in B} e_o$. El repartidor $c$
llega al restaurante en el instante $t_{arr,r}$ y, como en \cite{Reyes2018}, cada tiempo de servicio
se reparte mitad antes y mitad después del momento de recoger o de entregar. Si la comida ya está
lista, la recolección ocurre al cumplirse la primera mitad del servicio en el restaurante; si no, el
repartidor espera hasta que esté lista. Por eso el instante de recolección es
$t_{pickup}(c,B) = \max\bigl(e_B,\; t_{arr,r} + s_r/2\bigr)$. Después, el repartidor sale del
restaurante al terminar la segunda mitad del servicio, recorre el tramo hasta el primer cliente y le
entrega al cumplirse la primera mitad del servicio con él; cada entrega siguiente se calcula igual a
partir de la anterior.

\textbf{FCFS (línea base reactiva).} En cada época, las órdenes reveladas sin repartidor se procesan
por orden de colocación. Cada una se asigna al repartidor libre más cercano en línea recta al
restaurante (esta distancia solo sirve para escoger al repartidor); la ruta repartidor $\to$ restaurante
$\to$ cliente y sus tiempos se obtienen de OSRM, y el repartidor queda comprometido de inmediato. Cada
viaje lleva una sola orden.

\textbf{RH (política coordinada).} Sigue el esquema de \cite{Reyes2018} y en cada época $t$ ejecuta
cinco pasos. Primero, delimita las órdenes que entran a la optimización: de las órdenes ya reveladas y
aún no entregadas, $\mathcal{O}_t$, solo considera las que estarán listas dentro del horizonte de
asignación $\Delta_u$, es decir, $U_t = \{o \in \mathcal{O}_t : e_o \leq t + \Delta_u\}$; planear
órdenes cuya comida tardará mucho en estar lista sería poco útil en un entorno que cambia cada pocos
minutos.

Segundo, estima cuántas órdenes conviene juntar en cada viaje. El tamaño objetivo de agrupamiento
$Z_t$ compara la demanda próxima con la oferta próxima: el numerador cuenta las órdenes que estarán
listas en los siguientes $\Delta_1$ minutos, y el denominador, los repartidores en turno, libres u
ocupados, que quedarán disponibles en los siguientes $\Delta_2$ minutos, donde $e^{avail}_c$ es el
instante en que el repartidor $c$ termina su ruta actual:
\begin{equation}
  Z_t = \left\lceil \frac{|\{o \in \mathcal{O}_t : e_o \leq t + \Delta_1\}|}
                        {|\{c \in \mathcal{C}_{turno}(t) : e^{avail}_c \leq t + \Delta_2\}|} \right\rceil.
  \label{eq:zt}
\end{equation}
Si hay tantos repartidores como ordenes, $Z_t = 1$ y cada orden tiende a viajar sola; si la flota se
satura, $Z_t$ crece y se forman agrupamientos más grandes. $Z_t$ se acota en 5, y ningún agrupamiento
puede superar $B_{max} = 4$ órdenes.

Tercero, forma los agrupamientos de cada restaurante $r$. Sus órdenes $U_{t,r}$ se reparten en
\[
  m_r = \min\bigl(\lceil |U_{t,r}|/Z_t \rceil,\; |\Cdisp|\bigr)
\]
agrupamientos, de modo que cada uno reciba alrededor de $Z_t$ órdenes sin crear más agrupamientos que
repartidores disponibles. Las órdenes se toman por \textit{ready time}, y cada una se inserta en el
agrupamiento y la posición de la ruta que menos incrementan el costo
\begin{equation}
  \text{costo}(B, o, p) = T_{ruta}(B \oplus_p o) + \beta\, T_{servicio}(B \oplus_p o),
  \label{eq:costo}
\end{equation}
donde $B \oplus_p o$ es el agrupamiento con la orden $o$ insertada en la posición $p$, $T_{ruta}$ es
la duración de la ruta según OSRM y $T_{servicio}$ el tiempo de servicio acumulado. El peso $\beta$, que en
\cite{Reyes2018} representa la pérdida de frescura y se calibra en $[0, 1]$, se fija en 0.1 para que la
duración de la ruta domine el costo. Después, dos iteraciones de extracción y reinserción mueven
cada orden al agrupamiento donde su costo es menor.

Cuarto, asigna agrupamientos a repartidores. Los agrupamientos se clasifican en tres grupos de
prioridad: urgentes, cuando ningún repartidor puede entregar su orden más antigua dentro del objetivo
$\tau$; retrasados, cuando el objetivo aún es alcanzable pero ningún repartidor llega al restaurante
antes de que la comida esté lista; y normales. Los grupos se resuelven en ese orden, para que los
urgentes obtengan primero a los repartidores más cercanos. En cada grupo se resuelve un problema de
asignación lineal, en el que cada repartidor recibe a lo más un agrupamiento, con el peso
\begin{equation}
  w_{c,B} = \frac{|B|}{T_{delivery}(c,B)} - \theta\,\bigl(t_{pickup}(c,B) - e_B\bigr) - P_g.
  \label{eq:peso}
\end{equation}
El primer término mide la eficiencia del viaje en órdenes entregadas por hora, donde $T_{delivery}(c,B)$
es el tiempo desde que $c$ queda libre hasta la última entrega de $B$. El segundo penaliza, con
$\theta = 0.1$, los minutos que la comida espera lista en el restaurante, y $P_g$ resta 100 si el par
excede el máximo absoluto $\tau_{max}$ y 50 si la recolección se retrasa. La asignación de máximo peso
se obtiene de forma exacta con un algoritmo de tipo húngaro \cite{Kuhn1955}.

Quinto, el compromiso en dos etapas de \cite{Reyes2018} decide qué parte de la asignación se
ejecuta. Si el repartidor llega al restaurante y todas las órdenes estarán listas dentro de la época,
o si alguna orden lleva más de $\Delta_x$ minutos lista, se fija la ruta completa. En otro caso, el
repartidor solo se dirige al restaurante, como un preposicionamiento, y en la época siguiente el
agrupamiento puede crecer o cambiar de repartidor con las órdenes que se hayan revelado.

\subsection{Entorno de simulación reproducible}
La Figura~\ref{fig:arquitectura} resume el entorno. El extracto de OpenStreetMap de Baja California
Sur, modelo de datos viales abierto y de cobertura mundial \cite{Boeing2022}, se preprocesa con el
algoritmo \textit{Contraction Hierarchies} \cite{Geisberger2008}, que permite resolver cada consulta
de camino mínimo en milisegundos. El servidor OSRM \cite{LuxenVetter2011}, versión 5.27.1, atiende
las consultas de ruta del simulador por HTTP; su imagen oficial de Docker permite reproducir el
entorno en cualquier sistema operativo. El simulador, escrito en Python, avanza por épocas, guarda en
memoria las rutas ya consultadas y produce, para cada corrida, los registros por orden y por
repartidor de los que se derivan los indicadores.

\begin{figure}[htbp]
\centering
\begin{tikzpicture}[
  font=\footnotesize,
  caja/.style={draw=black!70, rounded corners=2pt, align=center, minimum height=1.05cm,
               text width=#1, fill=white},
  caja/.default=2.8cm,
  flecha/.style={-{Stealth[length=5pt]}, thick, black!70},
  capa/.style={draw=#1!70!black, rounded corners=4pt, fill=#1!8, inner xsep=7pt, inner ysep=7pt}]
  \node[caja=2.8cm] (osm)   at (0,0)      {Extracto\\OpenStreetMap};
  \node[caja=3.0cm] (ch)    at (3.55,0)    {Preprocesamiento\\\textit{Contraction}\\\textit{Hierarchies}};
  \node[caja=2.8cm] (osrm)  at (7.1,0)    {Servidor OSRM\\(API de rutas)};
  \node[caja=2.8cm] (gen)   at (0,-2.3)   {Generador de\\instancias (PPNH)};
  \node[caja=2.8cm] (sim)   at (3.55,-2.3) {Simulador por\\épocas ($f = 5$ min)};
  \node[caja=2.8cm] (pol)   at (7.1,-2.3) {Política\\FCFS o RH};
  \node[caja=2.8cm] (kpi)   at (10.9,-2.3) {Registros por orden\\y 21 indicadores};
  \node[caja=2.8cm] (est)   at (10.9,0)   {Análisis estadístico\\(Wilcoxon, Friedman)};
  \draw[flecha] (osm) -- (ch);
  \draw[flecha] (ch) -- (osrm);
  \draw[flecha] (gen) -- (sim);
  \draw[flecha, {Stealth[length=5pt]}-{Stealth[length=5pt]}] (sim) -- (pol);
  \draw[flecha, dashed, {Stealth[length=5pt]}-{Stealth[length=5pt]}] (osrm) --
     node[right, font=\scriptsize, align=left] {tiempos de\\viaje (HTTP)} (pol);
  \draw[flecha] (pol) -- (kpi);
  \draw[flecha] (kpi) -- (est);
  \begin{scope}[on background layer]
    \node[capa=blue, fit=(osm)(ch)(osrm)] {};
    \node[capa=green, fit=(gen)(sim)(pol)] {};
    \node[capa=orange, fit=(kpi)(est)] {};
  \end{scope}
\end{tikzpicture}
\caption{Arquitectura del entorno: enrutamiento sobre la red vial (azul), simulación (verde) y
análisis (naranja).}
\label{fig:arquitectura}
\end{figure}

\subsection{Generación de instancias}
Al no existir datos públicos de entregas de comida en La Paz, las instancias se generan de forma
sintética. Las llegadas siguen un proceso de Poisson no homogéneo con intensidad bimodal
\begin{equation}
  \lambda(t) = \lambda_0
  + A_m \exp\!\left(-\frac{(t-\mu_m)^2}{2\sigma^2}\right)
  + A_e \exp\!\left(-\frac{(t-\mu_e)^2}{2\sigma^2}\right),
  \label{eq:lambda}
\end{equation}
siguiendo el patrón bimodal observado en entregas de última milla \cite{Wu2024}. A falta de datos
locales, los parámetros son supuestos: ventana operativa de 08:00 a 19:00, picos en $\mu_m = $ 09:00
y $\mu_e = $ 17:00, $\lambda_0 = 0.6$ órdenes/min, $A_m = 2.4$, $A_e = 2.0$ y $\sigma = 45$~min. El perfil se escala para
obtener alrededor de mil órdenes por día (entre 988 y 1\,047 según la semilla). El tiempo de
preparación es $\max(5, \mathcal{N}(12, 3^2))$~min. Los 25 restaurantes y los clientes se ubican
por muestreo uniforme dentro de un polígono de 8.4~km$^2$ del centro urbano, y cada cliente se asocia
a un restaurante situado a entre 2 y 5~km en línea recta. Todos los repartidores comparten un turno
que inicia 15~min antes del primer pedido y termina una hora después del último \textit{ready time},
y parten del centroide de los restaurantes. La flota se dimensiona a partir de la razón de
saturación $\rho$, el número objetivo de pedidos por repartidor en la jornada:
\begin{equation}
  |\mathcal{C}| = \max\!\left(n_{\min},\; \left\lceil |\mathcal{O}| / \rho \right\rceil\right),
  \qquad n_{\min} = 10.
  \label{eq:flota}
\end{equation}
Fijar $\rho$, y no el número de repartidores, mantiene la misma carga relativa en todas las semillas
de un escenario aunque el número de órdenes varíe.

\subsection{Diseño experimental}
La Tabla~\ref{tab:parametros} reúne los parámetros, comunes a ambas políticas, y los cuatro
escenarios de saturación. Cada semilla (2025 a 2034) genera una instancia, y sobre ella se ejecutan
FCFS y RH con la misma flota; como el simulador no tiene componentes aleatorios propios, lo único que
cambia dentro de cada par es la política. Esto da 80 corridas sobre la red vial. Para evitar la
pseudorreplicación, cada métrica se promedia por semilla y la comparación se hace con la prueba de
rangos con signo de Wilcoxon pareada \cite{Wilcoxon1945} sobre las diez diferencias de cada
escenario, con intervalo de confianza $t$ al 95\,\% y la correlación biserial de rangos $r_{rb}$ como
tamaño de efecto. La prueba de Friedman \cite{Friedman1937} sobre la matriz de diferencias
$10 \times 4$ evalúa si el efecto cambia entre escenarios.

\begin{table}[htbp]
\centering
\footnotesize
\caption{Parámetros del simulador y escenarios de saturación.}
\label{tab:parametros}
\setlength{\tabcolsep}{4pt}
\begin{tabular}{@{}l c c @{\hspace{0.4cm}} l c c@{}}
\toprule
Parámetro & Símbolo & Valor & Parámetro & Símbolo & Valor \\
\midrule
Frecuencia de optimización & $f$ & 5 min & Penalización de frescura & $\beta$ & 0.1 \\
Horizonte de asignación & $\Delta_u$ & 20 min & Penalización por espera & $\theta$ & 0.1 \\
Anticipación (demanda, oferta) & $\Delta_1$, $\Delta_2$ & 20 min & Espera para forzar compromiso & $\Delta_x$ & 10 min \\
CtD objetivo / máximo & $\tau$, $\tau_{max}$ & 40/90 min & Tamaño máximo de agrupamiento & $B_{max}$ & 4 \\
Tiempos de servicio & $s_r$, $s_o$ & 4 min & Pago por orden / mínimo por hora & $p_1$, $p_2$ & \$10 / \$15 \\
\midrule
\multicolumn{6}{@{}l}{Escenarios: S1, $\rho = 18$ (55--59 repartidores); S2, $\rho = 26$ (38--41); S3, $\rho = 35$ (29--30); S4, $\rho = 52$ (19--21).} \\
\bottomrule
\end{tabular}
\end{table}

Para medir el efecto de la red vial se calculó la circuidad $\kappa = d_{OSRM}/d_{recta}$ sobre todos
los trayectos restaurante--cliente, y el experimento completo se repitió con dos controles en los que
el tiempo de viaje es la distancia en línea recta (Haversine) a velocidad constante: la velocidad
media que OSRM asigna a esos trayectos (651~m/min), que deja como única diferencia la geometría de
las rutas, y esa velocidad dividida entre $\kappa$ (472~m/min), que iguala la duración media de los
viajes. Cada control suma otras 80 corridas.

\section{Resultados y discusión}

\subsection{Calidad de servicio y cambio de la política más conveniente}
La Tabla~\ref{tab:calidad} muestra los indicadores de calidad de servicio. Con la flota holgada de S1,
FCFS entrega en promedio 1.4~min antes que RH: con repartidores de sobra, llevar cada pedido directo
del restaurante al cliente es más rápido que esperar a formar un agrupamiento, y en un agrupamiento
cada entrega espera a las anteriores. Aun así, el cumplimiento del SLA es prácticamente igual en ambas
políticas. A partir de S2 la relación se invierte: RH reduce el \textit{Click-to-Door} medio en
36.8\,\%, el percentil 90 en 46.0\,\% y eleva el cumplimiento del SLA en 31.5 puntos porcentuales. La
ventaja crece en S3 y S4, donde FCFS deja sin entregar al cierre el 11.3\,\% y el 39.9\,\% de los
pedidos, mientras que RH entrega prácticamente todos. Como el CtD y el SLA se calculan sobre órdenes
entregadas, la tabla favorece a FCFS en esos escenarios: medido sobre todas las órdenes generadas, su
cumplimiento cae a 6.4\,\% en S3 y 3.8\,\% en S4, frente a 57.6\,\% y 42.4\,\% bajo RH.

\begin{table}[htbp]
\centering
\footnotesize
\caption{Calidad de servicio, promedio de 10 semillas. $\Delta$ es el cambio de RH respecto a FCFS.}
\label{tab:calidad}
\setlength{\tabcolsep}{3pt}
\begin{tabular*}{\textwidth}{@{\extracolsep{\fill}}l rrr rrr rrr rr@{}}
\toprule
& \multicolumn{3}{c}{CtD medio (minutos)} & \multicolumn{3}{c}{CtD percentil 90 (minutos)} & \multicolumn{3}{c}{SLA (\%)} & \multicolumn{2}{c}{Sin entregar (\%)} \\
\cmidrule(lr){2-4} \cmidrule(lr){5-7} \cmidrule(lr){8-10} \cmidrule(l){11-12}
Esc. & FCFS & RH & $\Delta$ (\%) & FCFS & RH & $\Delta$ (\%) & FCFS & RH & $\Delta$ (pp) & FCFS & RH \\
\midrule
S1 & 28.93  & 30.31 & $+4.8$  & 43.33  & 43.80  & $+1.1$  & 84.6 & 83.4 & $-1.1$  & 0.0  & 0.0 \\
S2 & 56.07  & 35.42 & $-36.8$ & 97.91  & 52.87  & $-46.0$ & 38.1 & 69.6 & $+31.5$ & 0.1  & 0.0 \\
S3 & 108.49 & 45.36 & $-58.2$ & 169.54 & 79.10  & $-53.3$ & 7.3  & 57.6 & $+50.3$ & 11.3 & 0.0 \\
S4 & 224.88 & 73.63 & $-67.3$ & 321.10 & 186.18 & $-42.0$ & 6.2  & 42.5 & $+36.3$ & 39.9 & 0.2 \\
\bottomrule
\end{tabular*}
\end{table}

La Figura~\ref{fig:boxplot} muestra la distribución completa del \textit{Click-to-Door}. En S1 ambas
cajas son casi idénticas y quedan por debajo del objetivo de 40~min; desde S2, la distribución de
FCFS se alarga hacia valores de varias horas, mientras que la de RH se mantiene compacta y con la
mediana cerca del objetivo.

\begin{figure}[htbp]
\centering
\begin{tikzpicture}[y=0.0100cm, x=1cm, font=\footnotesize]
  \foreach \p/\esc in {0/S1 ($\rho=18$), 1/S2 ($\rho=26$), 2/S3 ($\rho=35$), 3/S4 ($\rho=52$)} {
    \begin{scope}[xshift=\p*4.15cm]
      \draw[black!60] (0,0) rectangle (3.6,390);
      \foreach \y in {0,100,200,300} \draw[black!12] (0,\y) -- (3.6,\y);
      \draw[red!70, dashed] (0,40) -- (3.6,40);
      \node[above] at (1.8,390) {\esc};
      \node[below] at (1.1,0) {FCFS};
      \node[below] at (2.5,0) {RH};
    \end{scope}}
  \foreach \y in {0,100,200,300} \node[left] at (0,\y) {\y};
  \node[rotate=90] at (-1.25,195) {\textit{Click-to-Door} (minutos)};
  % S1
  \caja{1.1}{10.9}{21.7}{26.4}{35.1}{55.1}{cfcfs}
  \caja{2.5}{11.4}{22.8}{28.5}{36.1}{56.1}{crh}
  % S2
  \begin{scope}[xshift=4.15cm]
  \caja{1.1}{11.9}{27.3}{54.3}{80.6}{121.7}{cfcfs}
  \caja{2.5}{11.4}{24.5}{32.5}{42.2}{68.7}{crh}
  \end{scope}
  % S3
  \begin{scope}[xshift=8.3cm]
  \caja{1.1}{11.9}{72.1}{109.1}{148.2}{199.9}{cfcfs}
  \caja{2.5}{11.4}{26.6}{36.5}{50.3}{85.8}{crh}
  \end{scope}
  % S4
  \begin{scope}[xshift=12.45cm]
  \caja{1.1}{14.9}{141.0}{266.0}{306.5}{367.5}{cfcfs}
  \caja{2.5}{12.3}{31.1}{43.8}{79.3}{151.5}{crh}
  \end{scope}
\end{tikzpicture}
\caption{Distribución del \textit{Click-to-Door} de las órdenes entregadas en las diez semillas
(cajas: cuartiles; bigotes: 1.5 veces el rango intercuartílico). La línea discontinua marca el
objetivo $\tau = 40$~min.}
\label{fig:boxplot}
\end{figure}

La Tabla~\ref{tab:wilcoxon} confirma, a nivel de instancia, el cambio de la política más conveniente. En S2, S3 y S4, RH
supera a FCFS en las diez semillas, tanto en CtD como en SLA, con $p = 0.0020$, el valor mínimo que
admite la prueba con $n = 10$, y $|r_{rb}| = 1$. En S1 la ventaja de FCFS en CtD es significativa
pero pequeña, y la diferencia en SLA no lo es. La prueba de Friedman confirma que la magnitud del
efecto cambia con la saturación ($\chi^2_F = 30.0$, $p < 0.001$). La Figura~\ref{fig:saturacion}
resume la respuesta a PI1: las curvas de CtD se cruzan entre $\rho = 18$ y $\rho = 26$, y por encima
de ese punto la ventaja de coordinar crece con la carga.

\begin{table}[htbp]
\centering
\footnotesize
\caption{Wilcoxon pareada por semilla ($n = 10$), $\bar{\Delta} = \bar{x}_{FCFS} - \bar{x}_{RH}$: favorecen a RH
$\bar{\Delta} > 0$ en CtD y $\bar{\Delta} < 0$ en SLA.}
\label{tab:wilcoxon}
\begin{tabular}{@{}ll rr l c r@{}}
\toprule
Esc. & Métrica & $\bar{x}_{FCFS}$ & $\bar{x}_{RH}$ & $\bar{\Delta}$ [IC 95\,\%] & $p$ & $r_{rb}$ \\
\midrule
S1 & CtD (minutos) & 28.93  & 30.31 & $-1.38$ [$-2.35$, $-0.41$]        & 0.0137 & $-0.85$ \\
S1 & SLA (\%)  & 84.6   & 83.4  & $+1.1$ [$-1.8$, $+4.1$]           & 0.2324 & $+0.45$ \\
S2 & CtD (minutos) & 56.07  & 35.42 & $+20.64$ [$+18.98$, $+22.31$]     & 0.0020 & $+1.00$ \\
S2 & SLA (\%)  & 38.1   & 69.6  & $-31.5$ [$-33.7$, $-29.2$]        & 0.0020 & $-1.00$ \\
S3 & CtD (minutos) & 108.49 & 45.36 & $+63.13$ [$+56.14$, $+70.13$]     & 0.0020 & $+1.00$ \\
S3 & SLA (\%)  & 7.3    & 57.6  & $-50.3$ [$-52.6$, $-48.0$]        & 0.0020 & $-1.00$ \\
S4 & CtD (minutos) & 224.88 & 73.63 & $+151.24$ [$+143.71$, $+158.78$]  & 0.0020 & $+1.00$ \\
S4 & SLA (\%)  & 6.2    & 42.5  & $-36.3$ [$-37.2$, $-35.3$]        & 0.0020 & $-1.00$ \\
\bottomrule
\end{tabular}
\end{table}

\begin{figure}[htbp]
\centering
\begin{tikzpicture}[font=\footnotesize]
  % Panel izquierdo: CtD medio
  \begin{scope}[x=0.16cm, y=0.018cm]
    \draw[black!60] (14,0) rectangle (56,240);
    \foreach \y in {0,50,100,150,200} {\draw[black!12] (14,\y) -- (56,\y); \node[left] at (14,\y) {\y};}
    \foreach \x in {18,26,35,52} {\draw[black!12] (\x,0) -- (\x,240); \node[below] at (\x,0) {\x};}
    \draw[red!70, dashed] (14,40) -- (56,40);
    \draw[cfcfs, very thick] (18,28.93) -- (26,56.07) -- (35,108.49) -- (52,224.88);
    \foreach \x/\y in {18/28.93, 26/56.07, 35/108.49, 52/224.88} \fill[cfcfs] (\x,\y) circle (2pt);
    \draw[crh, very thick] (18,30.31) -- (26,35.42) -- (35,45.36) -- (52,73.63);
    \foreach \x/\y in {18/30.31, 26/35.42, 35/45.36, 52/73.63} \fill[crh] (\x-0.45,\y-4.5) rectangle (\x+0.45,\y+4.5);
    \node[rotate=90] at (8.5,120) {CtD medio (minutos)};
    \node at (35,-30) {Pedidos por repartidor en la jornada ($\rho$)};
  \end{scope}
  % Panel derecho: SLA
  \begin{scope}[xshift=8.6cm, x=0.16cm, y=0.043cm]
    \draw[black!60] (14,0) rectangle (56,100);
    \foreach \y in {0,20,40,60,80,100} {\draw[black!12] (14,\y) -- (56,\y); \node[left] at (14,\y) {\y};}
    \foreach \x in {18,26,35,52} {\draw[black!12] (\x,0) -- (\x,100); \node[below] at (\x,0) {\x};}
    \draw[cfcfs, very thick] (18,84.6) -- (26,38.1) -- (35,7.3) -- (52,6.2);
    \foreach \x/\y in {18/84.6, 26/38.1, 35/7.3, 52/6.2} \fill[cfcfs] (\x,\y) circle (2pt);
    \draw[crh, very thick] (18,83.4) -- (26,69.6) -- (35,57.6) -- (52,42.5);
    \foreach \x/\y in {18/83.4, 26/69.6, 35/57.6, 52/42.5} \fill[crh] (\x-0.45,\y-1.9) rectangle (\x+0.45,\y+1.9);
    \node[rotate=90] at (8.5,50) {Cumplimiento del SLA (\%)};
    \node at (35,-12.5) {Pedidos por repartidor en la jornada ($\rho$)};
    % Leyenda
    \draw[cfcfs, very thick] (34,93) -- (38,93); \fill[cfcfs] (36,93) circle (2pt); \node[right] at (38,93) {FCFS};
    \draw[crh, very thick] (46,93) -- (50,93); \fill[crh] (47.55,91.1) rectangle (48.45,94.9); \node[right] at (50,93) {RH};
  \end{scope}
\end{tikzpicture}
\caption{\textit{Click-to-Door} medio y cumplimiento del SLA según la saturación de la flota. La línea
discontinua marca el objetivo $\tau = 40$~min.}
\label{fig:saturacion}
\end{figure}

\subsection{Eficiencia operativa y agrupamiento}
La Tabla~\ref{tab:eficiencia} explica el origen de la ventaja de RH. Aun en S1, RH reduce la espera
de la comida lista en el restaurante (RtP), porque en los picos de demanda la flota se ocupa por completo y
FCFS, con un pedido por viaje, acumula órdenes sin recoger, mientras que RH recoge varias en una sola
visita. Al agrupar, RH recorre entre 13.5\,\% y 49.4\,\% menos distancia por pedido y mantiene una
utilización de la flota menor que FCFS en S1--S3: con la misma flota, entrega lo mismo con menos
tiempo de conducción y, en S4, convierte ese margen en pedidos entregados. El tamaño de agrupamiento
crece con la saturación; en S1 una cuarta parte de los pedidos (25.4\,\%) viaja sola y en S4 el
80.9\,\% viaja en agrupamientos de tres o cuatro. El costo por pedido solo difiere en S1, donde muchos
repartidores reciben el pago mínimo y RH concentra las entregas en menos de ellos; desde S2 ambos
convergen al pago por orden.

\begin{table}[htbp]
\centering
\footnotesize
\caption{Eficiencia operativa y tamaño del agrupamiento en que viaja cada pedido (promedio de 10 semillas).}
\label{tab:eficiencia}
\setlength{\tabcolsep}{3pt}
\begin{tabular*}{\textwidth}{@{\extracolsep{\fill}}l rr rr rr r rr@{}}
\toprule
& \multicolumn{2}{c}{RtP medio (minutos)} & \multicolumn{2}{c}{Distancia/pedido (km)} & \multicolumn{2}{c}{Utilización (\%)} & Agrupamiento & \multicolumn{2}{c}{Costo/pedido (\$)} \\
\cmidrule(lr){2-3} \cmidrule(lr){4-5} \cmidrule(lr){6-7} \cmidrule(lr){8-8} \cmidrule(l){9-10}
Esc. & FCFS & RH & FCFS & RH & FCFS & RH & RH & FCFS & RH \\
\midrule
S1 & 7.78   & 4.94  & 4.71 & 4.08 & 37.7 & 31.2 & 2.54 & 10.80 & 11.59 \\
S2 & 34.92  & 9.33  & 4.88 & 3.60 & 55.2 & 41.4 & 2.80 & 10.00 & 10.03 \\
S3 & 87.35  & 18.62 & 5.07 & 3.17 & 66.8 & 51.4 & 3.03 & 10.00 & 10.00 \\
S4 & 203.75 & 46.20 & 5.20 & 2.63 & 67.3 & 67.8 & 3.30 & 10.00 & 10.00 \\
\bottomrule
\end{tabular*}
\end{table}

\subsection{Efecto de la red vial}
Sobre los 10\,184 trayectos restaurante--cliente de las diez instancias, la distancia por la red es
en promedio 3.30~km y en línea recta 2.39~km: los trayectos reales son 38\,\% más largos (circuidad
$\kappa = 1.38$), dentro del rango de 1.2 a 1.4 que reporta la literatura \cite{Boyaci2021}.

Para medir cuánto cambia la comparación, se repitió el experimento con tiempos de viaje en línea
recta. A la velocidad media de la red (651~m/min, primer control), ambas políticas parecen más
rápidas: en S2, el CtD de FCFS baja de 56.1 a 41.8~min y el de RH, de 35.4 a 30.3~min. Como FCFS
mejora más, la ventaja de RH se reduce en los cuatro escenarios ($p = 0.002$); en S3, por ejemplo,
pasa de 63.1 a 37.1~min. El segundo control baja la velocidad hasta que los viajes en línea recta
duran en promedio lo mismo que por la red ($651/1.38 \approx 472$~m/min); con él desaparece entre el
62 y el 84\,\% de esa reducción. La mayor parte del efecto se debe, por tanto, a que los trayectos
reales son más largos y cargan más a la flota, y el resto, a la forma de la red ($p \leq 0.004$).
Ignorar la red hace parecer más rápidas a ambas políticas y subestima el beneficio de coordinar, lo
que responde a la pregunta PI2.

\subsection{Reproducibilidad y costo de cómputo}
Toda la cadena (OpenStreetMap, OSRM y Python) es de código abierto, las instancias se regeneran a
partir de semillas fijas y cada corrida pasa verificaciones automáticas de integridad: ninguna orden
se recoge antes de estar lista ni se entrega antes de recogerse, no hay entregas duplicadas y ninguna
ruta recurrió a la estimación de respaldo en línea recta (de haber fallado OSRM). Con 1\,038 órdenes, simular un día toma
entre 1.8 y 2.7~segundos con FCFS y entre 57 y 80~segundos con RH, incluidas las consultas a OSRM, con menos de
225~MB de memoria por proceso, en una computadora portátil con procesador Apple M1 Pro y 16~GB de
RAM. Esto responde a PI3: la evaluación es reproducible y no requiere infraestructura especializada.

\subsection{Limitaciones}
El reloj del simulador avanza en épocas de 5~min. En cada época primero se asignan los pedidos y
después se registran las rutas terminadas, de modo que un repartidor que termina a mitad de una época
queda libre hasta la siguiente y espera entre 5 y 10~min su nuevo viaje. Como esta espera ocurre una
vez por viaje, perjudica más a FCFS, que hace un viaje por pedido. Para comprobar que no determina el
resultado, se repitieron las 80 corridas sobre la red vial con el orden inverso (primero se registran
las rutas terminadas y luego se asigna): FCFS sigue ganando en S1 y RH en S2--S4 en todas las
semillas, y el cruce permanece entre S1 y S2. La reducción del CtD baja a 14.6\,\% en S2, 41.4\,\% en
S3 y 63.0\,\% en S4, por lo que las magnitudes dependen de este detalle de implementación, pero las
conclusiones se mantienen. Además, la demanda es
sintética: sigue el patrón bimodal de la recolección de paquetes en LaDe con parámetros supuestos, no
calibrados con la demanda local de comida, y el tipo de ciudad influye en el
comportamiento de la última milla \cite{MunozVillamizar2021}. Los tiempos de preparación y de
servicio son homogéneos, la red se considera sin congestión y las ubicaciones se muestrean de forma
uniforme dentro del dominio urbano.

\section{Trabajo a futuro}
Quedan pendientes seis líneas de trabajo: calibrar la demanda, los tiempos de preparación y la
ubicación de restaurantes y clientes con datos locales de una plataforma de entregas o de encuestas;
explorar valores de $\rho$ entre 18 y 26 para ubicar con precisión el punto de cruce y evaluar
políticas híbridas que activen la coordinación solo cuando la carga lo justifique; incorporar
velocidades dependientes de la hora para representar la congestión; ampliar el dominio a toda la
mancha urbana, incluidos los trayectos que rodean la ensenada de La Paz; comparar RH con políticas
anticipatorias basadas en aprendizaje automático sobre el mismo entorno; y evaluar la escalabilidad
del simulador con instancias de mayor tamaño.

\section{Conclusiones}
En las instancias evaluadas del centro urbano de La Paz, coordinar los despachos no compensa cuando
la flota es holgada, pero sí a partir de una saturación moderada. Con una flota dimensionada para 18
pedidos por repartidor en la jornada, FCFS entrega 1.4~min antes que RH y el cumplimiento del SLA es
equivalente. Con 26 o más pedidos por repartidor, RH reduce el tiempo de entrega y mejora el
cumplimiento del servicio en todas las semillas, y cuando la flota es escasa entrega prácticamente
todos los pedidos que FCFS deja pendientes. Para un operador, esto sugiere que la decisión de
coordinar debe tomarse en función de la carga esperada por repartidor y no de forma fija.

La red vial importa para esta decisión. Los trayectos reales son en promedio 38\,\% más largos que
en línea recta, y un modelo euclidiano subestima tanto los tiempos de entrega como el beneficio de
coordinar. Finalmente, toda la evaluación puede reproducirse con herramientas abiertas en una
computadora portátil, lo que la hace accesible para operadores y desarrolladores de ciudades medias
que no cuentan con servicios comerciales de rutas.

\begin{thebibliography}{99}

\bibitem{Jazemi2023}
R. Jazemi, E. Alidadiani, K. Ahn y J. Jang, ``A review of literature on vehicle routing problems of
last-mile delivery in urban areas,'' \textit{Applied Sciences}, vol.~13, no.~24, art.~13015, 2023,
doi: 10.3390/app132413015.

\bibitem{Reyes2018}
D. Reyes, A. Erera, M. Savelsbergh, S. Sahasrabudhe y R. O'Neil, ``The meal delivery routing
problem,'' preprint, \textit{Optimization Online}, 2018. [En línea]. Disponible:
https://optimization-online.org/2018/04/6571/

\bibitem{Yildiz2019}
B. Yıldız y M. Savelsbergh, ``Provably high-quality solutions for the meal delivery routing
problem,'' \textit{Transportation Science}, vol.~53, no.~5, pp.~1372--1388, 2019,
doi: 10.1287/trsc.2018.0887.

\bibitem{Ulmer2021}
M. W. Ulmer, B. W. Thomas, A. M. Campbell y N. Woyak, ``The restaurant meal delivery problem: Dynamic
pickup and delivery with deadlines and random ready times,'' \textit{Transportation Science},
vol.~55, no.~1, pp.~75--100, 2021, doi: 10.1287/trsc.2020.1000.

\bibitem{Hildebrandt2022}
F. D. Hildebrandt y M. W. Ulmer, ``Supervised learning for arrival time estimations in restaurant
meal delivery,'' \textit{Transportation Science}, vol.~56, no.~4, pp.~1058--1084, 2022,
doi: 10.1287/trsc.2021.1095.

\bibitem{Boyaci2021}
B. Boyacı, T. H. Dang y A. N. Letchford, ``Vehicle routing on road networks: How good is Euclidean
approximation?,'' \textit{Computers \& Operations Research}, vol.~129, art.~105197, 2021,
doi: 10.1016/j.cor.2020.105197.

\bibitem{Pereira2021}
R. H. M. Pereira, M. Saraiva, D. Herszenhut, C. K. V. Braga y M. W. Conway, ``r5r: Rapid realistic
routing on multimodal transport networks with R5 in R,'' \textit{Findings}, 2021,
doi: 10.32866/001c.21262.

\bibitem{Kuhn1955}
H. W. Kuhn, ``The Hungarian method for the assignment problem,'' \textit{Naval Research Logistics
Quarterly}, vol.~2, no.~1--2, pp.~83--97, 1955, doi: 10.1002/nav.3800020109.

\bibitem{Boeing2022}
G. Boeing, ``Street network models and indicators for every urban area in the world,''
\textit{Geographical Analysis}, vol.~54, no.~3, pp.~519--535, 2022, doi: 10.1111/gean.12281.

\bibitem{Geisberger2008}
R. Geisberger, P. Sanders, D. Schultes y D. Delling, ``Contraction hierarchies: Faster and simpler
hierarchical routing in road networks,'' en \textit{Experimental Algorithms (WEA 2008)}, Lecture
Notes in Computer Science, vol.~5038, Berlín: Springer, 2008, pp.~319--333,
doi: 10.1007/978-3-540-68552-4\_24.

\bibitem{LuxenVetter2011}
D. Luxen y C. Vetter, ``Real-time routing with OpenStreetMap data,'' en \textit{Proc. 19th ACM
SIGSPATIAL Int. Conf. Advances in Geographic Information Systems}, Chicago, IL, EE.~UU., 2011,
pp.~513--516, doi: 10.1145/2093973.2094062.

\bibitem{Wu2024}
L. Wu \textit{et al.}, ``LaDe: The first comprehensive last-mile express dataset from industry,''
en \textit{Proc. 30th ACM SIGKDD Conf. Knowledge Discovery and Data Mining}, 2024,
pp.~5991--6002, doi: 10.1145/3637528.3671548.

\bibitem{Wilcoxon1945}
F. Wilcoxon, ``Individual comparisons by ranking methods,'' \textit{Biometrics Bulletin}, vol.~1,
no.~6, pp.~80--83, 1945, doi: 10.2307/3001968.

\bibitem{Friedman1937}
M. Friedman, ``The use of ranks to avoid the assumption of normality implicit in the analysis of
variance,'' \textit{Journal of the American Statistical Association}, vol.~32, no.~200,
pp.~675--701, 1937, doi: 10.1080/01621459.1937.10503522.

\bibitem{MunozVillamizar2021}
A. Muñoz-Villamizar, E. L. Solano-Charris, L. Reyes-Rubiano y J. Faulin, ``Measuring disruptions in
last-mile delivery operations,'' \textit{Logistics}, vol.~5, no.~1, art.~17, 2021,
doi: 10.3390/logistics5010017.

\end{thebibliography}

\end{document}
"
\documentclass[10pt]{article}

% ── Revisión doble ciega ────────────────────────────────────────────
\newif\ifblind
\ifdefined\conautores \blindfalse \else \blindtrue \fi

% ── Página (plantilla RIIM: carta; márgenes 2.0 / 2.0 / 3.5 / 2.25 cm;
%    encabezado y pie a 1.25 cm del borde) ──────────────────────────
\usepackage[letterpaper, left=2cm, right=2cm, top=3.5cm, bottom=2.25cm,
            headheight=1.45cm, headsep=0.8cm, footskip=0.62cm]{geometry}

% ── Tipografía ──────────────────────────────────────────────────────
\usepackage{fontspec}
\setmainfont{Montserrat}[
  Path=fonts/, Extension=.ttf,
  UprightFont=*-Regular, ItalicFont=*-Italic,
  BoldFont=*-SemiBold, BoldItalicFont=*-SemiBoldItalic]
\newfontfamily\semibold{Montserrat}[
  Path=fonts/, Extension=.ttf,
  UprightFont=*-SemiBold, ItalicFont=*-SemiBoldItalic, Ligatures=TeX]
\newfontfamily\arialfont{Arial}[Ligatures=TeX]
\usepackage[spanish, mexico, es-tabla, es-nodecimaldot, es-noshorthands]{babel}
\spanishplainpercent
\usepackage{amsmath, amssymb}
\usepackage{xcolor}
\definecolor{riimazul}{HTML}{003D7C}
\definecolor{cfcfs}{HTML}{1E88E5}
\definecolor{crh}{HTML}{F59E0B}

\usepackage{graphicx}
\usepackage{booktabs}
\usepackage{tabularx}
\usepackage{float}
\usepackage{flafter}
% Permite que tablas y figuras compartan página con el texto en lugar de ir a páginas solo de flotantes.
\renewcommand{\topfraction}{0.85}
\renewcommand{\bottomfraction}{0.7}
\renewcommand{\textfraction}{0.1}
\renewcommand{\floatpagefraction}{0.8}
\setcounter{topnumber}{3}
\setcounter{bottomnumber}{2}
\setcounter{totalnumber}{4}
\setlength{\textfloatsep}{12pt plus 2pt minus 2pt}
\setlength{\floatsep}{10pt plus 2pt minus 2pt}
\setlength{\intextsep}{10pt plus 2pt minus 2pt}
\usepackage{tikz}
\usetikzlibrary{arrows.meta, positioning, fit, backgrounds}
\usepackage{fancyhdr}
\usepackage{caption}
\usepackage[hidelinks]{hyperref}

% Cuerpo: Montserrat 10, interlineado sencillo, justificado, sin sangría, 8 pt después.
\renewcommand{\normalsize}{\fontsize{10}{12.2}\selectfont}
\normalsize
\setlength{\parindent}{0pt}
\setlength{\parskip}{8pt plus 1pt}
\setlength{\abovedisplayskip}{6pt}
\setlength{\belowdisplayskip}{6pt}
\setlength{\abovedisplayshortskip}{4pt}
\setlength{\belowdisplayshortskip}{4pt}

% Títulos de sección sin numerar: Montserrat SemiBold 12 pt azul; subtítulos SemiBold itálica 11 pt.
\makeatletter
\setcounter{secnumdepth}{0}
\renewcommand{\section}{\@startsection{section}{1}{0pt}{8pt}{0.1pt}
  {\semibold\fontsize{12}{14.6}\selectfont\color{riimazul}\raggedright}}
\renewcommand{\subsection}{\@startsection{subsection}{2}{0pt}{8pt}{0.1pt}
  {\semibold\itshape\fontsize{11}{13.4}\selectfont\addfontfeatures{LetterSpace=6.8}\raggedright}}
\makeatother

% Pies de figura y de tabla: 9 pt, centrados, "Figura N." / "Tabla N.".
\captionsetup{font={footnotesize}, labelfont={footnotesize}, labelsep=period,
  justification=centering, skip=4pt}
\captionsetup[table]{position=top}
\renewcommand{\footnotesize}{\fontsize{9}{11}\selectfont}

% Encabezado: logotipo RIIM (imagen de la plantilla recortada al 59.3 % izquierdo, como en Word) a la izquierda y "RIIM, Vol 3, Num 2, 2026" (Arial 8) a la derecha.
% Pie: número de página centrado; ISSN a la izquierda y URL a la derecha (9 pt).
\pagestyle{fancy}
\fancyhf{}
\renewcommand{\headrulewidth}{0pt}
\renewcommand{\footrulewidth}{0pt}
\fancyhead[C]{\makebox[\textwidth]{%
  \raisebox{-0.5\height}{\includegraphics[trim=0 0 335.53bp 0, clip, width=7.37cm, height=1.40cm]{img/riim_header.png}}%
  \hfill\raisebox{-0.5\height}{\arialfont\fontsize{8}{9.6}\selectfont RIIM, Vol 3, Num 2, 2026}}}
\fancyfoot[C]{\parbox[t]{\textwidth}{\centering\fontsize{9}{11}\selectfont\thepage\\[1pt]
  ISSN 3122-4806\hfill http://orizaba.tecnm.mx/revistaojs}}

% Listas compactas.
\makeatletter
\renewcommand{\@listI}{\leftmargin=1.2em \topsep=0pt \partopsep=0pt \itemsep=2pt \parsep=0pt}
\let\@listi\@listI
\makeatother

% Referencias IEEE (Arial 10, sangría francesa de 1.27 cm).
\makeatletter
\renewenvironment{thebibliography}[1]{%
  \section{Referencias}
  \list{\@biblabel{\@arabic\c@enumiv}}{%
    \arialfont\fontsize{10}{12}\selectfont
    \settowidth\labelwidth{\@biblabel{#1}}%
    \leftmargin=1.27cm \labelwidth=1.27cm \labelsep=0pt \itemindent=0pt
    \itemsep=0pt \parsep=0pt \topsep=0pt
    \def\makelabel##1{##1\hfil}%
    \usecounter{enumiv}\let\p@enumiv\@empty
    \renewcommand\theenumiv{\@arabic\c@enumiv}}%
  \sloppy\clubpenalty4000\widowpenalty4000\sfcode`\.\@m}
  {\def\@noitemerr{\@latex@warning{Empty `thebibliography' environment}}\endlist}
\makeatother

\hypersetup{pdftitle={Logística de última milla en redes viales reales: MDRP reproducible en La Paz, B.C.S.},
  pdfauthor={}, pdfsubject={RIIM}, pdfkeywords={enrutamiento dinámico, horizonte rodante, OpenStreetMap, simulación}}
\ifblind\else\hypersetup{pdfauthor={Alan José Marcos Sánchez Martínez}}\fi

% Notación.
\newcommand{\CtD}{\text{CtD}}
\newcommand{\Cdisp}{\mathcal{C}_{disp}(t)}

% Diagrama de caja: \caja{x}{mín}{Q1}{mediana}{Q3}{máx}{color}
\newcommand{\caja}[7]{%
  \draw (#1,#2) -- (#1,#3) (#1,#5) -- (#1,#6);
  \draw (#1-0.14,#2) -- (#1+0.14,#2) (#1-0.14,#6) -- (#1+0.14,#6);
  \filldraw[fill=#7!75, draw=black!70] (#1-0.3,#3) rectangle (#1+0.3,#5);
  \draw[line width=1pt] (#1-0.3,#4) -- (#1+0.3,#4);}

\begin{document}

% ── Título, autores y área ──────────────────────────────────────────
\begin{center}
  \setlength{\parskip}{0pt}
  \vspace*{6pt}
  {\semibold\fontsize{16}{19.5}\selectfont\addfontfeatures{LetterSpace=-3.1}%
   LOGÍSTICA DE ÚLTIMA MILLA EN REDES VIALES REALES: MDRP REPRODUCIBLE EN LA PAZ, B.C.S.\par}
  \vspace{18pt}
  \ifblind\else
  {\itshape\fontsize{9}{11}\selectfont
   Sánchez Martínez, Alan José Marcos\textsuperscript{1*}\par
   \textsuperscript{1}Tecnológico Nacional de México / Instituto Tecnológico de La Paz, m16310322@lapaz.tecnm.mx\par}
  \fi
  {\fontsize{9}{11}\selectfont Área de participación: Sistemas Computacionales\par}
\end{center}

\section{Resumen}
Las plataformas de entrega de comida deciden sin conocer los pedidos futuros, si asignan cada orden
al repartidor libre más cercano o si esperan unos minutos para agrupar órdenes y coordinar la
asignación. Este trabajo evalúa en qué condiciones conviene coordinar en una ciudad media mexicana,
La Paz, B.C.S., y cuanto cambia la evaluación al usar la red vial real. Se construyó un simulador de
código abierto que obtiene los tiempos de viaje de OSRM sobre OpenStreetMap y compara una política
reactiva (FCFS) con una heurística de horizonte rodante (RH) en cuatro niveles de saturación de la
flota, con diez instancias sintéticas por nivel. Con una flota holgada, FCFS entrega 1.4~min antes;
con 26 o más pedidos por repartidor en la jornada, RH reduce el tiempo de entrega en todas las
instancias, hasta en 67\,\%. Las distancias euclidianas subestiman los tiempos de entrega y el
beneficio de coordinar.

{\semibold\fontsize{12}{14.6}\selectfont\color{riimazul}Palabras clave}: \textit{enrutamiento dinámico, horizonte rodante, OpenStreetMap, simulación}

\section{Abstract}
\begin{otherlanguage}{english}\frenchspacing
Food delivery platforms must decide, without knowing future orders, whether to assign each order to
the nearest idle courier or to wait a few minutes, batch orders and coordinate the assignment. This
work evaluates under which conditions coordination pays off in a mid-sized Mexican city, La Paz,
B.C.S., and how much the evaluation changes when the real road network is used. An open-source
simulator was built that obtains travel times from OSRM over OpenStreetMap and compares a reactive
policy (FCFS) with a rolling-horizon heuristic (RH) at four fleet saturation levels, with ten
synthetic instances per level. With a slack fleet, FCFS delivers 1.4~min earlier; with 26 or more
orders per courier over the day, RH reduces delivery time in every instance, by up to 67\,\%.
Euclidean distances underestimate both delivery times and the benefit of coordination.

{\semibold\fontsize{12}{14.6}\selectfont\color{riimazul}Key words}: \textit{dynamic routing, rolling horizon, OpenStreetMap, simulation}
\end{otherlanguage}

\section{Introducción}
La última milla, el tramo que va del punto de recolección al domicilio del cliente, es uno de los
eslabones mas costosos y difíciles de la logística urbana. En zonas residenciales concentra grandes
volúmenes de entregas con horarios estrechos y condiciones que cambian durante el día
\cite{Jazemi2023}. La entrega de comida a domicilio agrava estas exigencias: los pedidos llegan sin
aviso, la comida pierde calidad con cada minuto de espera y el cliente espera recibirla en menos de
una hora.

En este contexto, una plataforma debe decidir de forma continua qué repartidor atiende cada orden y
en qué momento. Este problema, formalizado como \textit{Meal Delivery Routing Problem} (MDRP) por
Reyes \textit{et al.} \cite{Reyes2018}, admite dos estrategias básicas. La primera es reactiva: cada
orden se asigna en cuanto aparece al repartidor libre más cercano (\textit{First-Come, First-Served},
FCFS). La segunda es coordinada: el sistema espera unos minutos, agrupa órdenes del mismo restaurante
y decide de manera conjunta qué repartidor atiende cada agrupamiento, como en la heurística de
horizonte rodante (\textit{Rolling Horizon}, RH) propuesta en \cite{Reyes2018}. Coordinar tiene un
costo, porque algunas órdenes esperan a que se forme el agrupamiento o a que se entreguen otras antes
que ellas, y un beneficio, porque cada viaje atiende más pedidos. El balance entre ambos depende de
cuántos repartidores hay frente a la demanda.

La evidencia disponible sobre este balance proviene de entornos muy distintos a una ciudad media. RH
se evaluó con instancias que reproducen áreas metropolitanas de Estados Unidos y con tiempos de viaje
calculados como la distancia euclidiana multiplicada por un factor constante \cite{Reyes2018}. En
instancias pequeñas resueltas a optimalidad, su \textit{Click-to-Door} promedio queda apenas un
4\,\% por encima del óptimo \cite{Yildiz2019}, pero no se ha reportado cómo cambia su ventaja frente
a una política reactiva conforme varía el tamaño de la flota. Otras formulaciones del problema
incorporan la incertidumbre en los tiempos de preparación \cite{Ulmer2021} o la estimación del tiempo
de llegada, cuya desviación respecto a lo prometido pesa en la satisfacción del cliente tanto como la
rapidez misma \cite{Hildebrandt2022}.

La forma de medir los tiempos de viaje añade una segunda fuente de error. Las distancias por la red
vial suelen ser entre 1.2 y 1.4 veces mayores que en línea recta, y en problemas con ventanas de
tiempo estrechas, como la entrega de comida, una ruta factible con distancias euclidianas puede
resultar infactible sobre la red real \cite{Boyaci2021}. Una evaluación que ignore la red puede
equivocarse no solo en los tiempos de entrega, sino también en la carga efectiva de la flota y, por
lo tanto, en el beneficio de coordinar. Los servicios comerciales de rutas resuelven la precisión,
pero su costo de licencia limita su uso en investigación \cite{Pereira2021} y su funcionamiento
cerrado dificulta reproducir los experimentos.

La Paz, B.C.S., ilustra esta brecha: es una ciudad media sin datos públicos de operaciones de entrega
de comida, cuya red vial, condicionada por la costa y por pocas avenidas principales, obliga a
recorridos más largos que la línea recta. De ahí surge la pregunta que guía este trabajo:
\textit{en una ciudad media, ¿en qué condiciones de saturación de la flota conviene coordinar los
despachos en lugar de asignar cada pedido al llegar, y cuánto cambia esa evaluación al usar la red
vial real en lugar de distancias euclidianas?} La pregunta se divide en tres: cómo varía la ventaja
de RH con la saturación y a partir de qué nivel cambia la política más conveniente (PI1); cuánto más largos son los
trayectos por la red y cuánto cambia la comparación al ignorarla (PI2); y si la evaluación puede
hacerse de forma reproducible, con recursos modestos y solo con herramientas de código abierto (PI3).

Las contribuciones corresponden a esas tres preguntas: (i)~una comparación pareada de FCFS y RH en
cuatro niveles de saturación, con prueba estadística por instancia, que ubica el punto a partir del
cual coordinar compensa; (ii)~la medición de la circuidad de la red urbana de La Paz y dos escenarios
de control con distancias euclidianas que separan el efecto de la red del efecto de la política; y
(iii)~un entorno de simulación reproducible, construido sobre OpenStreetMap y OSRM, que ejecuta un
día de operación en alrededor de un minuto en una computadora portátil. El estudio no propone un
algoritmo nuevo ni compara RH con políticas del estado del arte: su objetivo es caracterizar cuándo
conviene coordinar en un entorno realista.

\section{Metodología}

\subsection{Formulación del problema y métricas}
Se adopta la formulación del MDRP de \cite{Reyes2018}. Un conjunto $\mathcal{R}$ de restaurantes con
ubicación fija recibe órdenes $o \in \mathcal{O}$, cada una con restaurante de origen $r_o$, instante
de colocación $a_o$, instante en que la comida está lista $e_o \geq a_o$ y punto de entrega $\ell_o$.
La orden solo se conoce a partir de $a_o$. Una flota $\mathcal{C}$ de repartidores, con turno
$[e_c, l_c]$, recoge y entrega las órdenes. Órdenes del mismo restaurante pueden formar un
agrupamiento $B$ que se recoge en una sola visita y se entrega en una ruta de varias paradas. Los
tiempos de servicio son $s_r$ en el restaurante, una vez por agrupamiento, y $s_o$ con cada cliente.

La métrica principal es el \textit{Click-to-Door} de cada orden,
$\CtD_o = t_{delivery,o} - a_o$. A partir de ella se reportan su media y su percentil 90, y el
cumplimiento del nivel de servicio (SLA, \textit{Service Level Agreement}), entendido como el
porcentaje de órdenes entregadas dentro del objetivo $\tau$:
\begin{equation}
  \text{SLA} = \frac{|\{o \in \mathcal{O}_{del} : \CtD_o \leq \tau\}|}{|\mathcal{O}_{del}|} \times 100,
  \label{eq:sla}
\end{equation}
donde $\mathcal{O}_{del}$ es el conjunto de órdenes entregadas. Como ambas métricas se calculan
sobre órdenes entregadas, se reporta también el porcentaje de órdenes sin entregar al cierre. Se
complementan con el tiempo que la comida lista espera a ser recogida (RtP, \textit{Ready-to-Pickup}),
$\text{RtP}_o = t_{pickup,o} - e_o$, la
distancia recorrida por pedido, la utilización de la flota (fracción del turno en conducción o en
servicio), el tamaño de agrupamiento y el costo por pedido, con un pago de $p_1$ por orden entregada
y un mínimo garantizado de $p_2$ por hora de turno \cite{Reyes2018}.

\subsection{Políticas de despacho}
El simulador avanza en epocas de $f = 5$~min para ambas políticas: en cada época, FCFS asigna una por
una las órdenes nuevas y RH resuelve su optimización. Las dos calculan los instantes de recolección y
entrega con las mismas reglas. Un agrupamiento $B$ solo puede recogerse cuando ya está lista la
última de sus órdenes, es decir, a partir del instante $e_B = \max_{o\in B} e_o$. El repartidor $c$
llega al restaurante en el instante $t_{arr,r}$ y, como en \cite{Reyes2018}, cada tiempo de servicio
se reparte mitad antes y mitad después del momento de recoger o de entregar. Si la comida ya está
lista, la recolección ocurre al cumplirse la primera mitad del servicio en el restaurante; si no, el
repartidor espera hasta que esté lista. Por eso el instante de recolección es
$t_{pickup}(c,B) = \max\bigl(e_B,\; t_{arr,r} + s_r/2\bigr)$. Después, el repartidor sale del
restaurante al terminar la segunda mitad del servicio, recorre el tramo hasta el primer cliente y le
entrega al cumplirse la primera mitad del servicio con él; cada entrega siguiente se calcula igual a
partir de la anterior.

\textbf{FCFS (línea base reactiva).} En cada época, las órdenes reveladas sin repartidor se procesan
por orden de colocación. Cada una se asigna al repartidor libre más cercano en línea recta al
restaurante (esta distancia solo sirve para escoger al repartidor); la ruta repartidor $\to$ restaurante
$\to$ cliente y sus tiempos se obtienen de OSRM, y el repartidor queda comprometido de inmediato. Cada
viaje lleva una sola orden.

\textbf{RH (política coordinada).} Sigue el esquema de \cite{Reyes2018} y en cada época $t$ ejecuta
cinco pasos. Primero, delimita las órdenes que entran a la optimización: de las órdenes ya reveladas y
aún no entregadas, $\mathcal{O}_t$, solo considera las que estarán listas dentro del horizonte de
asignación $\Delta_u$, es decir, $U_t = \{o \in \mathcal{O}_t : e_o \leq t + \Delta_u\}$; planear
órdenes cuya comida tardará mucho en estar lista sería poco útil en un entorno que cambia cada pocos
minutos.

Segundo, estima cuántas órdenes conviene juntar en cada viaje. El tamaño objetivo de agrupamiento
$Z_t$ compara la demanda próxima con la oferta próxima: el numerador cuenta las órdenes que estarán
listas en los siguientes $\Delta_1$ minutos, y el denominador, los repartidores en turno, libres u
ocupados, que quedarán disponibles en los siguientes $\Delta_2$ minutos, donde $e^{avail}_c$ es el
instante en que el repartidor $c$ termina su ruta actual:
\begin{equation}
  Z_t = \left\lceil \frac{|\{o \in \mathcal{O}_t : e_o \leq t + \Delta_1\}|}
                        {|\{c \in \mathcal{C}_{turno}(t) : e^{avail}_c \leq t + \Delta_2\}|} \right\rceil.
  \label{eq:zt}
\end{equation}
Si hay tantos repartidores como ordenes, $Z_t = 1$ y cada orden tiende a viajar sola; si la flota se
satura, $Z_t$ crece y se forman agrupamientos más grandes. $Z_t$ se acota en 5, y ningún agrupamiento
puede superar $B_{max} = 4$ órdenes.

Tercero, forma los agrupamientos de cada restaurante $r$. Sus órdenes $U_{t,r}$ se reparten en
\[
  m_r = \min\bigl(\lceil |U_{t,r}|/Z_t \rceil,\; |\Cdisp|\bigr)
\]
agrupamientos, de modo que cada uno reciba alrededor de $Z_t$ órdenes sin crear más agrupamientos que
repartidores disponibles. Las órdenes se toman por \textit{ready time}, y cada una se inserta en el
agrupamiento y la posición de la ruta que menos incrementan el costo
\begin{equation}
  \text{costo}(B, o, p) = T_{ruta}(B \oplus_p o) + \beta\, T_{servicio}(B \oplus_p o),
  \label{eq:costo}
\end{equation}
donde $B \oplus_p o$ es el agrupamiento con la orden $o$ insertada en la posición $p$, $T_{ruta}$ es
la duración de la ruta según OSRM y $T_{servicio}$ el tiempo de servicio acumulado. El peso $\beta$, que en
\cite{Reyes2018} representa la pérdida de frescura y se calibra en $[0, 1]$, se fija en 0.1 para que la
duración de la ruta domine el costo. Después, dos iteraciones de extracción y reinserción mueven
cada orden al agrupamiento donde su costo es menor.

Cuarto, asigna agrupamientos a repartidores. Los agrupamientos se clasifican en tres grupos de
prioridad: urgentes, cuando ningún repartidor puede entregar su orden más antigua dentro del objetivo
$\tau$; retrasados, cuando el objetivo aún es alcanzable pero ningún repartidor llega al restaurante
antes de que la comida esté lista; y normales. Los grupos se resuelven en ese orden, para que los
urgentes obtengan primero a los repartidores más cercanos. En cada grupo se resuelve un problema de
asignación lineal, en el que cada repartidor recibe a lo más un agrupamiento, con el peso
\begin{equation}
  w_{c,B} = \frac{|B|}{T_{delivery}(c,B)} - \theta\,\bigl(t_{pickup}(c,B) - e_B\bigr) - P_g.
  \label{eq:peso}
\end{equation}
El primer término mide la eficiencia del viaje en órdenes entregadas por hora, donde $T_{delivery}(c,B)$
es el tiempo desde que $c$ queda libre hasta la última entrega de $B$. El segundo penaliza, con
$\theta = 0.1$, los minutos que la comida espera lista en el restaurante, y $P_g$ resta 100 si el par
excede el máximo absoluto $\tau_{max}$ y 50 si la recolección se retrasa. La asignación de máximo peso
se obtiene de forma exacta con un algoritmo de tipo húngaro \cite{Kuhn1955}.

Quinto, el compromiso en dos etapas de \cite{Reyes2018} decide qué parte de la asignación se
ejecuta. Si el repartidor llega al restaurante y todas las órdenes estarán listas dentro de la época,
o si alguna orden lleva más de $\Delta_x$ minutos lista, se fija la ruta completa. En otro caso, el
repartidor solo se dirige al restaurante, como un preposicionamiento, y en la época siguiente el
agrupamiento puede crecer o cambiar de repartidor con las órdenes que se hayan revelado.

\subsection{Entorno de simulación reproducible}
La Figura~\ref{fig:arquitectura} resume el entorno. El extracto de OpenStreetMap de Baja California
Sur, modelo de datos viales abierto y de cobertura mundial \cite{Boeing2022}, se preprocesa con el
algoritmo \textit{Contraction Hierarchies} \cite{Geisberger2008}, que permite resolver cada consulta
de camino mínimo en milisegundos. El servidor OSRM \cite{LuxenVetter2011}, versión 5.27.1, atiende
las consultas de ruta del simulador por HTTP; su imagen oficial de Docker permite reproducir el
entorno en cualquier sistema operativo. El simulador, escrito en Python, avanza por épocas, guarda en
memoria las rutas ya consultadas y produce, para cada corrida, los registros por orden y por
repartidor de los que se derivan los indicadores.

\begin{figure}[htbp]
\centering
\begin{tikzpicture}[
  font=\footnotesize,
  caja/.style={draw=black!70, rounded corners=2pt, align=center, minimum height=1.05cm,
               text width=#1, fill=white},
  caja/.default=2.8cm,
  flecha/.style={-{Stealth[length=5pt]}, thick, black!70},
  capa/.style={draw=#1!70!black, rounded corners=4pt, fill=#1!8, inner xsep=7pt, inner ysep=7pt}]
  \node[caja=2.8cm] (osm)   at (0,0)      {Extracto\\OpenStreetMap};
  \node[caja=3.0cm] (ch)    at (3.55,0)    {Preprocesamiento\\\textit{Contraction}\\\textit{Hierarchies}};
  \node[caja=2.8cm] (osrm)  at (7.1,0)    {Servidor OSRM\\(API de rutas)};
  \node[caja=2.8cm] (gen)   at (0,-2.3)   {Generador de\\instancias (PPNH)};
  \node[caja=2.8cm] (sim)   at (3.55,-2.3) {Simulador por\\épocas ($f = 5$ min)};
  \node[caja=2.8cm] (pol)   at (7.1,-2.3) {Política\\FCFS o RH};
  \node[caja=2.8cm] (kpi)   at (10.9,-2.3) {Registros por orden\\y 21 indicadores};
  \node[caja=2.8cm] (est)   at (10.9,0)   {Análisis estadístico\\(Wilcoxon, Friedman)};
  \draw[flecha] (osm) -- (ch);
  \draw[flecha] (ch) -- (osrm);
  \draw[flecha] (gen) -- (sim);
  \draw[flecha, {Stealth[length=5pt]}-{Stealth[length=5pt]}] (sim) -- (pol);
  \draw[flecha, dashed, {Stealth[length=5pt]}-{Stealth[length=5pt]}] (osrm) --
     node[right, font=\scriptsize, align=left] {tiempos de\\viaje (HTTP)} (pol);
  \draw[flecha] (pol) -- (kpi);
  \draw[flecha] (kpi) -- (est);
  \begin{scope}[on background layer]
    \node[capa=blue, fit=(osm)(ch)(osrm)] {};
    \node[capa=green, fit=(gen)(sim)(pol)] {};
    \node[capa=orange, fit=(kpi)(est)] {};
  \end{scope}
\end{tikzpicture}
\caption{Arquitectura del entorno: enrutamiento sobre la red vial (azul), simulación (verde) y
análisis (naranja).}
\label{fig:arquitectura}
\end{figure}

\subsection{Generación de instancias}
Al no existir datos públicos de entregas de comida en La Paz, las instancias se generan de forma
sintética. Las llegadas siguen un proceso de Poisson no homogéneo con intensidad bimodal
\begin{equation}
  \lambda(t) = \lambda_0
  + A_m \exp\!\left(-\frac{(t-\mu_m)^2}{2\sigma^2}\right)
  + A_e \exp\!\left(-\frac{(t-\mu_e)^2}{2\sigma^2}\right),
  \label{eq:lambda}
\end{equation}
siguiendo el patrón bimodal observado en entregas de última milla \cite{Wu2024}. A falta de datos
locales, los parámetros son supuestos: ventana operativa de 08:00 a 19:00, picos en $\mu_m = $ 09:00
y $\mu_e = $ 17:00, $\lambda_0 = 0.6$ órdenes/min, $A_m = 2.4$, $A_e = 2.0$ y $\sigma = 45$~min. El perfil se escala para
obtener alrededor de mil órdenes por día (entre 988 y 1\,047 según la semilla). El tiempo de
preparación es $\max(5, \mathcal{N}(12, 3^2))$~min. Los 25 restaurantes y los clientes se ubican
por muestreo uniforme dentro de un polígono de 8.4~km$^2$ del centro urbano, y cada cliente se asocia
a un restaurante situado a entre 2 y 5~km en línea recta. Todos los repartidores comparten un turno
que inicia 15~min antes del primer pedido y termina una hora después del último \textit{ready time},
y parten del centroide de los restaurantes. La flota se dimensiona a partir de la razón de
saturación $\rho$, el número objetivo de pedidos por repartidor en la jornada:
\begin{equation}
  |\mathcal{C}| = \max\!\left(n_{\min},\; \left\lceil |\mathcal{O}| / \rho \right\rceil\right),
  \qquad n_{\min} = 10.
  \label{eq:flota}
\end{equation}
Fijar $\rho$, y no el número de repartidores, mantiene la misma carga relativa en todas las semillas
de un escenario aunque el número de órdenes varíe.

\subsection{Diseño experimental}
La Tabla~\ref{tab:parametros} reúne los parámetros, comunes a ambas políticas, y los cuatro
escenarios de saturación. Cada semilla (2025 a 2034) genera una instancia, y sobre ella se ejecutan
FCFS y RH con la misma flota; como el simulador no tiene componentes aleatorios propios, lo único que
cambia dentro de cada par es la política. Esto da 80 corridas sobre la red vial. Para evitar la
pseudorreplicación, cada métrica se promedia por semilla y la comparación se hace con la prueba de
rangos con signo de Wilcoxon pareada \cite{Wilcoxon1945} sobre las diez diferencias de cada
escenario, con intervalo de confianza $t$ al 95\,\% y la correlación biserial de rangos $r_{rb}$ como
tamaño de efecto. La prueba de Friedman \cite{Friedman1937} sobre la matriz de diferencias
$10 \times 4$ evalúa si el efecto cambia entre escenarios.

\begin{table}[htbp]
\centering
\footnotesize
\caption{Parámetros del simulador y escenarios de saturación.}
\label{tab:parametros}
\setlength{\tabcolsep}{4pt}
\begin{tabular}{@{}l c c @{\hspace{0.4cm}} l c c@{}}
\toprule
Parámetro & Símbolo & Valor & Parámetro & Símbolo & Valor \\
\midrule
Frecuencia de optimización & $f$ & 5 min & Penalización de frescura & $\beta$ & 0.1 \\
Horizonte de asignación & $\Delta_u$ & 20 min & Penalización por espera & $\theta$ & 0.1 \\
Anticipación (demanda, oferta) & $\Delta_1$, $\Delta_2$ & 20 min & Espera para forzar compromiso & $\Delta_x$ & 10 min \\
CtD objetivo / máximo & $\tau$, $\tau_{max}$ & 40/90 min & Tamaño máximo de agrupamiento & $B_{max}$ & 4 \\
Tiempos de servicio & $s_r$, $s_o$ & 4 min & Pago por orden / mínimo por hora & $p_1$, $p_2$ & \$10 / \$15 \\
\midrule
\multicolumn{6}{@{}l}{Escenarios: S1, $\rho = 18$ (55--59 repartidores); S2, $\rho = 26$ (38--41); S3, $\rho = 35$ (29--30); S4, $\rho = 52$ (19--21).} \\
\bottomrule
\end{tabular}
\end{table}

Para medir el efecto de la red vial se calculó la circuidad $\kappa = d_{OSRM}/d_{recta}$ sobre todos
los trayectos restaurante--cliente, y el experimento completo se repitió con dos controles en los que
el tiempo de viaje es la distancia en línea recta (Haversine) a velocidad constante: la velocidad
media que OSRM asigna a esos trayectos (651~m/min), que deja como única diferencia la geometría de
las rutas, y esa velocidad dividida entre $\kappa$ (472~m/min), que iguala la duración media de los
viajes. Cada control suma otras 80 corridas.

\section{Resultados y discusión}

\subsection{Calidad de servicio y cambio de la política más conveniente}
La Tabla~\ref{tab:calidad} muestra los indicadores de calidad de servicio. Con la flota holgada de S1,
FCFS entrega en promedio 1.4~min antes que RH: con repartidores de sobra, llevar cada pedido directo
del restaurante al cliente es más rápido que esperar a formar un agrupamiento, y en un agrupamiento
cada entrega espera a las anteriores. Aun así, el cumplimiento del SLA es prácticamente igual en ambas
políticas. A partir de S2 la relación se invierte: RH reduce el \textit{Click-to-Door} medio en
36.8\,\%, el percentil 90 en 46.0\,\% y eleva el cumplimiento del SLA en 31.5 puntos porcentuales. La
ventaja crece en S3 y S4, donde FCFS deja sin entregar al cierre el 11.3\,\% y el 39.9\,\% de los
pedidos, mientras que RH entrega prácticamente todos. Como el CtD y el SLA se calculan sobre órdenes
entregadas, la tabla favorece a FCFS en esos escenarios: medido sobre todas las órdenes generadas, su
cumplimiento cae a 6.4\,\% en S3 y 3.8\,\% en S4, frente a 57.6\,\% y 42.4\,\% bajo RH.

\begin{table}[htbp]
\centering
\footnotesize
\caption{Calidad de servicio, promedio de 10 semillas. $\Delta$ es el cambio de RH respecto a FCFS.}
\label{tab:calidad}
\setlength{\tabcolsep}{3pt}
\begin{tabular*}{\textwidth}{@{\extracolsep{\fill}}l rrr rrr rrr rr@{}}
\toprule
& \multicolumn{3}{c}{CtD medio (minutos)} & \multicolumn{3}{c}{CtD percentil 90 (minutos)} & \multicolumn{3}{c}{SLA (\%)} & \multicolumn{2}{c}{Sin entregar (\%)} \\
\cmidrule(lr){2-4} \cmidrule(lr){5-7} \cmidrule(lr){8-10} \cmidrule(l){11-12}
Esc. & FCFS & RH & $\Delta$ (\%) & FCFS & RH & $\Delta$ (\%) & FCFS & RH & $\Delta$ (pp) & FCFS & RH \\
\midrule
S1 & 28.93  & 30.31 & $+4.8$  & 43.33  & 43.80  & $+1.1$  & 84.6 & 83.4 & $-1.1$  & 0.0  & 0.0 \\
S2 & 56.07  & 35.42 & $-36.8$ & 97.91  & 52.87  & $-46.0$ & 38.1 & 69.6 & $+31.5$ & 0.1  & 0.0 \\
S3 & 108.49 & 45.36 & $-58.2$ & 169.54 & 79.10  & $-53.3$ & 7.3  & 57.6 & $+50.3$ & 11.3 & 0.0 \\
S4 & 224.88 & 73.63 & $-67.3$ & 321.10 & 186.18 & $-42.0$ & 6.2  & 42.5 & $+36.3$ & 39.9 & 0.2 \\
\bottomrule
\end{tabular*}
\end{table}

La Figura~\ref{fig:boxplot} muestra la distribución completa del \textit{Click-to-Door}. En S1 ambas
cajas son casi idénticas y quedan por debajo del objetivo de 40~min; desde S2, la distribución de
FCFS se alarga hacia valores de varias horas, mientras que la de RH se mantiene compacta y con la
mediana cerca del objetivo.

\begin{figure}[htbp]
\centering
\begin{tikzpicture}[y=0.0100cm, x=1cm, font=\footnotesize]
  \foreach \p/\esc in {0/S1 ($\rho=18$), 1/S2 ($\rho=26$), 2/S3 ($\rho=35$), 3/S4 ($\rho=52$)} {
    \begin{scope}[xshift=\p*4.15cm]
      \draw[black!60] (0,0) rectangle (3.6,390);
      \foreach \y in {0,100,200,300} \draw[black!12] (0,\y) -- (3.6,\y);
      \draw[red!70, dashed] (0,40) -- (3.6,40);
      \node[above] at (1.8,390) {\esc};
      \node[below] at (1.1,0) {FCFS};
      \node[below] at (2.5,0) {RH};
    \end{scope}}
  \foreach \y in {0,100,200,300} \node[left] at (0,\y) {\y};
  \node[rotate=90] at (-1.25,195) {\textit{Click-to-Door} (minutos)};
  % S1
  \caja{1.1}{10.9}{21.7}{26.4}{35.1}{55.1}{cfcfs}
  \caja{2.5}{11.4}{22.8}{28.5}{36.1}{56.1}{crh}
  % S2
  \begin{scope}[xshift=4.15cm]
  \caja{1.1}{11.9}{27.3}{54.3}{80.6}{121.7}{cfcfs}
  \caja{2.5}{11.4}{24.5}{32.5}{42.2}{68.7}{crh}
  \end{scope}
  % S3
  \begin{scope}[xshift=8.3cm]
  \caja{1.1}{11.9}{72.1}{109.1}{148.2}{199.9}{cfcfs}
  \caja{2.5}{11.4}{26.6}{36.5}{50.3}{85.8}{crh}
  \end{scope}
  % S4
  \begin{scope}[xshift=12.45cm]
  \caja{1.1}{14.9}{141.0}{266.0}{306.5}{367.5}{cfcfs}
  \caja{2.5}{12.3}{31.1}{43.8}{79.3}{151.5}{crh}
  \end{scope}
\end{tikzpicture}
\caption{Distribución del \textit{Click-to-Door} de las órdenes entregadas en las diez semillas
(cajas: cuartiles; bigotes: 1.5 veces el rango intercuartílico). La línea discontinua marca el
objetivo $\tau = 40$~min.}
\label{fig:boxplot}
\end{figure}

La Tabla~\ref{tab:wilcoxon} confirma, a nivel de instancia, el cambio de la política más conveniente. En S2, S3 y S4, RH
supera a FCFS en las diez semillas, tanto en CtD como en SLA, con $p = 0.0020$, el valor mínimo que
admite la prueba con $n = 10$, y $|r_{rb}| = 1$. En S1 la ventaja de FCFS en CtD es significativa
pero pequeña, y la diferencia en SLA no lo es. La prueba de Friedman confirma que la magnitud del
efecto cambia con la saturación ($\chi^2_F = 30.0$, $p < 0.001$). La Figura~\ref{fig:saturacion}
resume la respuesta a PI1: las curvas de CtD se cruzan entre $\rho = 18$ y $\rho = 26$, y por encima
de ese punto la ventaja de coordinar crece con la carga.

\begin{table}[htbp]
\centering
\footnotesize
\caption{Wilcoxon pareada por semilla ($n = 10$), $\bar{\Delta} = \bar{x}_{FCFS} - \bar{x}_{RH}$: favorecen a RH
$\bar{\Delta} > 0$ en CtD y $\bar{\Delta} < 0$ en SLA.}
\label{tab:wilcoxon}
\begin{tabular}{@{}ll rr l c r@{}}
\toprule
Esc. & Métrica & $\bar{x}_{FCFS}$ & $\bar{x}_{RH}$ & $\bar{\Delta}$ [IC 95\,\%] & $p$ & $r_{rb}$ \\
\midrule
S1 & CtD (minutos) & 28.93  & 30.31 & $-1.38$ [$-2.35$, $-0.41$]        & 0.0137 & $-0.85$ \\
S1 & SLA (\%)  & 84.6   & 83.4  & $+1.1$ [$-1.8$, $+4.1$]           & 0.2324 & $+0.45$ \\
S2 & CtD (minutos) & 56.07  & 35.42 & $+20.64$ [$+18.98$, $+22.31$]     & 0.0020 & $+1.00$ \\
S2 & SLA (\%)  & 38.1   & 69.6  & $-31.5$ [$-33.7$, $-29.2$]        & 0.0020 & $-1.00$ \\
S3 & CtD (minutos) & 108.49 & 45.36 & $+63.13$ [$+56.14$, $+70.13$]     & 0.0020 & $+1.00$ \\
S3 & SLA (\%)  & 7.3    & 57.6  & $-50.3$ [$-52.6$, $-48.0$]        & 0.0020 & $-1.00$ \\
S4 & CtD (minutos) & 224.88 & 73.63 & $+151.24$ [$+143.71$, $+158.78$]  & 0.0020 & $+1.00$ \\
S4 & SLA (\%)  & 6.2    & 42.5  & $-36.3$ [$-37.2$, $-35.3$]        & 0.0020 & $-1.00$ \\
\bottomrule
\end{tabular}
\end{table}

\begin{figure}[htbp]
\centering
\begin{tikzpicture}[font=\footnotesize]
  % Panel izquierdo: CtD medio
  \begin{scope}[x=0.16cm, y=0.018cm]
    \draw[black!60] (14,0) rectangle (56,240);
    \foreach \y in {0,50,100,150,200} {\draw[black!12] (14,\y) -- (56,\y); \node[left] at (14,\y) {\y};}
    \foreach \x in {18,26,35,52} {\draw[black!12] (\x,0) -- (\x,240); \node[below] at (\x,0) {\x};}
    \draw[red!70, dashed] (14,40) -- (56,40);
    \draw[cfcfs, very thick] (18,28.93) -- (26,56.07) -- (35,108.49) -- (52,224.88);
    \foreach \x/\y in {18/28.93, 26/56.07, 35/108.49, 52/224.88} \fill[cfcfs] (\x,\y) circle (2pt);
    \draw[crh, very thick] (18,30.31) -- (26,35.42) -- (35,45.36) -- (52,73.63);
    \foreach \x/\y in {18/30.31, 26/35.42, 35/45.36, 52/73.63} \fill[crh] (\x-0.45,\y-4.5) rectangle (\x+0.45,\y+4.5);
    \node[rotate=90] at (8.5,120) {CtD medio (minutos)};
    \node at (35,-30) {Pedidos por repartidor en la jornada ($\rho$)};
  \end{scope}
  % Panel derecho: SLA
  \begin{scope}[xshift=8.6cm, x=0.16cm, y=0.043cm]
    \draw[black!60] (14,0) rectangle (56,100);
    \foreach \y in {0,20,40,60,80,100} {\draw[black!12] (14,\y) -- (56,\y); \node[left] at (14,\y) {\y};}
    \foreach \x in {18,26,35,52} {\draw[black!12] (\x,0) -- (\x,100); \node[below] at (\x,0) {\x};}
    \draw[cfcfs, very thick] (18,84.6) -- (26,38.1) -- (35,7.3) -- (52,6.2);
    \foreach \x/\y in {18/84.6, 26/38.1, 35/7.3, 52/6.2} \fill[cfcfs] (\x,\y) circle (2pt);
    \draw[crh, very thick] (18,83.4) -- (26,69.6) -- (35,57.6) -- (52,42.5);
    \foreach \x/\y in {18/83.4, 26/69.6, 35/57.6, 52/42.5} \fill[crh] (\x-0.45,\y-1.9) rectangle (\x+0.45,\y+1.9);
    \node[rotate=90] at (8.5,50) {Cumplimiento del SLA (\%)};
    \node at (35,-12.5) {Pedidos por repartidor en la jornada ($\rho$)};
    % Leyenda
    \draw[cfcfs, very thick] (34,93) -- (38,93); \fill[cfcfs] (36,93) circle (2pt); \node[right] at (38,93) {FCFS};
    \draw[crh, very thick] (46,93) -- (50,93); \fill[crh] (47.55,91.1) rectangle (48.45,94.9); \node[right] at (50,93) {RH};
  \end{scope}
\end{tikzpicture}
\caption{\textit{Click-to-Door} medio y cumplimiento del SLA según la saturación de la flota. La línea
discontinua marca el objetivo $\tau = 40$~min.}
\label{fig:saturacion}
\end{figure}

\subsection{Eficiencia operativa y agrupamiento}
La Tabla~\ref{tab:eficiencia} explica el origen de la ventaja de RH. Aun en S1, RH reduce la espera
de la comida lista en el restaurante (RtP), porque en los picos de demanda la flota se ocupa por completo y
FCFS, con un pedido por viaje, acumula órdenes sin recoger, mientras que RH recoge varias en una sola
visita. Al agrupar, RH recorre entre 13.5\,\% y 49.4\,\% menos distancia por pedido y mantiene una
utilización de la flota menor que FCFS en S1--S3: con la misma flota, entrega lo mismo con menos
tiempo de conducción y, en S4, convierte ese margen en pedidos entregados. El tamaño de agrupamiento
crece con la saturación; en S1 una cuarta parte de los pedidos (25.4\,\%) viaja sola y en S4 el
80.9\,\% viaja en agrupamientos de tres o cuatro. El costo por pedido solo difiere en S1, donde muchos
repartidores reciben el pago mínimo y RH concentra las entregas en menos de ellos; desde S2 ambos
convergen al pago por orden.

\begin{table}[htbp]
\centering
\footnotesize
\caption{Eficiencia operativa y tamaño del agrupamiento en que viaja cada pedido (promedio de 10 semillas).}
\label{tab:eficiencia}
\setlength{\tabcolsep}{3pt}
\begin{tabular*}{\textwidth}{@{\extracolsep{\fill}}l rr rr rr r rr@{}}
\toprule
& \multicolumn{2}{c}{RtP medio (minutos)} & \multicolumn{2}{c}{Distancia/pedido (km)} & \multicolumn{2}{c}{Utilización (\%)} & Agrupamiento & \multicolumn{2}{c}{Costo/pedido (\$)} \\
\cmidrule(lr){2-3} \cmidrule(lr){4-5} \cmidrule(lr){6-7} \cmidrule(lr){8-8} \cmidrule(l){9-10}
Esc. & FCFS & RH & FCFS & RH & FCFS & RH & RH & FCFS & RH \\
\midrule
S1 & 7.78   & 4.94  & 4.71 & 4.08 & 37.7 & 31.2 & 2.54 & 10.80 & 11.59 \\
S2 & 34.92  & 9.33  & 4.88 & 3.60 & 55.2 & 41.4 & 2.80 & 10.00 & 10.03 \\
S3 & 87.35  & 18.62 & 5.07 & 3.17 & 66.8 & 51.4 & 3.03 & 10.00 & 10.00 \\
S4 & 203.75 & 46.20 & 5.20 & 2.63 & 67.3 & 67.8 & 3.30 & 10.00 & 10.00 \\
\bottomrule
\end{tabular*}
\end{table}

\subsection{Efecto de la red vial}
Sobre los 10\,184 trayectos restaurante--cliente de las diez instancias, la distancia por la red es
en promedio 3.30~km y en línea recta 2.39~km: los trayectos reales son 38\,\% más largos (circuidad
$\kappa = 1.38$), dentro del rango de 1.2 a 1.4 que reporta la literatura \cite{Boyaci2021}.

Para medir cuánto cambia la comparación, se repitió el experimento con tiempos de viaje en línea
recta. A la velocidad media de la red (651~m/min, primer control), ambas políticas parecen más
rápidas: en S2, el CtD de FCFS baja de 56.1 a 41.8~min y el de RH, de 35.4 a 30.3~min. Como FCFS
mejora más, la ventaja de RH se reduce en los cuatro escenarios ($p = 0.002$); en S3, por ejemplo,
pasa de 63.1 a 37.1~min. El segundo control baja la velocidad hasta que los viajes en línea recta
duran en promedio lo mismo que por la red ($651/1.38 \approx 472$~m/min); con él desaparece entre el
62 y el 84\,\% de esa reducción. La mayor parte del efecto se debe, por tanto, a que los trayectos
reales son más largos y cargan más a la flota, y el resto, a la forma de la red ($p \leq 0.004$).
Ignorar la red hace parecer más rápidas a ambas políticas y subestima el beneficio de coordinar, lo
que responde a la pregunta PI2.

\subsection{Reproducibilidad y costo de cómputo}
Toda la cadena (OpenStreetMap, OSRM y Python) es de código abierto, las instancias se regeneran a
partir de semillas fijas y cada corrida pasa verificaciones automáticas de integridad: ninguna orden
se recoge antes de estar lista ni se entrega antes de recogerse, no hay entregas duplicadas y ninguna
ruta recurrió a la estimación de respaldo en línea recta (de haber fallado OSRM). Con 1\,038 órdenes, simular un día toma
entre 1.8 y 2.7~segundos con FCFS y entre 57 y 80~segundos con RH, incluidas las consultas a OSRM, con menos de
225~MB de memoria por proceso, en una computadora portátil con procesador Apple M1 Pro y 16~GB de
RAM. Esto responde a PI3: la evaluación es reproducible y no requiere infraestructura especializada.

\subsection{Limitaciones}
El reloj del simulador avanza en épocas de 5~min. En cada época primero se asignan los pedidos y
después se registran las rutas terminadas, de modo que un repartidor que termina a mitad de una época
queda libre hasta la siguiente y espera entre 5 y 10~min su nuevo viaje. Como esta espera ocurre una
vez por viaje, perjudica más a FCFS, que hace un viaje por pedido. Para comprobar que no determina el
resultado, se repitieron las 80 corridas sobre la red vial con el orden inverso (primero se registran
las rutas terminadas y luego se asigna): FCFS sigue ganando en S1 y RH en S2--S4 en todas las
semillas, y el cruce permanece entre S1 y S2. La reducción del CtD baja a 14.6\,\% en S2, 41.4\,\% en
S3 y 63.0\,\% en S4, por lo que las magnitudes dependen de este detalle de implementación, pero las
conclusiones se mantienen. Además, la demanda es
sintética: sigue el patrón bimodal de la recolección de paquetes en LaDe con parámetros supuestos, no
calibrados con la demanda local de comida, y el tipo de ciudad influye en el
comportamiento de la última milla \cite{MunozVillamizar2021}. Los tiempos de preparación y de
servicio son homogéneos, la red se considera sin congestión y las ubicaciones se muestrean de forma
uniforme dentro del dominio urbano.

\section{Trabajo a futuro}
Quedan pendientes seis líneas de trabajo: calibrar la demanda, los tiempos de preparación y la
ubicación de restaurantes y clientes con datos locales de una plataforma de entregas o de encuestas;
explorar valores de $\rho$ entre 18 y 26 para ubicar con precisión el punto de cruce y evaluar
políticas híbridas que activen la coordinación solo cuando la carga lo justifique; incorporar
velocidades dependientes de la hora para representar la congestión; ampliar el dominio a toda la
mancha urbana, incluidos los trayectos que rodean la ensenada de La Paz; comparar RH con políticas
anticipatorias basadas en aprendizaje automático sobre el mismo entorno; y evaluar la escalabilidad
del simulador con instancias de mayor tamaño.

\section{Conclusiones}
En las instancias evaluadas del centro urbano de La Paz, coordinar los despachos no compensa cuando
la flota es holgada, pero sí a partir de una saturación moderada. Con una flota dimensionada para 18
pedidos por repartidor en la jornada, FCFS entrega 1.4~min antes que RH y el cumplimiento del SLA es
equivalente. Con 26 o más pedidos por repartidor, RH reduce el tiempo de entrega y mejora el
cumplimiento del servicio en todas las semillas, y cuando la flota es escasa entrega prácticamente
todos los pedidos que FCFS deja pendientes. Para un operador, esto sugiere que la decisión de
coordinar debe tomarse en función de la carga esperada por repartidor y no de forma fija.

La red vial importa para esta decisión. Los trayectos reales son en promedio 38\,\% más largos que
en línea recta, y un modelo euclidiano subestima tanto los tiempos de entrega como el beneficio de
coordinar. Finalmente, toda la evaluación puede reproducirse con herramientas abiertas en una
computadora portátil, lo que la hace accesible para operadores y desarrolladores de ciudades medias
que no cuentan con servicios comerciales de rutas.

\begin{thebibliography}{99}

\bibitem{Jazemi2023}
R. Jazemi, E. Alidadiani, K. Ahn y J. Jang, ``A review of literature on vehicle routing problems of
last-mile delivery in urban areas,'' \textit{Applied Sciences}, vol.~13, no.~24, art.~13015, 2023,
doi: 10.3390/app132413015.

\bibitem{Reyes2018}
D. Reyes, A. Erera, M. Savelsbergh, S. Sahasrabudhe y R. O'Neil, ``The meal delivery routing
problem,'' preprint, \textit{Optimization Online}, 2018. [En línea]. Disponible:
https://optimization-online.org/2018/04/6571/

\bibitem{Yildiz2019}
B. Yıldız y M. Savelsbergh, ``Provably high-quality solutions for the meal delivery routing
problem,'' \textit{Transportation Science}, vol.~53, no.~5, pp.~1372--1388, 2019,
doi: 10.1287/trsc.2018.0887.

\bibitem{Ulmer2021}
M. W. Ulmer, B. W. Thomas, A. M. Campbell y N. Woyak, ``The restaurant meal delivery problem: Dynamic
pickup and delivery with deadlines and random ready times,'' \textit{Transportation Science},
vol.~55, no.~1, pp.~75--100, 2021, doi: 10.1287/trsc.2020.1000.

\bibitem{Hildebrandt2022}
F. D. Hildebrandt y M. W. Ulmer, ``Supervised learning for arrival time estimations in restaurant
meal delivery,'' \textit{Transportation Science}, vol.~56, no.~4, pp.~1058--1084, 2022,
doi: 10.1287/trsc.2021.1095.

\bibitem{Boyaci2021}
B. Boyacı, T. H. Dang y A. N. Letchford, ``Vehicle routing on road networks: How good is Euclidean
approximation?,'' \textit{Computers \& Operations Research}, vol.~129, art.~105197, 2021,
doi: 10.1016/j.cor.2020.105197.

\bibitem{Pereira2021}
R. H. M. Pereira, M. Saraiva, D. Herszenhut, C. K. V. Braga y M. W. Conway, ``r5r: Rapid realistic
routing on multimodal transport networks with R5 in R,'' \textit{Findings}, 2021,
doi: 10.32866/001c.21262.

\bibitem{Kuhn1955}
H. W. Kuhn, ``The Hungarian method for the assignment problem,'' \textit{Naval Research Logistics
Quarterly}, vol.~2, no.~1--2, pp.~83--97, 1955, doi: 10.1002/nav.3800020109.

\bibitem{Boeing2022}
G. Boeing, ``Street network models and indicators for every urban area in the world,''
\textit{Geographical Analysis}, vol.~54, no.~3, pp.~519--535, 2022, doi: 10.1111/gean.12281.

\bibitem{Geisberger2008}
R. Geisberger, P. Sanders, D. Schultes y D. Delling, ``Contraction hierarchies: Faster and simpler
hierarchical routing in road networks,'' en \textit{Experimental Algorithms (WEA 2008)}, Lecture
Notes in Computer Science, vol.~5038, Berlín: Springer, 2008, pp.~319--333,
doi: 10.1007/978-3-540-68552-4\_24.

\bibitem{LuxenVetter2011}
D. Luxen y C. Vetter, ``Real-time routing with OpenStreetMap data,'' en \textit{Proc. 19th ACM
SIGSPATIAL Int. Conf. Advances in Geographic Information Systems}, Chicago, IL, EE.~UU., 2011,
pp.~513--516, doi: 10.1145/2093973.2094062.

\bibitem{Wu2024}
L. Wu \textit{et al.}, ``LaDe: The first comprehensive last-mile express dataset from industry,''
en \textit{Proc. 30th ACM SIGKDD Conf. Knowledge Discovery and Data Mining}, 2024,
pp.~5991--6002, doi: 10.1145/3637528.3671548.

\bibitem{Wilcoxon1945}
F. Wilcoxon, ``Individual comparisons by ranking methods,'' \textit{Biometrics Bulletin}, vol.~1,
no.~6, pp.~80--83, 1945, doi: 10.2307/3001968.

\bibitem{Friedman1937}
M. Friedman, ``The use of ranks to avoid the assumption of normality implicit in the analysis of
variance,'' \textit{Journal of the American Statistical Association}, vol.~32, no.~200,
pp.~675--701, 1937, doi: 10.1080/01621459.1937.10503522.

\bibitem{MunozVillamizar2021}
A. Muñoz-Villamizar, E. L. Solano-Charris, L. Reyes-Rubiano y J. Faulin, ``Measuring disruptions in
last-mile delivery operations,'' \textit{Logistics}, vol.~5, no.~1, art.~17, 2021,
doi: 10.3390/logistics5010017.

\end{thebibliography}

\end{document}
